\quellseite{230}
Denn GOtt soll man allein anbeten, und ihm allein dienen. Nun ist das Gelübde auch ein solcher Dienst; darum soll es niemand widerfahren denn ihm. Derhalben ist schon geantwortet, daß nichts gilt, was man den Heiligen gelobt. Dein Gelübde reimt sich mit diesem des Patriarchen Jakob nicht; darum führst du die Schrift falsch dazu, wenn du dich damit behelfen willst. Also fallen gar viel Gelübde hin, so man bisher in aller Welt gethan hat, nicht GOtt, sondern den Heiligen.

21. Der andere Mangel ist, daß hier steht ein solches Gelübde, das vorhin von GOtt gegeben ist, daß er es halten will. Denn er setzt merklich dazu: „Wenn mich GOtt wieder heim zum Vater bringen und behüten wird auf dem Wege, Essen und Trinken und Kleider geben, so soll dieser Ort ein Gotteshaus sein, und so viel du mir geben wirst, will ich dir den Zehnten geben.“ Solches alles, was er gelobt, will er zuvor haben, daß er es thun könne, und nicht eher. Darum streckt es sich nicht weiter, denn so ferne, daß die Güter da sind, und so er es wohl vermöge. Gibt er ihm wenig, so soll er wenig\footnote{Wittenberger und Jenaer: ein wenig.} nehmen; gibt er ihm viel, so soll er viel haben; gibt er ihm nichts, so darf er nichts verzehnten. Also, daß sich sein Gelübde nicht gründet denn auf die Güter, die er schon zehnfältig hat im Vorrathe.

22. Nun halte unsere Gelübde dagegen. Wenn du gelobst, Keuschheit zu halten, was hast du gelobt, denn ein Ding, das gar nicht in deiner Gewalt steht? „GOtt schuf sie ein Männlein und Fräulein, und sprach: Wachset und mehret euch“ [1 Mos. 1, 27. 28.]. Was darfst du nun geloben, nicht ein Mann oder ein Weib zu sein, weil dich GOtt also geschaffen hat; und nicht unfruchtbar also geschaffen, sondern daß du solltest fruchtbar sein? Hast du doch den Vorrath nicht, daß du es könntest halten, so du es doch wohl zehnfältig solltest im Vermögen haben. Warum ist es der Teufel und Tod, daß man gelobt, das uns GOtt nicht gibt, und wir nicht halten können. Du willst geloben, fromm zu sein, und bist ein Bube geboren. Was wäre es, wenn ein Krauter auf dem Bette gelobte, gesund zu sein und aufzustehen? Es heißt GOtt gespottet, wenn man ihm geben will, das wir nicht haben. Wenn er es zuvor gegeben hätte, und du darnach gelobtest, so wäre es eine Meinung.

23. Dazu kannst du viel weniger ewige Keuschheit geloben. Denn wenn du gleich heute fromm und keusch bist, morgen liegst du vielleicht im Drecke. Denn es ist alles durch und durch in GOttes Händen; darum kann ich es nicht weiter noch länger geloben, denn er gibt. Thut es doch Jakob in zeitlichem Gute, und will nichts weiter geloben, denn ihm GOtt geben will. Lieber GOtt, gilt dies Gelübde in solchem Gute nur so ferne, welches doch mehr in unsern Händen steht denn geistlich Gut: was sind wir denn für Narren, daß wir ohne alle Mittel die hohen, großen, seltsamen Güter dahin geloben ohne\footnote{„ohne“ fehlt in den Ausgaben, doch hat die Jenaer die Conjectur am Rande: ohne oder wider.} den Heiligen Geist, welcher alle geistlichen Güter in seinem Schooße hat? Daher siehst du, wie unsere Klostergelübde alle des Teufels sind. Also kannst du sie auf das Maul schlagen, und weisen, wie sie den Text falsch führen, daß sie ihn lernen recht ansehen. Also thue du nun auch. Willst du je geloben, Jungfrau zu bleiben, so gelobe es, so lange es GOtt gibt, und nicht weiter. Das ist, daß wir gelehrt haben, daß alle Klostergelübde sollten frei sein, daß ein jeglicher also spreche: „Ich habe das oder dies gelobt, bin es aber nicht weiter schuldig, denn sofern ich, es zu halten, habe.\footnote{nämlich als eine Gabe von GOtt.}“ Warum? Denn du bist nicht geschaffen zur Jungfrauschaft, sondern Mann oder Weib zu sein. Wird sie dir aber gegeben, so hast du es; wird sie dir genommen, so kannst du es nicht mehr geloben noch halten. Darum schließt dieser Text nicht wider uns, sondern für uns wider sie. Das sind die zwei Gebrechen an unsern Gelübden.

24. Der dritte [Mangel] ist nun der allerärgste; denn jene sind noch äußerlich, dieser aber ist innerlich. Denn hier fehlt es an der Person. Werde zuvor Jakob, und gelobe darnach auch; wie ich oft gesagt habe, man solle nicht auf die Werke der Heiligen sehen, sondern auf die Person und ihren Glauben. Denn aus dem Irrthum ist es gekommen, daß Bethel hernach nicht GOttes, sondern des Teufels Haus ist worden. Denn der Teufel hat diese Historie und Exempel auch stark geführt, und den Propheten genug zu schaffen gemacht, daß sie den falschen Verstand niederlegten. Wie sollte es
\quellseite{231}
der Teufel seiner angerichtet haben, denn daß er einen solchen Schein für sich hatte, und ließ also predigen: Hier ist Jakob, der heilige Vater, gelegen, hat GOtt selbst hören predigen, und gelobt, daß das ein Gotteshaus sollte werden; wie könnten wir nun besser thun, denn daß wir einen köstlichen Gottesdienst stiften, da man bete und opfere?\footnote{In den Ausgaben: opferte.} 2c.

25. Nicht gerne wollte ich, daß der Papst einen solchen scheinenden Grund für sich hätte, als diese gehabt haben wider die rechten Propheten. Denn wer dürfte dawider sagen? Der Text ist klar, daß er sagt: „Hier wohnet GOtt“; darum sagen sie: Wir wollen unserm heiligen Erzvater Jakob nachfolgen, haben GOttes Wort und das Tempel bei uns. Da sagten die Propheten wiederum: Nicht also, ihr fahrt unrecht. Aber wenn sie das Maul aufthäten, müßten sie des Todes sein, und noch dazu hören, sie hätten wider die Schrift geredet.

26. Was ist denn da der Mangel gewesen, daß GOtt solch Werk verworfen hat? Also haben die Propheten gepredigt: Liebe Herren, ihr seid noch lange nicht Jakob; es ist nicht die Meinung, daß man den bloßen Werken der Heiligen nachahme;\footnote{In den alten Ausgaben: „nachome“; Tenger (falsch): „nachthome“.} Gott hat der Person die Stätte gegeben und zugelassen, und ein Wohlgefallen an der Person und Stätte gehabt, wie er denn mancherlei Werk, Zeit und Stätte braucht durch seine Heiligen, will aber darum nicht, daß ein jeglicher hinlaufe, und dasselbe nachthue. Du bist ein ehelicher Mann oder Weib, da braucht dich GOtt zu; wenn du nun das wolltest lassen fahren, und einem andern folgen, so würde dir GOtt nichts darum danken. Es gilt nicht, Tempel ansehen, sondern den Glauben und Berufung GOttes. Jakob ist zu der Stätte berufen; darum darfst du ihm nicht folgen, es wäre denn, daß du auch berufen würdest, und GOtt mit dir redete. Also haben die heiligen Könige gethan. David hatte es wohl im Sinne, GOtt ein Haus zu bauen; er durfte es aber nicht thun, weil es ihm GOtt nicht befohlen hatte [2 Sam. 7, 5. ff.]. Salomo aber, sein Sohn, ward dazu berufen und bauete den Tempel zu Jerusalem [1 Chron. 18, 11. 12. 1 Kön. 6, 1.].

27. Darum ist das die Summa davon: Wo und wie dich GOtt ruft, da folge nach; wo aber nicht, so lasse es anstehen, wenn du auch aller Heiligen Exempel vor dir hättest. Aber es hat nichts geholfen. Unsere Narren haben kein GOttes Wort aus der Schrift für sich gehabt, und dennoch täglich eine Kirche über die andern aufgerichtet, und nach ihrem Traume Gott\footnote{In den Ausgaben: Gotte.} wollen anbinden, wo er wohnen sollte, nämlich, an den Stuhl zu Rom. Wenn sie einen solchen Text hätten mögen aufbringen, wie die Juden thaten, könnte ihnen die ganze Welt nicht wehren.

28. Also haben sich die Propheten mit dem Volke geschlagen, und auf dem Spruche [2 Mos. 20, 24.] gestanden: „An welchem Orte ich meines Namens Gedächtniß mache, da will ich zu dir kommen, und dich segnen.“ Damit hat GOtt alle Stätte aufgehoben, und keine gewisse gesetzt, sondern will also sagen: Sehet auf mich, und erwählet mir keine Stätte; sondern wo ich eine erwählen werde, da kommet hin, da will ich sein und predigen. Damit ist aller Welt begegnet, daß man nichts aufrichte zu Gottes Dienst, er thue und schaffe\footnote{schaffe = gebiete. Vgl. Walch, St. Louiser Ausg. Bd. IX, 1054, § 12.} es denn zuvor. Siehe nun alle Historien an, und merket, wie wir Narren sind gewesen, wenn man uns sagte: Da liegt St. Peter, St. Jakob, das heilige Grab, diese oder jene Heiligen; da ist GOtt gewesen, und hat die Stätte geheiliget; so bald wir das Geschrei gehört haben, sind wir zugelaufen als blinde und thörichte Leute, und haben mit Haufen Kirchen gestiftet. Ja, wenn wir einen Finger oder Hauptstückel vom Heiligthume überkommen konnten, richtete man sobald Altar und Capellen auf. Noch bleiben sie darauf stehen, und trotzen: Da liegt dennoch St. Peter, St. Paulus etc., die sind ja heilig, darum muß die Stätte, wo sie liegen, auch heilig sein.

29. Du aber antworte also: Ist denn das genug dazu, daß GOtt allda wohne, davon er nie geredet noch befohlen hat? Wenn gleich alle Heiligen mit einander allda lägen, wollte ich nicht davor niederknien noch sie ansehen der Meinung, daß ich dächte, ich wollte GOtt an dem Orte sonderlich dienen. Es liegt nichts daran, wo die Heiligen liegen, sondern daran liegt es, wo GOtt redet. Wer ist davon gebessert, daß so viel Heilige zu Rom liegen in der Höllenstuppe, da kein GOttes Wort, sondern
\quellseite{232}
eitel Teufelslehre und sein Reich ist? Noch läuft man dahin, und wenn man viel Heilige gesehen hat, soll es genug sein. Willst du des Seegens warten und rechte Gnade erlangen, so laufe dahin, da GOttes Wort ist, und hüte dich nur vor andern Stätten, da es nicht ist, wenn auch alle Engel da wären. Die Heiligen thun nichts, sondern das lebendige Wort thut es, das ist besser, denn aller Heiligen Erde und Kirchen, da sie liegen.

30.\footnote{Die Wittenberger und die Jenaer Ausgabe haben hier die Ueberschrift: „Verlegung der Gelübde“.} Also verlege\footnote{verlegen = widerlegen} nun, daß sie aufbringen, die Gelübde zu erhalten. Zum ersten: Jakob ist ein Mann im Glauben; werde du auch derselbe, und gelobe darnach. Zum andern gelobt er dem rechten GOtt und spricht: „So GOtt wird mit mir sein und mich behüten“ etc. Solche Worte redet niemand, denn der Heilige Geist. Denn wenn du also sprichst: Ich gelobe GOtt, der Himmel und Erde geschaffen hat, daß ich da will eine Kirche bauen, ist [es] so bald erlogen. Denn du hast keinen Geist, der dich es heiße, und geht aus deinem eigenen Willen und Andacht; [du] hast kein Zeichen noch Zeugniß davon; du solltest dein Lebelang gegen Gott nichts anfangen, du siehest denn gewiß, daß er dich es heiße. Wo nicht, so sprich frei, das hat der Teufel gethan.

31. Darum siehe den Text recht an. Ehe Jakob gelobt, zeigt ihm GOtt, daß er da wohnt, also, daß er eine äußerliche Ermahnung hat an dem Zeichen, daß er sieht, dazu auch ein innerlich Treiben des Geistes. Wenn dich nun Gott auch berufet mit solchen Zeichen, und dazu einen solchen Geist gibt, so thue auch dergleichen. So sieht man, wie das Exempel für uns ist, und eben zu Boden stößt alles, was man gelehrt hat von Gelübden.

32. Dazu habe ich auch mehr gesagt, daß unsere Gelübde den schändlichen Zweck an sich haben, das da heißt, sich darauf verlassen, und damit einen Weg gen Himmel wollen machen, und im höheren Stande sein denn der gemeine Mann. Als, wenn man sie fragt, warum sie in das Kloster gehen? antworten sie: Daß wir Christi Bräute werden. Ja, des Teufels Bräute werden sie. Christus buhlt nicht mit Fleisch und Blut, sondern will die Seele haben. So haben wir alle sowohl Seelen als sie, und ist darum zu thun, daß wir seine Braut werden in reinem Glauben, in welchem wir empfangen das Wort GOttes, welches ist der geistliche Same [Luc. 8, 11. 1 Petr. 1, 23.]. Das machen nun die gottlosen, tollen Prediger, daß das Volk soll meinen, sie haben einen besseren, seligeren Stand, denn der gemeine Stand der Christenheit. Das ist der Teufel und Tod, daß, wenn sonst nichts Anderes darinne wäre, denn solche unchristliche Meinung, wäre es doch greulich genug. Darum sind ihre Gelübde nicht Jakobs Gelübde, sondern des Teufels. Denn sie verleugnen den Glauben, und wollen ihre Keuschheit besser geachtet haben denn den Heiligen Geist, und einen armen Christenmann, da Christus mit allen seinen Gütern wohnt, geringer achten denn ihre Krone oder Kappen. Sie sind nicht werth, daß sie ihm die Füße waschen. Das sei genug von dem. [Wir] wollen nun den Text überlaufen, und sehen, wie und wohin Jakobs Gelübde lautet.

33. Zum ersten spricht er: „Der HErr soll mein GOtt sein.“ Da ist der Glaube angezeigt, als er sagt „der HErr“, welches der rechte Name GOttes ist. Er hatte ihn vorhin für seinen rechten GOtt, im rechten Glauben, noch spricht er: „Der HErr“, den ich schon habe, „soll mein GOtt sein.“ [Es] ist aber so viel gesagt: Ich will ihm den äußerlichen Dienst thun, nicht darum, daß er erst mein Herr sei, sondern mein GOtt; das ist, den HErrn habe ich im Herzen durch den Glauben, jetzt aber will ich\footnote{Jenaer: „ich in“ [d. i. ihn].} auch äußerlich bekennen und predigen vor der Welt, auf daß sie sehen, daß ich auch einen GOtt habe. Denn ein Christ muß GOtt im Herzen haben, also, daß er für sich nichts mehr bedarf. Doch muß er auch heraus brechen und bekennen, daß er einen GOtt habe, mit äußerlichem Gottesdienste. Was wäre es sonst hier gesagt, wenn er meinte, der HErr sollte erst sein GOtt werden, als hätte er ihn zuvor nicht für seinen GOtt gehalten.

34. So ist nun das\footnote{„das“ fehlt in der Erlanger.} die Meinung, daß er nun einen sonderlichen Ort, zu dem er denn berufen ist, nehmen will, und allda sein Gesinde, und die bei ihm sein werden, zusammen bringen, zu predigen, beten, und was mehr zu thun ist gegen GOtt. [Es] ist auf das kürzeste so viel,
\quellseite{233}
als spräche er: Ich will einen äußerlichen Dienst anrichten. Doch wollte er es nicht thun, denn als er dazu vermahnt war. So bleibt sein Gelübde in reinem Glauben; nicht, daß er dadurch fromm werden oder ein verdienstlich Werk thun wolle, sondern solchen Gottesdienst halten, allein daß andere Leute dadurch gebessert, und GOtt gepriesen werde. Also möchtest du auch thun, und geloben, daß er sollte dein GOtt auf einem Orte sein, zu predigen und anzurufen; nicht als ein Werk, damit du wollest den Himmel verdienen, sondern deinen GOtt und Glauben vor der Welt zu bekennen, und andern Leuten zu dienen.

35. Zum andern gelobt Jacob: „Alles, was du mir gibst, deß will ich dir den Zehnten geben“, nämlich, daß da würde gehalten ein Priester, der da predigte und lehrte. Denn denjenigen, die GOttes Wort lehrten, ist von Anfang gehalten, den Zehnten zu geben, welches jetzt im neuen Testamente nicht ist, sondern ist nachgelassen und heimgestellt der Liebe, daß man nicht darf Gesetze machen, sondern gebe so viel als noth ist, einen Prediger zu nähren und halten. Also lehrt St. Paulus Gal. 6, 6.: „Der da unterrichtet wird mit dem Worte, der theile mit allerlei Gutes dem, der ihn unterrichtet“; und 1 Cor. 9, 14.: „Die das Evangelium verkündigen, sollen sich auch des Evangelii nähren.“ Wo nun Christen sind, so lassen sie es nicht; aber da nicht Christen waren, und das Volk mit Zwang und Gesetzen regiert werden mußte, so mußte auch eine gewisse Summa bestimmt sein, wie viel man ihnen geben sollte. Es haben’s die auch müssen thun, die vor uns gewesen sind, daß man die Prediger nicht konnte versorgen auf die Weise. Die Liebe ist zu schwach unter dem gemeinen Haufen, sie thut es nicht.

36. Im Alten Testamente [4 Mos. 18, 24.] war es geboten, man sollte den Leviten den Zehnten geben; noch mußten sie sich zuletzt selbst ernähren, denn man wollte ihnen nicht mehr geben. Darum erdachten sie darnach einen falschen Gottesdienst, da gab jedermann wieder genug, wie es zu unsern Zeiten auch geht und gegangen ist. Denn, weil man den rechtschaffenen Predigern nicht gibt, daß sie sich ernähren, müssen sie zuletzt das Predigen fahren lassen, und sich mit Arbeiten selbst ernähren; so wird GOtt zur Plage und Strafe wieder lassen Prediger auftreten, die uns zum Teufel führen. Darum ist das ein edel Werk des frommen Vaters, daß er im Sinne war, GOttes Wort zu erhalten. Gelobe du auch also, und thue es ihm nach, so will ich es loben. Aber weil das Werk so gut ist, folgt ihm niemand; zu Klostergelübden aber gibt man mehr, denn genug ist. Das macht, es ist des Teufels Werk. Was GOttes Wort und Exempel ist, da will die Welt schlecht nicht an, läßt es zu einem Ohre ein-, zum andern ausgehen.

37. Also haben wir nun klärlich genug gesehen, was das Exempel in sich hält, und wie sie es mißbraucht haben, und wie es uns lehrt, daß wir auch schuldig sind, zu geben zu dem rechten Gottesdienste, nämlich, das Wort zu erhalten; doch so ferne als uns GOtt gibt, daß wir es vermögen, und nicht der Meinung, ein gut Werk zu thun, damit man etwas verdiene, sondern andern Leuten zu Nutzen, also, daß es ein Gelübde der Liebe sei, mehr denn des Glaubens. Also kannst du auch dem Nächsten alle Jahre so viel zu geben geloben, so lange du es vermagst, GOtt zu Ehren, dem, der dich GOttes Wort lehrt. Das wäre auch ein fein Gelübde, frei aus der Liebe. Daß also die drei Stücke bei einander seien in einem christlichen Gelübde: erstlich, daß es GOtt allein geschehe; zum andern, so viel wir es vermögen; zum dritten, dem Nächsten zu Nutze; und zuletzt, so ferne, daß der Glaube unversehrt bleibe. Was nun mehr im Texte steht, von dem Steine und Oele, damit ihn Jakob begoss, dient auf die geistliche Deutung der Historie, die wollen wir auch ein wenig rühren.

\subsection{Geistliche Deutung}
38. Der Stein, mit Oele begossen, ist Christus, unser einiges Haupt und Grundstein, wie ihn die Schrift auch nennt [1 Cor. 10, 4. Eph. 2, 20. 1 Petr. 2, 6.], auf welchen auch gedeutet hat des Patriarchen Schlaf; denn der Schlaf ist nichts Anderes denn der Glaube. Wenn die Sonne untergeht, finster und Nacht wird, daß der Leib schläft und nicht sieht, noch fühlt, die Seele aber wacht, damit die Figur christliches Lebens vorgebildet ist, welches also soll stehen, daß es nichts von der Welt sieht noch fühlt, lebt aber im Glauben: da sieht es GOtt und die Engel auf der Leiter. Also, daß der Glaube nichts ist denn ein rechtschaffener Traum, in welchem ein Mensch sieht, das sonst niemand
\quellseite{234}
sieht. Daß er nun den Stein unter den Häupten hat, deutet auf Jesum Christum, auf welchem unsere Gewissen ruhen. Denselben richtet er auf, als er erwacht, und gießt Oel darauf; das ist, wenn wir im Glauben sind und schlafen, machen wir darnach auf, brechen heraus durch äußerliches Bekenntniß; so richten wir denn den Stein auf, also, daß wir predigen von Christo, der uns tröstet; und gießen Oel darauf, das ist, daß solche Predigt ist von lauter Gnade und Barmherzigkeit, welche von Christo allein fließt und quillt.

39. Und damit ist eben angezeigt, was er für eine Kirche oder Gotteshaus hat wollen anrichten, nämlich allein dazu, daß man darinne predigte von Christi Gnade, welcher ihm in dem Schlafe und Traume von GOtt verheißen war; wie zuvor seinen Vätern, Abraham und Isaak; darum gelobt er auch eben dasselbige. Aber das wird niemand thun, er habe denn\footnote{In den alten Ausgaben: „denn er habe“.} die Leiter zuvor gesehen, das ist, wer nicht vorhin Christum kennt, der wird nicht wohl von ihm predigen. Denn die Leiter bedeutet auch Christum, auf welchem die Propheten und die ganze Schrift auf und absteigt, das ist, es stimmt alles auf ihn, und geht alles auf der Leiter. Wer nun\footnote{Wittenberger und Erlanger: den.} Christum so kennt, und sieht, wie alle Schrift auf ihn geht, da wird ein gewaltiger Prediger aus. Das hat diese Figur auf das kürzeste in sich, also, daß es alles gehe auf das Predigtamt von Christo.



\section{Das neunundzwanzigste Capitel}
\begin{schriftzitat}
V. 1—8. Da hub Jakob seine Füße auf und ging in das Land, das gegen Morgen liegt. Und sahe sich um, und siehe, da war ein Brunnen auf dem Felde, und siehe, drei Heerden Schafe [lagen] dabei. Denn von dem Brunnen mußten die Heerden trinken, und lag ein großer Stein vor dem Loche des Brunnens. Und sie pflegten die Heerden alle daselbst [zu] versammeln, und den Stein von dem Brunnenloche [zu] wälzen, und die Schafe [zu] tränken, und thaten alsdann den Stein wieder vor das Loch an seine Stätte. Und Jakob sprach zu ihnen: Brüder, wo seid ihr her? Sie antworteten: Wir sind von Haran. Er sprach zu ihnen: Kennet ihr auch Laban, den Sohn Nahor? Sie antworteten: Wir kennen ihn wohl. Er sprach: Gehet es ihm auch wohl? Sie antworteten: Es gehet ihm wohl; und siehe, da kommt seine Tochter Rahel mit den Schafen. Er sprach: Es ist noch hoch Tag, und ist noch nicht Zeit das Vieh einzutreiben, tränket die Schafe, und gehet hin und weidet sie. Sie antworteten: Wir können nicht, bis daß alle Heerden zusammen gebracht werden, und [wir] den Stein von des Brunnens Loche wälzen, und also die Schafe tränken.
\end{schriftzitat}

1. Hier beschreibt Moses, wie der Patriarch Jakob in das Land Mesopotamien gekommen ist, und daselbst Weib und Kinder überkommen hat. Da siehest du abermal, daß GOtt seiner lieben Heiligen Legenden läßt wunderlich beschreiben; daß so eben daher erzählt wird, wie Jakob gegangen ist, und auf dem Wege nichts denn einen Brunnen und drei Heerden Schafe gefunden hat, und wie er mit den Hirten redet. Welches sich doch gar läßt ansehen, als ein unnütz, unnöthig Ding zu beschreiben, so er doch wohl andere Dinge und größere Stücke hätte können finden, die etwas Ansehens hätten. Aber das ist es, so wir oft gehört haben, daß es alles darum geschieht, auf daß Gott der Vernunft die Augen aussteche, und sie blende, daß man nicht sehe nach großem, scheinendem Wesen und Werken, sondern, wie uns die ganze Schrift weiset auf die Werke, die GOtt thut. Wenn wir dahin kämen, und ließen es gehen, wie es GOtt gemacht hat, so ginge es recht.

2. Als, da Adam und Eva im Paradiese waren, da war kein selbsterdacht, menschlich Werk, sondern ging alles, wie es GOtt gepflanzt und die Natur gegeben hatte. Das waren auch die allerbesten Werke. Nun aber die Natur also gefallen ist, sind wir gerathen auf unsere eigenen Werke, und ist die Welt toll und thöricht darauf, will nur mit andern Werken umgehen, denn [die]
\quellseite{235}
GOtt gepflanzt hat, also, daß GOtt in allen Worten und Werken damit handelt, daß er es in Schwang bringe, zu gehen, wie er es geschaffen hat. Das ist nun ein groß Erkenntniß, welches auch großen Heiligen fehlt. Wir bekennen alle im Glauben: Ich glaube an GOtt Vater, allmächtigen Schöpfer Himmels und Erden. Wer es also glaubt, der weiß viel.

3. Darum ist auch dies also beschrieben, daß Jakob kommt an einen Ort zu den Hirten, und redet von Sachen, die man ansieht, als gehen sie GOtt gar nichts an. Wenn er einen heiligen Mann beschreiben wollte, sollte er nicht vielmehr also schreiben: Er kam an einen Ort und lehrte die Leute gute Werke thun, hieß sie zur Kirche gehen, Messe hören und beten; und so fortan von großen geistlichen Werken: so führet er herein, und weiß nichts zu sagen, denn wie er mit den Hirten spricht von Schafen zu tränken. Wir haben es genug gehört, wie Gott Lust hat, mit solchem Narrenwerke umzugehen. Was haben wir davon, daß er sie heißt die Schafe tränken, und hernach den Stein vom Brunnen wälzt? [Es] ist wohl der Rede werth, daß er solch Ding schreibe, und die Welt lasse daran lesen, gleich als hätte man sonst nichts zu lesen. Das sind nun die rechten natürlichen Werke, welche GOtt in die Natur gepflanzt hat, und dazu geordnet, daß der Mensch das Vieh regiere, gehe damit um, und suche dadurch seine Nahrung, Essen, Trinken, und Futter; die werden beschrieben von den großen Heiligen. Wollte GOtt, daß [es] uns auch so einginge, so würden wir sehen, welche Narren wir sind.

4. So ist nun das ganze Capitel eitel Narrenwerk vor der Welt, und eine Predigt, die gar nichts zur Sache dient. Christus selbst und Paulus predigen also, wie man sieht fast in allen Episteln. St. Paulus, wenn er von christlichen Werken will lehren, weiß er nichts zu sagen, denn wie der Mann seines Weibes und Kinder warten soll [Eph. 5, 25. (Col. 3, 19. Eph. 6, 4.)], das Weib den Mann fürchten und in Ehren halten, und Kinder ziehen [1 Tim. 2, 12.], der Knecht dem Herrn gehorsam sein [Eph. 6, 5. f.], die Kinder den Eltern [Eph. 6, 1—3.], noch gehen wir mit andern Werken um. Die großen Apostel wußten nichts Besseres zu predigen; sollten sie aber jetzt schreiben, müßten sie vorhin noch eine Weile studiren auf den hohen Schulen, und viel Geld geben, daß sie Doctores würden. Die Vernunft kann es wohl besser treffen, wie man GOtt dienen soll; da machen sie große Bücher davon, und disputiren, welche die besten Werke sind.

5. Wir aber mögen GOtt wohl danken, wenn wir können, daß uns die Schrift vor die Thüre kommt, und zeigt uns recht an, was GOtt wohlgefällt, was er selbst darf loben in seinen höchsten Heiligen, was die Apostel dürfen preisen und predigen. [Es] ist billig, daß wir das Maul zuhalten, und sagen: Wir wissen es nicht besser, wäre etwas Besseres, hätte er es uns auch wohl lassen schreiben und predigen. [Es] ist alles darum zu thun, daß GOtt hat mit diesem Buche wollen wehren allem Jammer, der da gefolgt ist\footnote{In den Ausgaben: hat.} (wie man in folgenden Büchern sieht), und erhalten, daß er schafft und geschaffen hat; also, daß ein jegliches seines Dinges, Standes und Amtes warte. Willst du etwas Anderes heißen suchen, so fehlst du. Das gibt beide, die Erfahrung und die Schrift; noch sind wir toll und thöricht, daß wir solches in Wind schlagen, und auf andere Dinge fallen. Das sei davon geredet; mehr weiß ich nicht daraus zu machen. Folgt weiter im Texte:

\begin{schriftzitat}
V. 9—30.

Als er noch mit ihnen redete, kam Rahel mit den Schafen ihres Vaters; denn sie hütete der Schafe. Da aber Jakob sahe Rahel, die Tochter Labans, seiner Mutter Bruder, und die Schafe Labans, seiner Mutter Bruder, trat er hinzu und wälzte den Stein von dem Loche des Brunnens, und tränkte die Schafe seiner Mutter Bruders, und küssete Rahel, hub auf seine Stimme und weinete, und sagte ihr an, daß er ihres Vaters Bruder wäre und Rebekka’s Sohn. Da lief sie und sagte es ihrem Vater an. Da aber Laban hörte von Jakob, seiner Schwester Sohne, lief er ihm entgegen, und herzte und küssete ihn, und führte ihn in sein Haus. Da erzählte er ihm alle diese Geschichte. Da sprach Laban zu ihm: Wohlan, du bist mein Bein und mein Fleisch, bleibe bei mir einen Monden lang. Darnach sprach er aber zu Jakob: Wiewohl du mein Bruder bist, sollst du mir darum umsonst dienen? Sage an, was soll dein Lohn sein? Laban aber hatte zwo Töchter, die älteste hieß Lea und die jüngste Rahel. Aber Lea hatte ein blödes Gesicht, Rahel war schön gestaltet und schön von Angesicht. Dieselbige gewann Jakob lieb und sprach: Ich will dir sieben
\quellseite{236}
Jahre um Rahel, deine jüngste Tochter, dienen. Laban antwortete: Es ist besser, ich gebe dir sie, denn einem andern, bleibe bei mir. Also dienete Jakob um Rahel sieben Jahre, und däuchten ihn, als wären es einzelne Tage, so lieb hatte er sie. Und Jakob sprach zu Laban: Gib mir her mein Weib, denn die Zeit ist hier, daß ich beiliege. Da lud Laban alle Leute des Orts, und machte eine Hochzeit. Des Abends aber nahm er seine Tochter Lea, und brachte sie hinein, und er beschlief sie. Und Laban gab Lea, seiner Tochter, Silpa zur Magd. Des Morgens aber, siehe, da war es Lea. Und er sprach zu Laban: Warum hast du mir das gethan? Habe ich dir nicht um Rahel gedienet? Warum hast du mich denn betrogen? Laban antwortete: Man thut nicht also in unserm Lande, daß man die Jüngste ausgebe vor der Aeltesten. Halte diese Woche aus, so will ich dir diese auch geben, um den Dienst, den du mir noch andere sieben Jahre dienen sollst. Jakob that also, und hielt die Woche aus. Da gab er ihm Rahel, seine Tochter, zum Weibe, und Laban gab Rahel, seiner Tochter, Bilha zur Magd. Also lag er auch bei mit Rahel, und hatte Rahel lieber denn Lea, und dienete ihm fürder die andern sieben Jahre.
\end{schriftzitat}

6. Da beschreibt er, wie Jakob nun gekommen ist zu Laban, der ihn nennt seinen Bruder, so sie doch Geschwister Kinder mit einander waren. Hier sollen wir aber Frau Hulda, der tollen Närrin, der Vernunft, antworten. Wie haben sich die Hochgelehrten daran gestoßen, daß der Patriarch so zu einem Narren wird, und nimmt zwei Weiber, ja, zwo Schwestern auf einmal! Dazu, welches viel närrischer ist, daß er um Rahel sieben Jahre dient, und sie so lieb hat, daß er sich läßt dünken, es wäre nur einer Woche lang, daß billig unsere Geistlichen müssen sagen, es sei lauter Thorheit. Ich müßte es auch selbst sagen, wenn ich ein Mönch wäre, welche man nicht lehrt denken, was die Welt thut, sondern nur, was die Engel im Himmel thun; und sollen hier lesen, wie der Narr um die Metze buhlt sieben ganze\footnote{Walch und die Erlanger: ganzer.} Jahre, und hält es für sieben Tage. Reime dich, so lange mit Metzen buhlen, und doch heilig sein. Aber hier steht die Schrift; wer es besser machen kann, der thue es, wir wollen dieweil zusehen.


7. Was wollen wir nun daraus machen? Wir wollen ihn im Drecke lassen stecken, und bekennen, daß er ein Narr sei, daß er bisher sein Leben also gewagt hat, und nun erst hinkommt und hängt sich an eine Metze; und, als er betrogen wird, nimmt er die andere Schwester, und dient um diese noch sieben Jahre. Aber daß Jakob solches nicht aus Vorwitz gethan habe, ist daraus wohl zu merken, daß er nun zum wenigsten bei siebenzig Jahren ist alt gewesen, daß ihm ja der Kitzel sollte vergangen sein. Denn Esau, sein Bruder, nahm ein Weib, da er vierzig Jahre alt war; er aber harrt noch dreißig Jahre. Und wiewohl es der Text nicht gibt, wie alt er gewesen ist, doch, wenn man es soll rechnen zurück von Joseph her, so findet man, daß es fast läuft auf siebenzig Jahre, ja, noch weiter, schier auf achtzig Jahre.

8. Das thue nun dem Jakob nach, und lasse sie einen Spott daraus machen und sagen: er habe nichts gethan, denn gebuhlt und bei Frauen gelegen. Gib uns aber auch solche Leute, eben von denen, die so hoch die Jungfrauschaft preisen und treiben, der so lange Keuschheit gehalten habe. So viel als ihrer unter dem Papstthum sind, wenn man sie zusammen schmiedete, sollte man nicht Einen finden, der bis in das vierzigste Jahr Keuschheit gehalten hätte: noch wollen sie viel sagen von der Jungfrauschaft, und alle Welt tadeln; stecken selbst im Drecke bis an die Ohren, und soll nichts sein, daß dieser fromme Jakob siebenzig Jahre geharrt hat. Daß er nicht sollte auch Fleisch und Blut gehabt haben, ist nicht zu glauben. Denn er die Rahel so bald lieb gewinnt, als er meint, daß sie ihm werden soll; noch hat er so lange Zeit sich enthalten, daß ich den Jakob dürfte setzen über alle Jungfrauen. Denn so lange sich zu enthalten, muß eine Gnade und Geist sein, der da groß ist. Darum gilt es nicht, daß man so überhin läuft, und hält es für Leichtfertigkeit, und den Ernst nicht ansieht, den er gehabt hat. Wenn ich ein Exempel sollte setzen zur Keuschheit, so wollte ich eben diesen Jakob setzen; setze du mir einen dagegen im Alten und Neuen Testamente, deß Keuschheit also gepriesen sei. Darum ist GOtt wunderbarlich mit den Seinen, macht die Welt zu Narren, daß sie nicht weiß, wie sie daran ist.

9. Wiederum, in solcher großen Gnade, daß er so lange rein gelebt hat, ist nun das andere
\quellseite{237}
Merkt, daß GOtt dennoch läßt die Natur in ihm bleiben. Ob er darinne genarrt habe, daß er die Rahel so lieb hatte, und sieben, ja, vierzehn Jahre um sie dient, lasse ich gehen. Ich habe auch gerne, daß die Heiligen mit im Schlamme stecken, wie mir. Aber das ist wahr, daß GOtt neben dem Geiste und der Gnade die Natur läßt bleiben, wie er sie geschaffen hat, damit den großen Heiligen das Maul gestopft werde, die die Heiligen so hoch heben, als hätten sie keine natürliche Bewegung gefühlt. Zwar, wenn man es von mir schriebe, würde mich es wohl verdrießen, daß man von mir lesen sollte, wie ich mit Wehen gebuhlt hätte; und Moises schreibt es von dem heiligen Manne, der so lange ohne ein\footnote{„ein“ fehlt in der Wittenberger und in der Erlanger.} Weib gelebt, und nun ein alter Mann war.

10. Nun, was sollen wir daraus machen? Es sei gleich ein Ueberlanges\footnote{überlang = mehr als nöthig, überflüssig.} von ihm gewesen, ist uns doch dabei angezeigt, daß Gnade läßt die Natur bleiben; also, daß wenn du ein Christ bist, darfst du darum nicht gar ausziehen, was Fleisch und Blut ist. GOtt läßt dich Mann oder Weib bleiben; wie er es nun weiter ordnet, daß der eheliche Stand soll sein, Früchte zu zeugen, das hebt [die] Gnade nicht auf, sondern bestätig'ts, und macht, daß es alles seliglich wird. Wird nun jemand die Gnade gegeben, wie Jakob, so lange ohne Weib zu bleiben, wohl ihm; wo nicht, so weiß er wohl, wo er hin gehört.

11. GOtt hat dem Jakob die Gnade sonderlich gegeben, siebenzig Jahre rein zu leben; darnach, als er den Spruch gehört, daß durch seinen Samen sollte die Welt ausgebreitet und voll werden, und nun sieht, daß [es] Zeit ist, stellt er sich dazu. Darum hat er gelobt in solchem Geiste, daß er freilich nicht genarrt hatte, und ein solches ohne GOttes Wort gethan.

12. Das ist es nun, das ich oft gepredigt habe, die Gewissen frei zu machen. Denn wir sind nun von GOttes Gnaden so gelehrt, daß wir wissen, warum das Weib geschaffen ist. Das ist genug verstanden, man darf nicht mehr davon sagen; allein, daß die Gewissen gerettet werden, so mit falscher Lehre gefangen sind, daß man schier gemeint hat, es wäre eine Schande, ein Weib oder Mann zu nehmen. Es muß je also zugeben, daß man Lust und Liebe zusammen habe. Ohne Liebe ist es eitel Jammer und Noth; [es] kommt sonst noch Unglücks genug. Darum will der Heilige Geist nicht zürnen, ob eine Metze einen Knaben lieb hat, und wiederum; also, daß sie Lust haben, ehelich zu leben. Er hat es dem Jakob geschenkt, und nicht allein das, sondern auch lassen schreiben; nicht um seinetwillen, sondern um der tollen Prediger willen, die es so enge gespannt haben, und uns zu trösten. Das sind die zwei Stücke, daß GOtt den so begnadet mit Keuschheit, und doch bei der Natur läßt bleiben, wie sie geschaffen ist.

13. Wie wollen wir aber das entschuldigen, daß er zufährt, und nimmt zwo Schwestern? Das können wir nicht leugnen. Es ist ja Sünde bei dem Pabste, die man mit Feuer muß büßen; so hart, wenn es geschähe in solchem Falle, daß einer mit seines Weibes Schwester fiele, darf er nimmer sein Weib beschlafen. Hier siehe nun, dieser Text gilt mehr denn hunderttausend Pabste; und weil die Schrift Jakob nicht straft, mußt du mir das lassen ohne Sünde gethan sein. Weiter sagt der Text auch: „Da der HErr sahe, daß Lea gehasset ward, machte er sie fruchtbar, und Rahel unfruchtbar.“ Da bestätigt er den Irrthum, und will es haben, daß er die zwo Schwestern habe und behalte, und mache die Eheleuten, die sich so feindlich darüber martern, zu Narren. Wahr ist es, dem Exempel ist nicht also nachzufolgen, weil GOtt ein Sonderliches mit dem Manne thut, das er mit andern nicht gethan hat; sondern darum sage ich es, daß man hierinne sehe GOttes Werk, und wenn ein solcher Fall käme, daß man sich wüßte darein zu richten. Denn ich weiß wohl, wie man die Gewissen getrieben und gehandelt hat unter des Pabstes Gesetz.

14. Das erste nun, hier zu sagen, ist das: Wenn es nicht Gebrauch gewesen wäre in dem Lande, solches zu thun, so würde es Jakob ohne Zweifel nimmermehr gethan haben. Denn sein Herz stund also, daß er Rahel wollte haben zu einem Weibe, und nicht beide Schwestern. Daß ihn nun Laban täuscht, und legt ihm vorerst Lea bei, an Rahels Statt, ist ohne seine Schuld geschehen; [er] kommt also dazu, daß er nicht weiß wie, thut als ein frommer Mann thun muß, weiß nicht anders zu thun. Hätte er nicht wollen evangelisch fahren, und bei seinem Rechte bleiben, so hätte er gesagt: Da hab dir deine
\quellseite{238}
Tochter wieder; willst du mich narren, so narre mich auf deinen Schaden. Er wollte sich aber nicht rächen; wiewohl der Vater handelte wie ein Schalk, dennoch will er sie es nicht lassen entgelten, und thut ihr die Freundschaft, daß er sie behält. Derhalben ist es ein Zeichen, daß es auch des Landes Gewohnheit ist gewesen, und nicht so seltsam, daß niemand, denn er allein, solches gethan soll haben. Dazu bestätigt es GOtt auch, daß er es geschehen läßt. Darum hat er nicht unrecht noch sträflich gethan.

15. Doch, weil es nicht mehr Landes Sitte ist, und GOtt nicht befiehlt, oder Ursache gibt, dergleichen zu thun, soll man es lassen. Auch hat er es hernach aufgehoben, wie 3 Mos. 18, 18. steht: Du sollt kein Weib nehmen sammt ihrer Schwester, ihre Scham zu blößen, weil sie noch lebt. Unter den Heiden aber ist es gewesen, und Mosis Gesetz war die Zeit noch nicht. Auch wenn es gleich verboten wäre gewesen, so ist er doch ohne seinen Willen und Wissen dazu gekommen, also, daß doch das Gebot hier nicht Raum hätte. Denn GOtt hat nicht gewollt, daß dies Mägdlein zu Schanden würde. So streckt sich das Gebot nicht weiter, denn wo einer aus Muthwillen seines Weibes Schwester nähme; er aber hat sie allein darum behalten, daß er sie bei Ehren erhielte. Die, so da\footnote{Erlanger: die nu.} Christen sind, wissen wohl, was sie thun sollen; doch, wenn sich so ein Fall begäbe, müßte ich auch nichts anders zu rathen, denn wie Jakob gethan hat: auf daß man nicht dürfte des Papsts Gefangener sein, und seines Weibes nicht gebrauchen. So sei nun das die Summa, daß Jakob ein wunderlicher Mann, voll Geistes, und doch auch Fleisch und Blut ist wie andere.

16. Wie aber, das hernach folgt, daß Lea gehaßt, und Rahel die Liebste sei gewesen? Antwort: Er muß ihr ja nicht gar feind gewesen sein, hat ja sechs Söhne mit ihr, und mit Rahel noch nichts; dennoch sagt der Text: „Er hatte Rahel lieber denn Lea.“ Da läßt GOtt Fleisch und Blut mit unterlaufen. Es sind ja auch zwei fromme Kinder gewesen die Weiber, noch bleibt Fleisch und Blut in allen dreien, wie wir noch mehr hören werden. Die Ursache bringt der Text mit, als er sagt: „Rahel war schön gestaltet, und schön von Angesicht; Lea aber hatte ein blöde Gesicht.“ Nicht weiß ich, ob er von blöden Augen, oder vom ganzen Angesicht redet; als mir sagen, mager, dürr und bleich von Angesicht, verfallen und nicht voll; welches mir am besten gefällt. Rahel aber ist eine hübsche, glatte Metze, und völlig gewesen.

17. Siehe nun, ist das nicht ein seltsamer Heiliger? Er sollte so geistlich sein, daß er nicht mehr ansähe, denn daß sie ein Weib wäre; noch hat er so großen Unterschied, und sieht bald nach der Schönsten; zu der hat er Lust, zur andern nicht. So tief läßt ihn GOtt noch im Fleisch stecken, daß sein Herz mehr an dieser hangt denn an jener. Lea war die erste Braut, sollte sich lassen tragen, noch mußte sie hernach gehen; das Gesinde im Hause, auch Jakob selbst, hat an Rahel gehalten, und [sie] Frau im Hause lassen sein; niemand hat viel auf Lea gegeben. Das hat sie auch wohl verdrossen; [sie] hat gedacht: Mein Vater hat auch nicht allzumohl bei mir gethan, hätte mich wohl besser können berathen. Solche Gebrechlichkeit läßt GOtt bleiben und schreiben von den großen heiligen Kindern, daß man von ihnen muß sagen und predigen, wie sie unweis und neidisch gewesen sind.

18. Was thut nun GOtt dazu? [Er] läßt sie alle drei straucheln, und immer Fleisch und Blut mitlaufen. Noch hat er so reine Augen, sieht nicht darauf, welche die Hübscheste ist oder nicht, und gibt der hübschen Frau kein Kind; aber das arme Mägdlein muß ihrer bald fünf oder sechs haben. Er hat die Weiber recht lieb, sieht nichts, denn daß sie Weiber sind, fragt nicht darnach, wie sie gestaltet sind. Ueber das, so thut er auch das: Wenn er sieht, daß man eine zum Aschenbrödel machen will, so fährt er zu, und gibt ihr Früchte, den andern keine, und hat die am liebsten, die die Häßlichste und Verachtetste ist, die muß die rechte Mutter sein. Denn von der Lea ist Christus gekommen, aus dem Stamme Juda. Wenn Rahel noch so hübsch wäre, so hilft sie es nicht; was Christi sein soll, das muß von ersten herunter und demüthig werden. Lea muß manches stolze Wort gelitten haben, auch vom Gesinde im Hause, und wird immerdar gefressen sein, daß die andere regiert hat. GOtt ist aber so geschickt, daß er herunter sieht, auch bei seinen Heiligen, die alle seine Kinder sind, und hält sich mehr zu denen, so die schwächsten sind, denn die sich etwas mehr dünken lassen. Darum spricht nun der Text:
\quellseite{239}
\begin{schriftzitat}
V. 31—35.

Da aber der HErr sahe, daß Lea gehasset ward, machte er sie fruchtbar, und Rahel unfruchtbar. Und Lea ward schwanger, und gebar einen Sohn, den hieß sie Ruben, und sprach: Der HErr hat angesehen mein Elend, nun wird mich mein Mann lieb haben. Und ward abermal schwanger, und gebar einen Sohn, und sprach: Der HErr hat gehöret, daß ich gehasset bin, und hat mir diesen auch gegeben, und hieß ihn Simeon. Abermal ward sie schwanger, und gebar einen Sohn und sprach: Nun wird sich mein Mann wieder zu mir thun, denn ich habe ihm drei Söhne geboren; darum hieß sie ihn Levi. Zum vierten ward sie schwanger, und gebar einen Sohn und sprach: Nun will ich dem HErrn danken; darum hieß sie ihn Juda, und hörte auf mit Gebären.
\end{schriftzitat}

19. Den ersten Sohn nennt sie Ruben, das ist, ein Sohn des Gesichtes, oder Sehesohn; als sollte sie sagen: Ich elende Magd, weil mich GOtt ansieht, als ein armes, verlassenes Kind, so will ich den Sohn davon nennen, daß mich GOtt angesehen hat. Und siehe, wie Fleisch und Blut regieret. Denn das ist ein recht fleischlicher Gedanke, da\footnote{Erlanger: Fleisches Gedanken, daß.} sie spricht: „Nun wird mich mein Herr lieb haben“, denn ich bin ja die erste Mutter, bringe den ersten Sohn. Sie meint, sie wolle auch Frau sein; es fehlt aber. Den andern Sohn heißt sie Simeon, und spricht: „Der HErr hat gehöret, daß ich feindselig oder gehasset bin.“ Das ist nun fein, daß sie also kann sagen: GOtt hat mir den Sohn auch gegeben. Das sind Worte eines großen Glaubens, als spräche sie: Ob ich gleich der Aschenbrödel bin, so sieht er doch herab, und hört meine Noth.

20. Als sie nun den dritten Sohn dazu kriegt, heißt sie ihn Levi, das ist, Zugethan oder Zuhalter. Als sollte sie sagen: Jakob hält sich zu Rahel, und wohnet bei ihr, aber nun wird er sich zu mir auch halten. Sie hätte ihn gerne von Rahel gerissen, daß sie möchte die Liebste sein. Das wird nun jene auch verdrossen haben, daß sie nicht auch solche Kinder haben. Darauf kriegt sie noch einen Sohn, der heißt Juda, ein Dankbarer, darum, daß sie sprach: „Nun will ich dem HErrn danken“ oder bekennen, und nachsagen, was er gethan hat.

21. Nun siehe, wie die Schrift handelt. Was geht mich’s an, wie viel die Väter Kinder gekriegt, und wovon sie ihnen Namen gegeben haben? Müssen wir eben das alles hören und sehen, wie das ganze Buch von Kinderzeugen und Gebären redet? Oder, was soll ich davon hören, wie die Frauen auch hochmüthig gewesen sind, und jegliche wollte die Beste sein? Es ist alles, das ich gesagt habe: GOtt gefällt nichts denn sein eigen Werk. Also siehst du, was er hält von Kindertragen, daß er so viel läßt davon schreiben, also, daß da beide, sein Werk und Wort, uns offenbart ist; welches die Welt nicht sieht, ja, für lauter Schimpf und Lächer(e)i hält, und sich schämt, davon zu schreiben. Aber GOtt hat Lust dazu, daß er sich so hernieder lasse und von dem Dreck solche Ding mache, und was die Welt für köstlich hält, läßt im Dreck liegen.

22. Laß es nun Narrheit sein, Kinder zeugen, und uns mit den heiligen Vätern Narren bleiben, und auf GOttes Werk warten, wozu er\footnote{Walch und die Erlanger: es.} uns geschaffen hat, daß wir darinne bleiben, und für sein Werk halten, und groß achten, ob es schon vor der Welt schimpflich und ungeachtet ist. [Es] ist uns genug, daß es ihm wohlgefällt; wir werden es nicht besser machen, denn er es gemacht hat. So haben wir nun vier Söhne des heiligen Patriarchen, von den andern wollen wir nun auch hören. Die geistliche Deutung wollen wir sparen, bis wir baß hinein kommen, und diese Historie gar haben.
\quellseite{240}

\section{Das dreißigste Capitel}
\begin{schriftzitat}
V. 1—13. Da Rahel sahe, daß sie Jakob nichts gebar, neidete sie ihre Schwester, und sprach zu Jakob: Schaffe mir auch Kinder, wo nicht, so sterbe ich. Jakob aber ward sehr zornig auf Rahel, und sprach: Bin ich denn an GOttes Statt, der dir wehrt deines Leibes Frucht? Sie aber sprach: Siehe, da ist meine Magd Bilha, lege dich zu ihr, daß sie mir auf meinen Schooß gebäre, und ich doch durch sie erbauet werde; und gab ihm also Bilha, ihre Magd, zum Weibe. Und Jakob legte sich zu ihr. Also ward Bilha schwanger, und gebar Jakob einen Sohn. Da sprach Rahel: GOtt hat meine Sache gerichtet, und meine Stimme erhöret, und mir einen Sohn gegeben; darum hieß sie ihn Dan.\footnote{Randglosse: Dan heißt Richter.} Abermal ward Bilha, Rahels Magd, schwanger, und gebar Jakob den andern Sohn. Da sprach Rahel: GOtt hat es gewandt mit mir und meiner Schwester, und ich nehme überhand, und hieß ihn Naphtali.\footnote{Randglosse: Naphtali heißt verwechselt, angewandt, umgekehrt, wenn man das Widerspiel thut, Ps. 18, 27.: „Mit den Verkehrten verkehrest du dich.“} Da nun Lea sahe, daß sie aufgehöret hatte mit Gebären, nahm sie ihre Magd Silpa, und gab sie Jakob zum Weibe. Also gebar Silpa, Lea’s Magd, Jakob einen Sohn. Da sprach Lea: Rüstig, und hieß ihn Gad. Darnach gebar Silpa, Lea’s Magd, Jakob den andern Sohn. Da sprach Lea: Wohl mir, denn mich werden selig preisen die Töchter, und hieß ihn Asser.
\end{schriftzitat}

1. Im vorigen Capitel haben wir genug gehört, wie die arme Lea mußte der Aschenbrödel sein, Rahel aber die rechte Frau, und wie GOtt darin angezeigt hat, wie er das hoch hebt, was wir gering achten, so gar, daß er das Urtheil fest hält, auch unter den Heiligen. Denn Jakob war ein frommer Mann, voll Geistes; noch war Lea nichts Sonderliches bei ihm, Rahel aber war die rechte Frau, von der er viel hielt. In solchem fleischlichen Sinne läßt ihn GOtt noch gehen; aber läßt ihn auch getrost fehlen, und richtet viel anders denn Jakob.

2. Darum müssen wir uns, wenn wir gleich heilig sind, deß versehen, daß Gott nicht allewege das empor hebt, das wir empor heben, und verachtet, das wir verachten. Er wird’s uns nicht überheben, weil er es auch seiner Mutter nicht überhoben hat, und allen Aposteln, welche auch oft gute Gedanken gehabt, und gemeint, es wäre recht, und sollte so hinaus gehen; aber er macht es anders. Solches thut er noch immerdar, daß ja das Recht bleibe, daß er unsern Sinn breche. Darum ist dieses auch geschehen, daß er Leam ansieht, und fruchtbar macht, die andere aber nicht, wiewohl Jakob meint, sie sollte die rechte Mutter werden, bis so lange, daß sie sich auch demüthigt.

3. Hier sehen wir abermal ein langes wunderliches Capitel, welches, wenn es die Vernunft auf das tiefste ansieht, so ist es doch eitel Narrenwerk, und schier verdrießlich, daß man mit dem Viehe soll umgehen, und so viel davon halten, das doch nichts redet denn von Ziegen und Schafen und von Kinderzeugen; wie könnte er es doch närrischer vorlegen, denn das ist? sonderlich, wenn es soll von heiligen Leuten geschrieben sein, als diese sollen sein.

4. Aber ich kann nicht dawider, da liegt es; wer es besser\footnote{Erlangen: heißen.} kann, der thue es, wir wollen zusehen. Wenn St. Hieronymus davon sollte schreiben, so würde er sagen, man müßte nicht bei der schlechten Historie bleiben, sondern nur den geistlichen Verstand darin suchen; das wollen wir lassen. Ich habe also gesagt, man solle also lesen und hören, daß man in allen GOttes Werken nicht mehr ansehe denn seinen Willen, Augen, Ohren, und alle Sinne zuthun, und nicht weiter fragen. Ist es GOttes Werk und Wille, so bete es an, und trage es auf den Händen. Er hätte auch gar wohl andere Dinge gewußt zu beschreiben, wenn es ihm gefallen hätte.

5. Was wollen wir aber dennoch dazu sagen, daß Jakob vier Weiber nimmt zuhauf, zwo Schwestern, und zwo Mägde? Ist es nicht ein hurischer Abenteurer, daß ihm nicht genüget an zweien, sondern noch zwo dazu haben muß? Ich habe im vorigen Capitel [§ 8] gesagt, daß keine Keuschheit wohl beweisen ist in dem, daß er siebenzig Jahre gelebt, und keines Weibes schuldig
\quellseite{241}
ist worden; jetzt aber, als er ein alter Narr ist, und natürlich nicht viel mehr taugt, Frucht zu zeugen, nimmt er vier Weiber. Die erste Antwort ist die: Weil es die Schrift nicht tadelt, sondern so mit Fleiß beschreibt, und zeigt, wie ihn GOtt gesegnet, und von den vier Weibern zwölf Patriarchen gibt, gebühret uns das Maul zuzustopfen, und [zu] sagen: Es gefällt mir wohl, weil es ihm gefallen hat.

6. Ich halte aber, sollten viel heilige Leute (als Hieronymus, Augustinus, Hilarius, und auch wohl St. Peter) zu ihren Zeiten solche Werke gesehen haben\footnote{So Walch und die Erlanger. „haben“ fehlt in der Jenaer. Im Original und in der Wittenberger: geschehen wäre.} von einem solchen Manne, wäre [es] ihnen zu hoch gewesen, daß sie sich nicht sollten davor entsetzt haben, es wäre denn, daß ihnen GOttes Geist sonderlich eingegeben hätte, und gesagt, es wäre recht. Denn es ist ja ein hoch Werck, über alle Vernunft, daß es auch wohl den lieben Heiligen hätte mögen ärgerlich sein. Es ist nun in die Welt gebracht, daß niemand sagen darf, es sei unrecht. Aber wenn es vor unsere Augen gestellt würde, so würde ich mit allen andern darüber zum Narren.

7. Also macht GOtt die Welt (und auch zu Zeiten die hochverständigen geistlichen Leute) zu Narren mit solchen Werken, daß man nicht weiß, was er im Sinne hat. Darum ist sein Regiment wunderlich, wie St. Paulus sagt aus vollem Geiste. Als er es lange hin und her ansah, läßt er es fallen, und fährt mit dem Spruche herein [Röm. 11, 33.]: „O welch eine Tiefe des Reichthums beide der Weisheit und Erkenntniß GOttes! wie gar unbegreiflich sind seine Gerichte, und unerforschlich seine Wege.“ Summa, es ist gar dahin gerichtet, daß er gerne GOtt wollte sein, und wir dafür hielten, was er machte, daß [es] recht wäre. Aber das thun wir nicht, sondern nehmen Gesetze vor uns, und schließen: Wer das nicht thut, ist unrecht. Wie aber, wenn er sagt: Wie denn, wenn du das Gesetz nicht recht verstehst, und ich wollte es anders machen, denn du willst oder begreifst?

8. Das ist die erste Ursache, daß der\footnote{„er“ fehlt in der Jenaer.} Moses solch Ding beschreibt, also, daß er es nicht häßlicher könnte beschreiben. Er hat es gethan, darum ist es recht. Wie denn, wenn wir es auch thäten, wäre es nicht auch recht? Oder ist es uns zum Exempel geschrieben, daß wir auch also thun sollen? Nein, denn ich habe ja oft genug gesagt, man solle nicht auf die Werke, sondern auf die Person Jakobs sehen. Bist du auch ein Fetel,\footnote{Fetel = Jakob.} so thue es auch. GOtt will es nicht leiden, daß wir die Werke unter die Hände nehmen; darum stellt er sich so närrisch mit der Heiligen Werken, uns zu wehren, daß wir nicht darauf fallen. Siehe vor auf den Glauben, Geist und Wort, so ihm gegeben ist. Wenn du nun eine solche Person bist, so thue in GOttes Namen, was dir vorkommt.

9. Denn siehe darauf, wie er thut. Sucht er es doch nicht, und ist kein Vorwitz da, sondern Rahel spricht zu ihm: Gib mir auch Kinder, oder ich muß sterben. Die gute Metze hätte gerne Kinder gehabt, sonderlich weil sie sah, daß ihre Schwester so fruchtbar war, und wäre ihr große Schande und Herzeleid gewesen, daß sie ohne Kinder sollte bleiben. Sie stand auf der Predigt, die sie von Jakob oft gehört hatte, wie durch seinen Samen allerlei Volk sollte gesegnet werden; das mußte und glaubte sie; darum hätte sie lieber begehrt zu sterben, denn ohne Frucht zu bleiben. [Sie] hat müssen also denken: Ich bin die rechte Frau; was thut mir denn GOtt für eine Schalckheit, daß ich nicht werth bin, die Mutter zu sein? Wie es noch heutiges Tages einem christlichen Herzen möchte zu Muthe sein, wenn es ihm so ginge, daß es würde erschrecken, und denken: Auweh! ich sehe wohl, daß GOtt das Wort nicht hat von mir geredet, und soll mich nichts angehen; er hält sich zu meiner Schwester, und verachtet mich. Also wird sie das Wort gemüthigt haben in die Asche hinein, und ihr gar wehe gethan, daß sie muß sagen: Kriege ich nicht auch Kinder, so muß ich vor Leide sterben. Da ist ihre Freude rein hinweg genommen.

10. Dazu spricht der Text, welches noch härter ist, wie sie ihre Schwester geneidet hat, und auf sie verdrossen. Da ist noch eben stark Fleisch und Blut mit unter gelaufen, wie in allen Christen. Sie hat sich gefürchtet, daß GOtt nicht nach ihr fragte; auch hat sie Sorge gehabt (als die Weiber thun), sie möchte nicht lange die Liebste bleiben. Solche närrischen, fleischlichen Theidinge läßt GOtt noch in ihnen stecken, daß man sieht, wie sie alle gebrechlich sind gewesen,
\quellseite{242}
wie mir alle. Wiederum wird Jakob auch läunisch, ja, wie der Text sagt, sehr zornig auf sie, und spricht: Meinst du denn, daß ich an GOttes Statt sei, und soll dir Kinder geben, wenn er dir keine gibt? Hier sieht man wohl, daß er ihnen oft gepredigt hat, und wie sie GOtt und sein Wort wohl erkannt haben. Will aber so viel sagen:

11. Du weißt je wohl, was GOtt ist, daß keines Menschen Kraft ist, Kinder zu tragen, sondern GOttes Hand und Gewalt. Was kann ich dazu, wenn er dir keines geben will? Das ist zwar ein kleiner Trost, den er ihr da gibt; aber das verdrießt ihn, als einen heiligen Mann, daß sie so unvernünftig aus Ungeduld herfährt, und redet, als habe sie GOttes vergessen, [es] sei Jakobs Schuld, daß sie keine Kinder hat.

12. Was thut sie aber, als sie sieht, daß sie nicht fruchtbar wird? „Da ist meine Magd (spricht sie), lege dich zu ihr, daß sie mir auf meinen Schooß gebäre, und ich doch durch sie erbauet werde.“ Das ist zu der Zeit der Brauch gewesen unter den Heiden, daß Knechte und Mägde eigene Leute waren; also, daß sie die Herren um Geld kauften für ihr eigen Gut, wie Pferde und Kühe. Also hat es auch Moses im Gesetze verfaßet, wenn Knechte und Mägde, die man zusammen gegeben hatte, Kinder zeugten, daß sie alle des Herrn eigen waren. Wie auch noch der Brauch sein soll in des Türken Landen. Darum sagt Rahel allhier: Die Magd ist mein, lasse sie mir auf meinen Schooß gebären. Denn, weil sie die Magd ihrem Manne gibt und zulegt, so ist alles, was sie gebiert, nach dem Rechte ihr eigen; [sie] will ihr aber keinen fremden und eigenen Mann geben, sondern von ihrem Jakob Kinder haben; läßt sie Magd bleiben, und sie bleibt Frau. Es ist ein wunderlich Recht, nicht des Papsts, sondern GOttes. Hat sie aber auch Recht, daß sie ihren Mann hinweg gibt? Ohne ihren Willen hätte es Jakob nicht gethan; weil er nun zuhört, und die Magd behält, so muß es recht sein; sonst hätte er nicht dazu bewilligt. Es sind eitel ernste Sachen, aus keinem Vorwitz geschehen.

13. Dieselbe Magd zeugt nun der Frau zween Söhne; den ersten heißt sie Dan, das ist, ein Richter. [Sie] läßt ihm nicht die Magd einen Namen geben, noch sich des Kindes annehmen, sondern will selbst Mutter sein, auf daß sie durch die Magd erbauet werde, und ein Haus habe, daß man könnte sagen: Das ist Rahels Kind. Darum nennt sie ihn nun Dan, und spricht: „GOtt hat meine Sache gerichtet.“ Als sollte sie sagen: „Ich bin in der Schmach gelegen; nun aber bin ich auch glücklich, daß man mich auch Mutter kann heißen; wiewohl durch ein ander Weib, doch durch die, die mein ist. So genau sucht sie es, daß sie nur auch fruchtbar möge sein. Denn im alten Testamente war es die größte Schande, wo ein Weib nicht fruchtbar war; [es] ließ sich ansehen, als wäre ihr GOtt nicht gnädig, daß er sie des Segens nicht theilhaftig machte. Den andern Sohn von der Magd nennt sie auch selbst, und heißt ihn Naphtali, das heißt umgekehrt oder gewechselt. [Sie] will also sagen: GOtt sei gelobt! Meine Schwester hat aufgehört, sie steht stille, und ich fahre fort und stehe mit großem Ehren, nun will ich wohl oberste Frau bleiben; ist abermahl ein menschlicher und weibischer Gedanke.

14. Wiederum hat Lea auch eine weibische Anfechtung, will der Schwester nichts zuvor geben, und, weil sie aufgehört hat zu gebären, nimmt sie auch ihre Magd, und gibt sie Jakob. Die zeugt nun einen Sohn, den hieß sie Gad, und sprach: „Rüstig“; als sollte sie sagen: Es geht noch fein rüstig und frisch zu. Denn Gad heißt rüstig oder reisig zum Streite; davon kommt Gedud, ein Kriegsknecht. Darnach noch einen, der heißt Asser, selig; denn „nun werden mich die Weiber selig sprechen“, sagt sie, daß ich von Jakob so viel Kinder kriege.

15. Das ist die Historie von den zwei Mägden, die den zwey Frauen müssen Kinder zubringen, und doch Mägde bleiben; wiewohl sie Jakobs Eheweiber sind. Ich habe zuvor [§ 12] gesagt, wie das recht sei, daß der große heilige Mann dennoch die Weiber so hart hält, und nach gestrengem Rechte fährt, läßt sie nicht genießen, daß sie seine Weiber sind, daß er sie los gäbe und frei machte, welches je billig und freundlich anzugehen wäre. O wie hat GOtt das Volk im Zwange gehalten, daß es nicht zu frei und muthwillig würde. Wäre auch wohl gut, daß solcher Zwang jetzt auch wäre, denn wir wollen allesamt Junker sein; daher auch das Gesinde den Uebermuth und Muthwillen treibt, daß niemand mit ihnen umkommen kann; da wird nimmermehr kein recht Regiment aus. Es ist wohl unfreundlich anzusehen, daß man die Leute eigen macht, wie das Vieh. Ist es doch auch
\quellseite{243}
greulich, daß man Liebe an Galgen hängt; man muß aber die Leute zwingen und zähmen, man hat dennoch genug zu schaffen, daß man Frieden habe. Denn es ist unmöglich, daß die, so nicht GOttes Kinder, und dazu ungläubig sind, sollten etwas Gutes thun, sonderlich wenn man ihnen Raum und Willen läßt. Derhalben ist jetzt kein Regiment, Zucht, noch Ehre mehr, läuft alles von einem Lande zum andern, und richtet alle Büberei an.

16. Darum haben es nun die heiligen Leute also gehalten, daß [es] nicht gut ist, solchen [Leuten] Raum und Zaum zu lassen; sie würden sonst zu stolz, und unbändig. Auch, ob sie Jakob gleich hätte wollen frei machen, hätte es nicht wohl gestanden, weil es des Landes Brauch also war; darum wollte er kein neues Recht aufbringen. Man muß je eine Ordnung unter den Leuten halten, daß nicht jedermann thue, was er will; wie es jetzt zugeht, daß Herren und Frauen dem Gesinde mehr nachgeben müssen denn ihren Kindern, und oft mehr Recht lassen, denn sie selbst haben. Darum ist es ein wohlgeordnetes, köstliches Regiment gewesen. Es ist wohl nicht fein, daß eines des andern Gefangener sei, aber noth ist es, äußerlich, weltlich Regiment zu erhalten. Also hätten es die Väter ihrethalben wohl lassen gehen; aber weil man das Gesinde nicht wohl anders konnte regieren und im Zwange behalten, haben sie es nach Landes Sitten mit gehalten, und dabei lassen bleiben. Auch ist mir nicht Zweifel, daß diese zwo Mägde fromme und heilige Weiber gewesen sind. [Es] ist eine große Ehre, daß sie die zwölf Patriarchen haben helfen tragen. [Es] folgt nun weiter in der Historie.

\begin{schriftzitat}
V. 14—24. Ruben aber ging aus zur Zeit der Weizenernte, und fand Alraun auf dem Felde, und brachte sie heim seiner Mutter Lea. Da sprach Rahel zu Lea: Gib mir der Alraun deines Sohnes ein Theil. Sie antwortete: Hast du nicht genug, daß du mir meinen Mann genommen hast, und willst auch die Alraun meines Sohnes nehmen? Rahel sprach: Wohlan, laß ihn diese Nacht bei dir schlafen um die Alraun deines Sohnes. Da nun Jakob des Abends vom Felde kam, ging ihm Lea hinaus entgegen, und sprach: Bei mir sollst du liegen, denn ich habe dich erkauft um die Alraun meines Sohnes. Und er schlief die Nacht bei ihr. Und GOtt erhörete Lea, und sie ward schwanger, und gebar Jakob den fünften Sohn, und sprach: GOtt hat mir gelohnet, daß ich meine Magd meinem Manne gegeben habe, und hieß ihn Isaschar. Abermal ward Lea schwanger, und gebar Jakob den sechsten Sohn, und sprach: GOtt hat mich wohl berathen, nun wird mein Mann wieder bei mir wohnen, denn ich habe ihm sechs Söhne geboren; und hieß ihn Sebulon. Darnach gebar sie eine Tochter, die hieß sie Dina. GOtt gedachte aber an Rahel, und erhörete sie, und machte sie fruchtbar. Da ward sie schwanger, und gebar einen Sohn, und sprach: GOtt hat meine Schmach von mir genommen, und hieß ihn Joseph, und sprach: GOtt wolle mir noch einen andern Sohn dazu thun.
\end{schriftzitat}

17. Hier [er]hebt sich abermal ein Aergerniß, daß sich die zwo heiligen Schwestern, Rahel und Lea, neiden und hassen um den Mann Jakob, auf daß man sehe, wie sie auch gebrechlich und Menschen gewesen sind, uns zu Trost, die wir Sünder sind. Aber daneben wird angezeigt, wie die eheliche Liebe sei die größte und höchste Liebe auf Erden, daß sie um des ehelichen Gemahls willen verläßt, verachtet, ja, auch haßt und meidet Schwester, Bruder, Vater und Mutter. Nun, die Lea klagt, daß Rahel nehme ihr den Mann; will damit ihr Recht nicht vergeben, sondern auch Weib sein; ja, sie wollte wohl gerne allein und die Vornehmste sein, denn sie spricht: „Meinen Mann nimmst du mir“, gerade als wäre sie das rechte Weib. Aber es hilft nicht, sie mußte Nebenfrau sein, und Jakob bei Rahel lassen wohnen täglich und am meisten. Doch thut ihr Rahel eine Freundschaft, und läßt Jakob zu Nacht bei Lea liegen. Solch Geheiß, Neid, Klagen und wiederum Versöhnen wird sich oft haben begeben, daß Jakob genug mit ihnen hat zu schaffen gehabt. Es wird aber einmal hier angezeigt, zu deuten, wie das Haushalten sei\footnote{In den Ausgaben: habe.} zugegangen, wie mit allen andern frommen GOttes Kindern; wie auch Petrus [1. Ep. 3, 7.] bekennt, daß ein Weib schwach sei, und St. Paulus [1 Cor. 7, 10.] die Ehelichen lehrt, sich selbst versöhnen. Denn sie wissen, daß [es] ohne Gebrechen, Zorn, Ungeduld und dergleichen, nicht abgeht im ehelichen Leben, viel mehr zu der Zeit, da viel Weiber Einen Mann haben gehabt.
\quellseite{244}

18. Die Allrun, welche auf hebräisch heißen Dubaim, halten etliche, es seien Violen gewesen. Aber es mögen nicht die blauen oder gelben Violen sein, denn hier sagt der Text, Ruben habe sie in der Weizenernte gefunden, wie auch Salomon in seinem Hohenliede [Cap. 7, 13.] spricht: „Dudaim, die geben ihren Geruch.“ Aus dem scheint es, daß es sind seine Blumen gewesen, die um die Zeit der Ernte wohl riechen, als bei uns die Nelken\footnote{„als bei uns die Nelken“ fehlt in der Wittenberg.}, Lilien und Rosen. Aber was [es] für Blumen sind, weiß man nicht gewiß. Daß aber alle Dolmetscher haben Allrun daraus gemacht, welchen ich gefolgt habe, weil sonst keine gewisse Blume zu nennen war, halte ich für nichts. Was sollte Allrun riechen, und so wohl gefallen den Weibern? Ich nehme Nelken\footnote{Jenaer: „Nelichen“; Wittenberger: „nelglichen“; Erlanger: „Nelichen“.}, oder weiße Lilien dafür. Es scheint, weil niemand gewiß gewußt hat, was Dudaim heißt, daß etwa ein Jude habe unser gepottet, und gesagt, es heißt Allrun; wie St. Hilarius im Worte Hosianna, und viele andere mehr auch betrogen und gespottet sind.

19. Was ist es aber vonnöthen, solche Geschichte zu schreiben, wie sich Weiber um Blumen hadern? Keine andere Noth, denn daß GOtt, wie oft gesagt, nicht viel nach unsern Werken fragt, sondern seiner lieben Heiligen Werke gefallen ihm, wie gering und nichtig sie immer sind vor der Vernunft. Denn hier siehst du dennoch, neben dem geringen Werke, wie volles Glaubens die Lea ist, daß GOtt ihr Gebet erhört, und macht sie schwanger; und sie lobt und dankt und GOttes Gnade erkennt, als spiele GOtt mit diesen Blumen bei seinen Kindern; und richtet gleichwohl dadurch so große Dinge aus, daß zween Erzväter und die Tochter Dina davon kommen, mit großem Lobe und Ehren GOttes. Was aber diese Geschichte und Blumen bedeuten, lasse ich jetzt anstehen bis zur andern Zeit, ist auch aus den vorigen Deutungen wohl zu erdenken, wer Lust dazu hat.

\begin{schriftzitat}
V. 25—43.

Da nun Rahel den Joseph geboren hatte, sprach Jakob zu Laban: Laß mich ziehen und reisen an meinen Ort und in mein Land. Gib mir meine Weiber und meine Kinder, darum ich dir gedienet habe, daß ich ziehe; denn du weißest, was für einen Dienst ich dir gethan habe. Laban sprach zu ihm: Kann ich nicht Gnade vor deinen Augen finden? Ich spüre, daß mich GOtt gesegnet hat um deinetwillen. Stimme das Lohn, das ich dir geben soll. Er aber sprach zu ihm: Du weißest, wie ich dir gedienet habe, und was du für Vieh hast unter mir. Du hattest wenig, ehe ich her kam, nun aber ist es ausgebreitet in die Menge, und der HErr hat dich gesegnet um meinetwillen; und nun, wann soll ich auch für mein Haus schaffen? Er aber sprach: Was soll ich dir denn geben? Jakob sprach: Du sollst mir nichts überall geben, sondern so du mir thun willst, das ich sage, so will ich wiederum weiden, und hüten deiner Schafe. Ich will heute durch alle deine Herde gehen, und sondere du von dannen alle fleckete und bunte Schafe, und alle schwarze Schafe unter den Lämmern. Was nun bunt und flecket unter den Ziegen\footnote{Handglosse: Ziegen. Du mußt hier dich nicht irren, daß Moses das kleine Vieh jetzt Ziegen, jetzt Lämmer, jetzt Böcke heißt, wie dieser Sprache Art ist. Denn er will so viel sagen, daß Jakob habe alles weiße, einfarbige Vieh behalten, und alles bunte und schwarze Laban gethan. Also nun bunt von dem einfarbigen Viehe käme, das sollte sein Lohn sein. Es ward Laban froh und hatte die Natur für sich, daß von einfarbigen nicht viel bunte natürlich kommen. Aber Jakob half der Natur mit Kunst, daß die einfarbigen viel bunte trugen. Durch diese Geschicht ist bedeutet, daß durch das Evangelium werden die Seelen von den Gesetztreibern und Werkheiligen abgeführt, darinnen sie bunt, sprenklicht und flecket, das ist, mit mancherlei Gaben des Geistes geziert werden. Röm. 12, 6. 1 Cor. 12, 4., daß unter dem Gesetze und Werken nur die Untüchtigen bleiben. Denn Laban heißt weiß oder gleichend, und bedeutet der Gleisner Haufen in den schönen Werken, auch göttliches Gesetzes.} sein wird, das soll mein Lohn sein. So wird mir meine Gerechtigkeit zeugen heute oder morgen, wenn es kommt zu meinem Lohn vor dir; also, daß, was nicht fleckt und bunt ist unter den Ziegen, und was schwarz ist unter den Lämmern, das sei ein Diebstahl bei mir. Da sprach Laban: Siehe da, es sei, wie du gesagt hast, und sonderte des Tages aus sprenkliche und bunte Böcke, und alle fleckete und bunte Ziegen, wo nur was Weißes daran war, und alles, was schwarz war unter den Lämmern, und that es unter die Hand seiner Kinder, und machte Raum dreier Tagereisen weit, zwischen ihm und Jakob. Also weidete Jakob die übrige Herde Labans. Jakob aber nahm Stäbe von grünen Weiden, Haseln und Kastanien, und schälte weiße Streifen daran, und legte die Stäbe, die er geschälet hatte, in die Tränkrinnen, vor die Herden,
\quellseite{245}
die da kommen mußten zu trinken, daß sie empfangen sollten, wenn sie zu trinken kämen. Also empfingen die Heerden über den Stäben, und brachten sprencklichte, fleckte und bunte. Da scheidete Jakob die Lämmer, was nicht bunt, und alles, was schwarz war, und that sie zu Haufen, unter die Heerde Laban, und machte ihm eine eigene Heerde, die that er nicht zu der Heerde Labans. Wenn aber der Lauf der Frühlinge Heerde war, legte er die Stäbe in die Rinnen vor die Augen der Heerde, daß sie über den Stäben empfingen. Aber in der Spätlinger Laufe legte er sie nicht hinein. Also wurden die Spätlinge des Laban, aber die Frühlinge des Jakob. Daher ward der Mann über die Maaßen reich, daß er viel Schaafe, Mägde und Knechte, Kameele und Esel hatte.
\end{schriftzitat}

20. Da beschreibt Moses, wie Jakob mit seinem Schwäher handelt, und ihn betrüget um die Schaafe. Hier müssen wir abermal antworten denen, so es nach dem Fleische ansehen, als habe Jakob schändlich gehandelt mit seinem Herrn. Der Text ist aber etwas finster der Sprache halben, darum wollen wir die Meinung auf unser Deutsch sagen, wie es zugegangen ist. Zum ersten hat Jakob gedient um Rahel sieben Jahre, wie gehört ist. Und als die aus waren, betrog ihn Laban, und gab ihm die Lea. Danach mußte er noch sieben Jahre dienen um Rahel, die gab er ihm für seinen Lohn. Da nun die vierzehn Jahre um waren, und er keinen Lohn verdient hatte, denn Hülle und Fülle, Futter und Decke, und doch seines Herrn Gut merklich gebessert, wie er selbst bekennt, werden sie der Sache eins des Lohnes halben, und spricht Jakob: Wir wollen lassen bezähmen\footnote{bezähmen lassen = schalten und walten lassen.} Gott und die Natur, was mein und dein sein solle. Da hast du einen Haufen Schaafe, und thue heraus, was bunt ist, und gib mir die einfarbigen. Das schreibt Moses mit umschweifenden Worten; die bunten Schaafe heißt er halb weiß und schwarzgefleckte oder sprencklichte, die weiß oder schwarze Male und Flecken, doch sonst Eine Farbe haben. Das heißen wir alles bunt. Danach soll unser Gedinge also gehen: Die weißen, einfarbigen Schaafe, und was dieselben für einfarbige bringen, soll auch dein sein; was aber Buntes geboren wird, soll mein Lohn sein.

21. Also machte er nun zwo Heerden, nahm zu sich die einfarbigen, und was bunt war, that er den Kindern. Und gibt es also vor: Nun habe ich die einfarbigen Schaafe allein, was nun von diesen bunte\footnote{Wittenberger und Jenaer: bunten.} kommt, soll mein sein. Das hört Laban gerne, meinte, Jakob hätte es wohl getroffen, wollte sich selbst betrügen; denn es war nicht nach gemeinlichem, natürlichem Laufe, daß von einfarbigem Viehe viel Buntes sollte kommen. Aber Jakob braucht einer List, läßt Laban die bunten alle hinnehmen, und unter seiner Kinder Hände thun, drei Tagereisen von dannen; er aber bleibt mit den einfarbigen, nimmt aber Stecken oder Stäbe, und schält sie halb mit Streifen. Da sich nun die Schaafe züchten und laufen, legt er ihnen die Stäbe vor, und als sie die bunten Stäbe ansehen, empfingen sie auch bunte und sprencklichte Schaafe. Mit dieser Kunst brachte er zuwege, daß von den einfarbigen so viel bunte und wenig einfarbige kamen.

22. Da aber das halbe Jahr um war, und Laban sahe, daß Jakobs Heerde groß ward, und seiner weniger: Ei, sprach er, es taugt nicht also, wir wollen wechseln; ich will diesmal die bunten nehmen, und behalte du die einfarbigen. Jakob war es zufrieden, und ließ [es] nach der Natur bleiben, legte ihnen keine bunten Stäbe vor. Nach dem Jahre sahe Laban abermal, daß Jakob die meisten Schaafe kriegte; da wandelt er es abermal mit dem Lohne, daß Jakob wieder die bunten behielte. Der Wechsel geschah nun zehnmal nach einander. Denn Laban war geizig, wollte allezeit die meisten Schaafe haben; aber Jakob brauchte der Kunst, daß es ihm immerdar mehr trug denn Laban. Das ist die Meinung des ganzen Textes, auf das deutlichste gesagt.

23. Nun ist die Frage: ob Jakob recht gethan habe, daß er so listiglich handelt? bringt also die meisten Schaafe davon, daß [es] jener nicht weiß; denn er hätte es ihm sonst nicht gestattet. [Es] ist auch wohl zu denken, daß Jakob fromm Gesinde gehabt habe, daß sie ihn nicht verrathen haben. Denn es ist ohne Zweifel zu glauben, Jakob wird sein Haus also regiert haben, daß er GOttes Gebot und Wort gepredigt, und nichts Unrechtes im Hause gelitten hat. Was wollen wir nun dazu sagen, weil der Text so klar steht, wie er seinen Vortheil sucht, und dem Laban
\quellseite{246}
abbricht? Wenn es unser einer thäte, so gälte es gar nichts; oder, sollte es jetzt gelten, so wollten wir getrost stehlen. Also sind die Historien im Alten Testamente schier das Meiste alle, daß Vernunft müßte schließen, es sei nicht recht gehandelt. Warum schreibt er es denn also vor, als sei es recht und wohl gethan, dazu segnet und benedeiet ihn darüber?

24. Antwort: Daß er das Spiel getrieben habe aus des Heiligen Geistes Eingeben und Treiben, ist klar und unwidersprechlich aus dem folgenden Texte, da es GOtt bestätigt, und heißt ihn die Weiber, Gesinde und Vieh nehmen und davon ziehen. Wo nun GOttes Wort geht, soll man nicht fragen, ob es recht sei, sondern getrost hinan gehen; was er heißt, das soll recht sein. Darum, weil es hier bestätigt wird von GOtt und getrieben, soll man keine andere Ursache hören, denn daß es recht gethan sei, [er] auch Recht hatte zu dem Gute [auf] mancherlei Weise. Er hatte sieben Jahre gedient um die Töchter; er aber täuscht ihn, und hält ihn noch sieben Jahre auf, daß er nicht mehr kriegt, denn beide Töchter; gibt ihm nichts, damit er sich nähre. War es nicht Sünde und Schande, daß er ihm die Wezen so theuer gab, und gibt ihm nichts dazu, wie den Töchtern gebührt, denn was sie verdienen, als die Fremden und Mäagde, wie sie im folgenden Capitel [V. 15.] klagen?

25. Das war nun zweierlei groß Unrecht: Daß ihm Jakob so lange umsonst dienen muß, und er ihm beide Töchter anhängt und nichts mitgibt. Wenn da Recht Recht wäre gewesen, hätte man ihn sollen zwingen, daß er die Töchter versorgete, und gäbe ihnen zu, was sich gebührte. Jakob war fromm, wollte nicht had[e]rn, sondern litt es mit Geduld vierzehn Jahre lang. Das thue ihm nach, so wollen wir dich auch lassen nehmen und stehlen. Dazu ist das auch eine Ursache, daß er ihm schuldig war für die Zeit lang, da er mit seinen Weibern und Kindern diente. So rechnet er es ab. Wenn es gedingt wäre gewesen, hätte er ihm jährlich eine bestimmte Zahl Schafe müssen geben; das hätte die ganze Zeit über einen großen Haufen getragen, und sich jährlich gemehrt. Was thut er aber? Wiewohl er solch Recht hat vor GOtt, thut er es dennoch nicht; sondern harrt, bis ihm GOtt Ursache dazu gibt, und gibt es ihm selbst ein.

26. Auch hatten sie, die Töchter, sammt ihren Kindern gut Recht, solches zu sich zu bringen, weil eine jegliche nun sieben Jahre neben ihm gedienet hatte; wie sie hernach [Cap. 31, 15.] sagen: „Er hat uns gehalten wie die Fremden, denn er hat uns verkauft und unsern Lohn verzehrt.“ Sie waren ehelich, hatten Kinder und Gesinde; darum war es Zeit, auch auf ihr Haus zu denken, weil er ihnen doch nichts geben wollte. Also zeigt Moses an, wie der fromme Jakob gar viel Unrecht hat müssen leiden, und immer im Kreuze gelebt, und doch still geschwiegen, sich deß getröstet, daß ihm GOtt zugesagt hatte, er wollte ihn nicht lassen. Heißt das nicht gelassen, daß er so lange mit saurer Arbeit dienen muß, und dazu Ueberlast und Unrecht leiden? O wie könnten wir so übel leiden! er aber hat daraus gelernt, wie ein wunderbarer Gott es wäre.

27. Ueber das handelt der Laban noch weiter schälklich mit ihm, in dem, daß er ihm zehnmal seinen Lohn verwechselt, und ihm nicht gönnen will, was ihm Gott bescheert, so er doch durch ihn reich worden war; [er] machte es so lange und viel, daß [es] GOtt nicht mehr leiden wollte. Also muß man der Heiligen Werk ansehen, daß man sich nicht daran stoße. Es wäre noch nicht große Sünde, wenn einer jetzt so lange und schwerlich dienen sollte, und keinen Lohn kriegte, ob er mit solcher Weise und Meinung etwas von des Herrn Gute zu sich brächte. Es gebührt ihm doch, und ist vor GOtt sein. Kann er es nun füglichere Weise davon bringen, so hat er das Seine genommen. Also haben die Kinder Israel auch gethan, als sie aus Egypten zogen; wie im andern Buche Mosis [Cap. 11, 2. 12, 35.] steht. Als sie mit langer und harter Arbeit geschmächt waren, und zwo Städte gebauet und fest gemacht, und die Egypter auf das allerunbarmherzigste mit ihnen umhielten, hieß sie GOtt mit allem, was sie hatten, davon ziehen, und darüber allerlei silbern und gülden Gerathe von den Egyptern nehmen und entführen. Summa: Verdienter Lohn ist nicht des Herrn, sondern des Knechtes; kann er es hinweg bringen, daß [es] der Herr nicht weiß, so nehme er es immer hin. Denen,\footnote{Wittenberger und Jenaer: Denn.} die da Christen sind, darf man nicht viel Regeln vorschreiben, noch ihre Werke meistern; die andern aber gehen uns nichts an. Das sei dies Capitel; wollen nun die Historie vollends hinaus führen.
\quellseite{247}

\section[Das einunddreißigste Capitel]{Das einunddreißigste Capitel\footnote{In Buchwalds „Andreas Brachs handschriftliche Sammlung“ etc. begegnen wir S. XXXVIII der Notiz, daß aus den Nachschriften Stephan Roths in Zwickau hervorgehe, „daß Luther am Sonntag Judica (13. März) 1524 bei Capitel 31 stand“.}}

\begin{schriftzitat}
V. 1—10. Und es kamen vor ihn die Reden der Kinder Labans, daß sie sprachen: Jakob hat alle unsers Vaters Gut zu sich gebracht, und von unsers Vaters Gute hat er solchen Reichthum zuwege gebracht. Und Jakob sahe an das Angesicht Labans, und siehe, es war nicht gegen ihn, wie vorhin. Und der HErr sprach zu Jakob: Zeuch wieder in deiner Väter Land und zu deiner Freundschaft, ich will mit dir sein. Da sandte Jakob hin, und ließ rufen Rahel und Lea auf das Feld bei seine Heerde, und sprach zu ihnen: Ich sehe eures Vaters Angesicht, daß es nicht gegen mir ist wie vorhin. Aber der GOtt meines Vaters ist mit mir gewesen. Und ihr wisset, daß ich aus allen meinen Kräften eurem Vater gedienet habe. Und er hat mich getäuschet, und nun zehnmal meinen Lohn verändert. Aber GOtt hat ihm nicht gestattet, daß er mir Schaden thäte. Wenn er sprach: Die bunten sollen dein Lohn sein, so trug die ganze Heerde bunte. Wenn er aber sprach: Die sprenglichten sollen dein Lohn sein, so trug die ganze Heerde sprenglichte. Also hat GOtt die Güter eures Vaters ihm entwandt, und mir gegeben. Denn wenn die Zeit des Laufs kam, hob ich meine Augen auf, und sahe im Traume, und siehe, die Böcke sprangen auf die sprenglichten, flecketen und bunten Schafe.
\end{schriftzitat}

1. So haben wir gehört, wie Jakob ganzer zwanzig Jahre gedienet hatte, welches je ein langer Dienst war einem solchen Manne, der Weib und Kind hatte, und nichts Eigenes zu Lohn kriegte, ohne was ihm GOtt zuletzt bescherte. So haben auch die beiden Kinder viel müssen leiden, wiewohl sie seine eigenen Töchter waren. Rahel, wiewohl sie die Liebste war dem Jakob, doch hat sie auch müssen ein Aschenbrödel sein, gegen der Söhne Weiber. Das hat sieben Jahre gewährt, daß sie mußten Mägde sein, und kein Recht im Hause haben. Was manche böse Tücke und Stücke haben sie müssen leiden, und schweigen! Denn die andern Weiber sind Heiden gewesen, stolze Metzen, haben gelebt nach Fleisch und Blut, also, daß sie alle mit Jakob ein schwer, elend Leben geführt haben.

2. Doch ist er nicht aufgebrochen, hat auch nicht davon geeilt, sondern geharrt, bis ihm GOtt sagte, welcher ihn vorhin [Cap. 28, 15.] hatte heißen ins Land ziehen, und verheißen, er wollte ihn mit Frieden wieder heim führen. Darauf hat er nun täglich gewartet, [es] hat sich aber sehr lange verzogen.

3. Wer nun dies Capitel liest, der muß darauf Achtung haben, daß das Wort GOttes, so hernach steht, sei vorhin geschehen. Der Text hebt an, als habe er die zwei Weiber zuvor zu sich gerufen, ehe GOtt mit ihm redet; aber es ist eine gute Weile vor[her] geschehen, zu der Zeit, als sich das Gedinge angefangen hat mit Laban, im vierzehnten oder fünfzehnten Jahre; wie er hernach [V. 11. 12.] zu verstehen gibt, als er sagt: „Der Engel GOttes kam zu mir im Traume, und sprach: Hebe auf deine Augen, und siehe, die Böcke springen auf die sprenglichten, flecketen und bunten Schafe“ etc. Die lieben Heiligen, was sie gethan haben, das haben sie gethan im Worte GOttes; was sie das nicht geheißen hat, haben sie nicht gethan. Darum, als das Wort hier kommt, und heißt ihn wieder hinweg ziehen, folgt er ihm, und macht sich auf, geht immer dem Worte nach, ist aber voll Leidens und Kreuzes.

4. Es kamen die Reden der Kinder Labans vor ihn, spricht Moses, daß sie sprachen: „Jakob hat alle unsers Vaters Gut zu sich bracht.“ Welch ein neidisch, geizig Volk sind sie gewesen; er hat alles zu sich gebracht, meinen sie; ist ihnen leid, daß Jakob soll einen eigenen Heller haben. Das müßte er hören, und wo er nicht GOttes Wort gehabt hätte, würde er gedacht haben, sie würden ihm seine Weiber und alle sein Gut nehmen. Doch ficht er das Gut nicht an, sondern spricht hernach [V. 31.] zu Laban: „Ich fürchtete, du würdest deine Töchter von mir reißen.“

5. Weiter sagen sie: „Und von unsers Vaters
\quellseite{248}
Gute hat er solchen Reichthum zuwege gebracht.“ Welche seine Gedanken sind das! Jakob hatte Laban reich gemacht, wie er selbst bekennt, daß ihn GOtt um seinetwillen gesegnet habe; noch dürfen sie sagen: Wenn des Vaters Gut nicht wäre, so hätte er keinen Heller. So muß er noch die Schmach dazu leiden, über alles Unrecht und Gewalt, die sie ihm gethan hatten. Was er nun weiter sagt, als er mit den Weibern redet, ist leicht und wohl erklärt.

\begin{schriftzitat}
V. 11—16.

Und der Engel GOttes sprach zu mir im Traume: Jakob. Und ich antwortete: Hier bin ich. Er aber sprach: Hebe auf deine Augen, und siehe, die Böcke springen auf die sprenklichten, fleckten, und bunten Schafe; denn ich habe alles gesehen, was dir Laban thut. Ich bin der GOtt zu Bethel, da du den Stein gesalbet hast, und mir daselbst ein Gelübde gethan. Nun mache dich auf, und ziehe aus diesem Lande, und ziehe wieder in das Land deiner Freundschaft. Da antwortete Rahel und Lea, und sprachen zu ihm: Wir haben doch kein Theil noch Erbe in unseres Vaters Hause, hat er uns doch gehalten, als die Fremden; denn er hat uns verkauft, und unsern Lohn verzehret. Darum hat GOtt unserm Vater entwandt seinen Reichthum zu uns und zu unsern Kindern. Alles nun, was GOtt dir gesagt hat, das thue.
\end{schriftzitat}

6. Das ist alles nach einander darum erzählt, daß wir wissen, daß der gute Vater Jakob recht gethan hat, und GOttes Befehl gehabt; denn „GOtt hat die Güter Labans entwandt (spricht er) und mir gegeben“. Und hier siehst du, daß GOtt seine lieben Heiligen ja wohl versucht, läßt den Jakob in Angst und Noth zwanzig Jahre; aber endlich kommt er und tröstet ihn.

7. Darum sollen wir auch lernen das Wort fest halten mit Geduld, und nicht davon lassen, wenn es auch hundert Jahre währet. Denn, daß er verzeucht, thut er aus Gnaden, uns zum Besten, auf daß der Glaube stark und groß werde, und er desto reichlicher gebe, was er verheißen hat. Denn das ist seine Art, daß er langsam kommt; er kommt aber wohl; als er auch langsam strafet, strafet aber greulich, gibt Frost und Raum genug, sich zu erkennen und zu bessern. Darum müssen die Auserwählten um der Bösen willen harren und geduldig sein, auf daß er beistehe mit seiner Ehre, daß er langmüthig ist, und nicht bald strafet. Also währt der Frommen Geduld, und dagegen der Bösen Schalckheit lange, aber GOtt kommt zu rechter Zeit, und vergilt es auf beiden Seiten reichlich.

8. So spricht er nun: „Ich habe alles gesehen, was dir Laban thut.“ Da macht es GOtt mit einem Worte alles schlecht und recht, was Jakob gethan hat, und verdammt alles, was Laban gethan hat. Laß nun sagen, wer da will, daß Jakob den Laban getäuscht habe. Denn es ist je so viel geredet, als sollte er sagen: Du hast ihm recht gethan, und er dir unrecht. Das hören die zwo Schwestern, Rahel und Lea, und klagen auch: „Wir haben doch kein Theil noch Erbe mehr in des Vaters Hause“ etc. Sollen die Töchter von ihrem Vater so reden? Sie haben aber so viel versucht, daß sie frei durften sagen: „Er hat uns gehalten als die Fremden“; [wir] haben auch an dem Gute kein Theil überall; er reißt es uns alles ab, und gibt es andern. Welches ihnen auch oft wird wehe gethan haben, und sie\footnote{„sie“ fehlt in der Erlanger.} zuweilen auch gemurrt werden haben, und ungeduldig gewesen [sein]. Aber Jakob hat sie vermählt und getröstet, daß sie es ließen gehen.

9. Dazu, sagen sie, hat er uns nicht gehalten als Kinder, sondern hat uns verkauft, das ist, für Mägd gehalten, die man kauft und verkauft, und hält wie das Vieh, daß alles, was sie tragen und arbeiten, nicht ihr, sondern des Herrn ist, und ihnen nichts gibt denn Futter und Kleider. „Darum (sprechen sie) hat GOtt es unserm Vater entwandt“ etc. Das Wort zeigt auch an, daß sie fromme Kinder gewesen sind, die GOtt und sein Werk erkannt haben. So weit geht diese Historie. Folgt nun ein ander Stück, wie es ihnen auf der Reise gegangen ist.

\begin{schriftzitat}
V. 17—35.

Also machte sich Jakob auf, und lud seine Kinder und Weiber auf Kameele, und führte weg alle sein Vieh, und alle seine Habe, die er zu Mesopotamia erworben hatte, daß er käme zu Isaak, seinem Vater, in das Land Canaan (Laban aber war gegangen seine Heerde zu scheren), und Rahel stahl ihres Vaters Bilder. Also stahl Jakob dem Laban zu Syrien das Herz,\footnote{Handloß: Stahl das Herz. Das Herz stehlen ist Hebräisch geredet so viel, als etwas thun hinter eines andern Wissen, bedeutet aber, daß die Gläubigen den rechten Kern GOttes Wortes fassen, daß die Widersacher nimmer gewahr werden.}
\quellseite{249}
damit, daß er ihm nicht ansagte, daß er flöhe. Also floh er, und alles, was sein war, machte sich auf und fuhr über das Wasser, und richtete sich nach dem Berge Gilead. Am dritten Tage ward es Laban angesagt, daß Jakob flöhe. Und er nahm seine Brüder zu sich, und jagte ihm nach sieben Tagereisen, und erwischte ihn auf dem Berge Gilead. Aber GOtt kam zu Laban, dem Syrer, im Traume des Nachts, und sprach zu ihm: Hüte dich, daß du mit Jakob nicht anders redest denn Gutes. Und Laban nahete zu Jakob. Jakob aber hatte seine Hütten aufgeschlagen auf dem Berge. Und Laban mit seinen Brüdern schlug seine Hütten auch auf, auf dem Berge Gilead. Da sprach Laban zu Jakob: Was hast du gethan, daß du mein Herz gestohlen hast, und hast meine Töchter entführt, als die durch das Schwert gefallen wären? Warum hast du das verborgen, daß du flöhest, und hast dich mir entflohen, und hast es mir nicht angesagt, daß ich dich hätte geleitet mit Freuden, mit Singen, mit Pauken und Harfen? Und hast mich nicht lassen meine Kinder und Töchter küssen. Du hast thörlich gethan. Und ich hätte, GOtt Lob, wohl so viel Macht, daß ich euch könnte Uebels thun; aber eures Vaters GOtt hat gestern zu mir gesagt: Hüte dich, daß du mit Jakob nicht anders denn Gutes redest. Und weil du denn ja wolltest ziehen, und begehrest so fast deines Vaters Haus, warum hast du mir meine Götter gestohlen? Jakob antwortete, und sprach zu Laban: Ich fürchtete mich, du würdest deine Töchter von mir reißen. Bei welchem aber du deine Götter findest, der sterbe hier vor unsern Brüdern. Suche das Deine bei mir, und nimm es hin. Er wußte aber nicht, daß sie Rahel gestohlen hatte. Da ging Laban in die Hütte Jakobs und Lea’s, und der beiden Mägde, und fand nichts. Und ging aus der Hütte Lea’s in die Hütte Rahels. Da nahm Rahel die Bilder und legte sie unter die Streu der Kameele, und setzte sich darauf. Laban aber betastete die ganze Hütte, und fand nichts. Da sprach sie zu ihrem Vater: Mein Herr, zürne nicht, denn ich kann nicht aufstehen gegen dir, denn es geht mir nach der Frauen Weise. Also suchte er, und fand die Bilder nicht,\footnote{Die folgenden Worte, welche in der Bibel und in der Jenaer stehen, fehlen in der Erlanger.} wie fast er suchte.
\end{schriftzitat}


10. Das Herz stehlen, ist eine hebräische Weise zu reden, wenn man hinter dem Rücken ein Ding ausrichtet, daß ein anderer nichts darum weiß. Da hebt sich aber eine Schuld an, daß sich Jakob so heimlich mit Weib und Kind davon stiehlt, und dazu Rahel die Büberei thut, und stiehlt dem Vater die silbernen Götzen. Fragt man abermal, ob es auch recht gehandelt sei? Der Vater hat ihr nichts gegeben, wie er schuldig war, so hat sie gedacht: Ich bin ja auch ein Kind im Hause; [sie] müßte ja etwas davon bringen, hat zum andern nicht kommen können. Ob es aber wohl gethan sei, weiß ich nicht. Man sagt, sie habe es darum gethan, daß sie der Abgötterei weniger machte in des Vaters Hause; doch weiß ich es nicht. Es mag wohl ein weibischer Fehler gewesen sein; doch liegt nicht Macht daran; ob sie gleich mehr gestohlen hätte, müßte [es] doch nicht schaden. Hat sie unrecht gethan, wollen wir sie auch nicht entschuldigen.

11. Das ist aber zu sehen, daß ihnen Laban nacheilte; in welchem aber ein Stück Jakobs\footnote{Erlanger: seines.} Glaubens angezeigt wird. Denn es muß ein kühner Muth und trotziger Glaube gewesen sein, daß er mit so vielen, Weib, Kind, Gesinde und allem Vieh, aufbricht, durch das Land so weit zu ziehen, sonderlich weil er es hinter des Laban Rücken thut. [Er] hat müssen denken: Er wird sich aufmachen, und uns nachziehen, daß er alles, was hier ist, nehme und fange, und auf das greulichste mit uns fahre. Aber es enthält und schüzt ihn nichts denn GOttes Wort, wiewohl sich alles anders läßt ansehen. Das war ein kühner, starker Glaube, der doch hernach schwach wird werden. So befiehlt er nun die Sache GOtt so stark, daß er Weib, Kind und Gesinde führt über das fremde Land, weil er doch fürwahr wußte, daß ihm Laban nacheilen würde. Weil er nun also auf GOtt trauet, läßt er ihn nicht, eben da Laban auf ihn erhitzt, und im Sinne hat, er wolle ihn angreifen und das Müthlein an ihm fühlen. Denn wo er es nicht also gemeint hätte, hätte ihn GOtt nicht dürfen im Schlafe warnen, und gebieten, daß er ihm Leid thun sollte; so kommt er selbst, und schafft so viel, daß ihn Laban muß mit Frieden lassen.

12. Also siehe, wie GOtt seine Heiligen wunderlich läßt allein stehen; aber wenn es zum
\quellseite{250}
Treffen kommt, da die Noth ist, ist er so bald da. Zuvor hat er Jakob zwanzig Jahre gelassen, und gethan, als fragte er nichts nach ihm; jetzt kommt er unversehens und hilft. Wenn wir meinen, er habe unser vergessen, so kommt er zu rechter Zeit so eben, daß wir es nicht besser könnten wünschen. Darum sprach er zu Jakob: „Ich habe alles gesehen, was dir Laban gethan hat.“ Wie, hast du es gesehen, und doch geschwiegen? Ja, jetzt komme ich, und will dich schützen. Also auch Christus im Matthäo [Cap. 10, 30.]: „Alle eure Haare auf dem Haupte sind gezählet.“ Trotz dem, der uns ein Haar krümme; noch stellt er sich, als wollte er uns gar lassen hinweg reißen. Das ist nun die Kunst, daß man GOttes Werk kennen lerne.

\begin{schriftzitat}
V. 36—55. Und Jakob ward zornig und zankte sich mit Laban, antwortete, und sprach zu ihm: Was habe ich mißgehandelt oder gesündiget, daß du so auf mich erzürnt bist? Du hast alle meinen Hausrath betastet, was hast du deines Hausrathes gefunden? Lege es dar vor meinen und deinen Brüdern, daß sie zwischen uns beiden richten. Zwanzig Jahre bin ich bei dir gewesen, deine Schafe und Ziegen sind nicht unfruchtbar gewesen. Die Widder deiner Heerde habe ich nie gegessen. Was die Thiere zerrissen, brachte ich dir nicht, ich mußte es bezahlen, du fordertest es von meiner Hand, es wäre mir des Tages oder des Nachts gestohlen. Des Tages verschmachtete ich vor Hitze, und des Nachts vor Frost, und mein Schlaf wich von meinen Augen. Also habe ich zwanzig Jahre in deinem Hause gedienet, vierzehn um deine Töchter, und sechs um deine Heerde, und hast mir meinen Lohn zehnmal verändert. Wo nicht der GOtt meines Vaters, der GOtt Abrahams, und die Furcht\footnote{Randglosse: Furcht. Jakob nennt hier den GOtt Isaaks: „Furcht“, darum, daß Isaak gottesfürchtig war und GOttes Diener.} Isaak auf meiner Seite gewesen wäre, du hättest mich leer lassen ziehen. Aber GOtt hat mein Elend und Arbeit angesehen, und hat dich gestern gestraft. Laban antwortete, und sprach zu Jakob: Die Töchter sind meine Töchter, und die Kinder sind meine Kinder, und die Heerden sind meine Heerden, und alles, was du siehst, ist mein. Was kann ich meinen Töchtern heute, oder ihren Kindern thun, die sie geboren haben? So komm nun, und laß uns einen Bund machen, ich und du, der ein Zeugniß sei zwischen mir und dir. Da nahm Jakob einen Stein, und richtete ihn auf zu einem Mal, und sprach zu seinen Brüdern: Leget Steine auf. Und sie nahmen die Steine, und machten einen Haufen, und aßen auf demselben Haufen. Und Laban hieß ihn Jegar Sahadutha, Jakob aber hieß ihn Gilead.\footnote{Randglosse: Gilead heißt ein Zeugenhaufen, und bedeutet die Schrift, da viel Zeugnisse von GOtt häufig innen sind.} Da sprach Laban: Der Haufe sei heute Zeuge zwischen mir und dir (daher heißt man ihn Gilead), und sei eine Warte. Denn er sprach: Der HErr sehe darein zwischen mir und dir, wenn wir von einander kommen, wo du meine Töchter beleidigest, oder andere Weiber dazu nimmst. Es ist hier kein Mensch mit uns; siehe aber, GOtt ist der Zeuge zwischen mir und dir. Und Laban sprach weiter zu Jakob: Siehe, das ist der Haufe, und das ist das Mal, das ich aufgerichtet habe zwischen mir und dir; derselbe Haufe sei Zeuge, und das Mal sei auch Zeuge, wo ich herüber fahre zu dir, oder du herüber fährest zu mir über diesen Haufen und Mal, zu beschädigen. Der GOtt Abraham, und der GOtt Nahor, und der GOtt ihrer Väter, sei Richter zwischen uns. Und Jakob schwur ihm bei der Furcht seines Vaters Isaak. Und Jakob opferte ein Opfer auf dem Berge, und lud seine Brüder, das Brod zu essen. Und da sie gegessen hatten, blieben sie auf dem Berge über Nacht. Des Morgens aber stund Laban frühe auf, küssete seine Kinder und Töchter, und segnete sie, und zog hin, und kam wieder an seinen Ort.

Cap. 32, 1. 2. Jakob aber zog seinen Weg. Und es begegneten ihm die Engel GOttes, und da er sie sahe, sprach er: Es sind GOttes Heere; und hieß dieselbige Stätte Mahanaim.
\end{schriftzitat}

13. Da erzählt Jakob sein Leben her, das er geführt hatte nun zwanzig Jahre, darinne wir abermal sehen, was der lieben Heiligen Leben sei: nichts denn Jammer und Herzeleid, wie er sagt, daß er Tag und Nacht keine Ruhe gehabt habe, vor Hitze und Frost verschmachtete sei. Das ist seine Legende; da ist noch nichts von großem, heiligem Wesen, da wir von sagen. [Es] sind eitel Werke eines armen, betrübten Schäfers, der sich hat müssen drücken und brüden, und viel leiden von seinem Herrn und Gesinde, und sonst allerlei zufällig Ding. Denn wiewohl er gewisse
\quellseite{251}
Zusagung hatte, daß ihn GOtt wollte schützen, und bei ihm sein, dennoch läßt er ihn gehen, wie einen andern Menschen, daß ihm auch die Wölfe die Schafe fressen, und allerlei Noth leidet.\footnote{Wahl und die Erlanger: leiden.} Wiewohl es alles scheint, als sei es Menschending; wie denn GOtt seine Heiligen führt, daß man es nicht sieht. Das ist seine Kurzweil gewesen, daß er nicht viel fröhliche Tage seine Lebetage gehabt hat.

14. So erzählt er nun sein Unglück, und kommt auf die Verheißung seines GOttes, und nennt ihn einen GOtt Abrahams, und die Furcht Isaak; das ist auf hebräischer Sprache Weise geredet, welche Gottesfurcht heißt eigentlich Gottesdienst. Wir haben es zu scharf und spitzig gedeutet, wenn die Schrift von Gottesfurcht redet, daß wir allezeit haben großen Unterschied gemacht de timore filiali et servili; aber auf das einfältigste ist es nichts, denn GOtt mit dem Herzen innwendig, und mit auswendigem Wesen dienen, welches darinne stehet, daß man ihn in Ehren halte, und sich vor ihm scheue, thue und lasse nichts, ohne was man weiß, daß [es] ihm wohlgefalle. Darum muß er hier den Namen haben, daß er Isaak Furcht heißt, das ist, den Isaak vor Augen gehalten, und sich vor ihm gescheuet, hat gethan und gelassen, was ihm gefallen hat; wiewohl er es nicht allezeit so eben getroffen hat. Denn man muß sie auch Menschen lassen bleiben, doch also, daß sie mehrern Theils von reinem Glauben GOtt dienen haben. Also muß man die Furcht in der Schrift verstehen, daß es nicht deute auf Furcht oder Schrecken, sondern daß es sei das ganze Leben und Wesen, das da geht in Ehren und Scheu vor GOtt; denn es wird niemand GOtt dienen, denn wer sich vor ihm fürchtet. Das ist das Vornehmliche in diesem Capitel zu sehen, darnach ist weiter nichts Sonderliches, denn wie sie einen Bund machen, und einander die Letzte\footnote{Letze = Letze, Valete, Abschied.} geben.

15. Nun sollen wir auch die geistliche oder heimliche Deutung der Historie suchen, können wir sie anders finden. Also habe ich vorgesagt: Wer da die Schrift geistlich auslegen will, oder in einem verborgenen Sinne, soll vor allen Dingen aufsehen, daß er es also treffe, daß es sich reime mit dem Glauben, oder, wie St. Paulus [Röm. 12, 7.] lehrt, daß [es] dem Glauben ähnlich sei; wo anders, so taugt es nichts. Was heißt denn, dem Glauben ähnlich sein? Das heißt es, wenn man die Leute nicht vom Glauben führt, und nichts Anderes lehrt, denn daß der Glaube bleibe. Denn es geht mit Glauben, gleich wie Paulus [1 Cor. 3, 10—13.] sagt: „Ich habe den Grund gelegt, als ein weiser Baumeister. Ein jeglicher aber sehe zu, wie er darauf baue. Es kann zwar niemand einen andern Grund legen, außer dem, der gelegt ist, welcher ist Jesus Christus. So aber jemand darauf bauet Gold, Silber, Edelsteine, Holz, Heu, Stoppeln, so wird eines jeglichen Werk offenbar werden.“

16. Das ist alles vom Predigtamte gesagt, daß, wer in der Schrift fahren will, und wohl auslegen, der fahre ja also, daß er nichts Anderes lehre, denn daß da eben\footnote{Wittenberger und Jenaer: ähnlich.} sei der Lehre vom Glauben, welche allein gegründet ist und steht auf Christum. Wenn nun jemand zufährt, und predigt einen Menschentraum, als irgend von St. Barbara, wie sie von GOtt erworben habe, wer ihr fastet und feiert, der soll ohne das Sacrament nicht sterben, [das] ist auch gepredigt und gebauet auf den Glauben, als auf Christum, denn man schreiet sie aus, als eine Märtyrin Christi, und will die Christen damit unterrichten, wie man leben soll; aber da siehe zu, wie sich’s auf den Glauben reimt.

17. Wie reimt es sich nun? Also, daß mein Fasten solle mein Trost sein; so will der Glaube keinen Trost wissen, denn den einigen Christum. Nun mag solche Lehre auch von heiligen Leuten gepredigt werden, denn St. Paulus leugnet nicht, daß sie den Grund lassen bleiben; ja, er sagt eben von denen, die den Grund bauen, ob sie gleich Heu oder Stoppeln darauf bauen; wie St. Bernhard, Augustinus, Hieronymus und viele andere, die noch auf dem Narrenwerke standen, daß Möncherei etwas wäre, und damit umgingen (wiewohl sie von Christo sonst recht predigten). Das war Stroh und Stoppeln auf den Glauben gebauet. Aber GOtt konnte es sein mit Feuer hinweg brennen; wie St. Bernhardo geschah. Da das Feuer über ihn kam, das ist, die Versuchung und der Tod, sprach er: Ach, ich habe schändlich gelebt, und mein Leben sündlich zugebracht; ich weiß aber einen Trost, daß GOtt ein zerbrochen und zerschlagen Herz
\quellseite{252}
nicht verachten will [Ps. 51, 19.]. Item: Du hast das Himmelreich mit zweierlei Recht gewonnen: einmal gebühret dir’s Erbes halben, weil du GOttes Sohn bist; zum andernmal hast du es auch verdient durch dein Leiden und Tod, da du es nicht bedürftest, und mir geschenkt. Dieser Mann hatte es lange getrieben mit strengem und hartem Leben und Werken; aber jetzt, da das Stündlein kam, verdorrt es alles, und fällt frei dahin, als das verloren ist; er aber blieb auf dem Grunde, und ward selig, doch durch das Feuer; dadurch muß es [zu]vor gehen, das ist, er muß vor dem Tode also erschrecken, daß jenes alles zu Boden fällt und zu Pulver wird, er aber allein auf dem Grunde bleibe.

18. Also auch St. Augustinus führte ein fein Leben, Regel und Gesetz, davon sich noch die Mönche rühmen; aber am Ende seines Lebens fühlte er, daß nichts half, und mußte die sieben Bußpsalmen singen; da vergaß er seines Lebens, als hätte er nie nichts Gutes gethan. Das sehen wir nicht an, wie der Heiligen gute Werke in der Anfechtung, gleichwie\footnote{Die Worte: „Der Heiligen ... gleichwie“ und gleich folgend: „durch das Feuer“ fehlen in der Erlanger.} das Stroh und Stoppeln durch das Feuer verzehrt ist, daß sie zuletzt zum Kreuze gekrochen sind, und auf dem bloßen Grunde haben müssen bleiben. Was sollte uns denn St. Barbara oder andere Heilige helfen in Todesnöthen? Was wäre mir geholfen, wenn es gleich wahr wäre, daß diese Jungfrau erlangt hätte, daß ich nicht ohne das Sacrament stürbe? Gerade so viel, als der Monstranz, darinne es behalten wird.

19. Das habe ich darum gesagt, daß man sehe, wie alle Predigt dem Glauben solle ähnlich sein, wie es sich reime; daß es immer der Lehre Art sei, daß sie den Glauben treibe. Also haben wir von Anfang dieses Buches nichts gefunden, denn daß sich alles auf den Glauben und Wort richtet; ob gleich immer ander Ding und neue Historie vorgehalten wird, noch artet es sich immer darnach, also, daß der Glaube durch und durch das Hauptstück bleibt. Es sind beide, Historien und Figuren, nichts denn Gold und Silber, auf den einigen Grund gebauet, daß immer einerlei gelehrt wird, obwohl Person, Geschicht,\footnote{Erlanger: „obs wohl persönlich geschicht“.} Zeit und Weise anders und anders ist. Also müssen wir nun des Jakob Historie auch handeln, und erstlich die zwo Hauptpersonen ansehen, darauf es sich alles bezieht, nämlich Jakob und Laban.

20. Jakob (wie gebührt ist) heißt ein Untertreter darum, daß Christus (welches Figur und Bild Jakob ist) solch Regiment führt, daß er mit Füßen tritt alles, was nicht GOttes Wort ist, aller Welt Weisheit, Klugheit, Frömmigkeit, Gewalt und Kraft.

21. Dagegen ist Junker Laban, der heißt auf deutsch blankt oder weiß, gleichend, wie das Licht herblickt, oder scheint in einem Becken. So bringt der Name selbst mit, was er sei und deute, nämlich ein schön gleisender Heuchler, wie der Harnisch glänzt, wenn die Sonne darauf scheint.

22. Darum deutet er das widersinnische Reich, als gewesen sind die Juden, welche das Gesetz Mosis hatten, und Hohepriester, die das Volk sollten lehren und regieren nach dem, das Moses geschrieben hat; welches, weil es nicht mehr konnte geben, denn allein den äußerlichen Sünden steuern und wehren, und sie darauf fielen mit Werken, mußten eitel Gleisner daraus werden. Wie auch noch, wenn noch so viel geistliche Leute wären, die sich zwingen, und führen ein äußerlich schön, ehrbar Leben, das fein helle daher glänzt, und vor den Augen scheint, sind sie doch inwendig im Herzen verzweifelte Buben; denn das Herz kann nimmermehr von Werken und Geboten fromm werden. So ist Laban ein Bild der gleisenden\footnote{Die Worte: „gleisenden“ und (gleich folgend) „des Gesetzes“ fehlen in der Erlanger.} Leute, hat den Namen und Figur des Volkes des Gesetzes, sonderlich des jüdischen Volkes.

23. Solches sieht man auch in des Mannes Gebärden, da findest du zwei Stücke. Zum ersten, daß er Güter hat, die ihm seine Tochter stiehlt. Das andere, daß er seinen Eidam und seine Töchter so schändlich handelt. Denn alle gleisenden und heuchlerischen Heiligen müssen die Art haben, und also thun. Zum ersten, daß sie dem rechten GOtt wollen auf das fleißigste dienen, haben auch allein den Schein, daß sie GOttes Diener sind. Denn der gute Jakob hat keine Götzen, sein Gottesdienst steht in keiner äußerlichen Geberde oder Weise; [er] ist nichts denn ein Hirte, der der Schafe hütet. Dieser aber hat eine eigene Capelle und sonderliche Götzen, und hat doch den rechten GOtt angebetet, wie er im Texte selbst sagt: „Der GOtt Abraham, und der GOtt Nahor, und der GOtt
\quellseite{253}
ihrer Väter sei Richter zwischen uns.“ Also soll man alle Abgötterei deuten in dem jüdischen Volke. Denn so toll sind die Leute nicht gewesen, daß sie gemeint hätten, daß die Bilder, Holz und Steine, der rechte GOtt wären; die Juden sonderlich, die doch viel Abgötterei anrichteten; aber das haben sie gemeint, dem rechten, natürlichen GOtt damit zu dienen. Das ist aber der Fehler daran, daß sie keinen Befehl dazu hatten, sondern aus eigenem, falschem Wahn thaten, unterstehen sich, GOtt einen solchen Dienst anzurichten, das soll ihm gefallen; und weil es ihnen wohlgefällt, und alle für köstlich achten, stehen sie darauf, und denken, sie haben etwas verdient, und seien fast wohl daran.

24. Er aber sagt also: Ich will keinen Gottesdienst, den ich nicht befohlen habe; ich will dein Werk nicht haben, und von dir ungekauft sein, noch dir verkaufen. Sollst du meine Gnade, Geist und Güter haben, so mußt du es lauter umsonst haben, und sonst nichts haben, darauf du dich verlässest. Noch sind sie so thöricht, und fahren zu, taufen nicht allein ihnen, sondern auch für andere Leute, wie unsere Geistlichen gethan haben; welches nicht GOtt gedient, sondern auf das schändlichste gelästert und geschändet ist. Das ist die Labanstotte, die GOtt allein wollen dienen, haben aber nur den Schein davon, daß alles muß gehen unter dem Namen des rechten GOttes.

25. Derhalben mußt du solcher Weise der Schrift gewohnt, die alle Abgötterei dahin deutet, daß die Bilder nichts sind, denn eine äußerliche Weise, GOtt zu dienen. Und wiewohl wir jetzt solche Bilder und Götzen nicht haben, ist es doch demselbigen gleich. Denn ich bin ebensowohl abgöttisch, wenn ich in ein Kloster gehe, und schreibe mir solche Regeln vor, binde mich an eine eigene Weise, darnach zu leben, und GOtt [zu] dienen, davon GOtt nicht geboten, noch je geredet hat. Da diene ich nicht GOtt, sondern der Regel, die ich erdichtet habe, und gebe es dennoch GOtt zu. Solcher Abgötterei ist die Welt voll, wohl, daß ihm ein jeglicher GOtt abmalt nach seinem Wahne und Gutdünken, und will ihn ehren mit selbsterdachtem Dienst.

26. Also bleibt Laban ein Abgöttischer, und macht ihm eigene Götzen. [Es] ist auch wohl zu denken, er werde seine Töchter nicht wohl gezogen haben. Wäre Jakob nicht gekommen, so wären sie wohl in der Abgötterei geblieben; er hat sie aber heraus gerissen, und GOtt gelehrt recht erkennen. Denn, daß Rahel verstanden habe, daß es Narrenwerk wäre, beweist wohl, daß sie die Götzen unter die Streu legt, und setzt sich darauf. [Es] ist auch ein Großes, daß sie den gleisenden Gottesdienst also kann verachten. Denn wir sehen, wie sauer es uns noch wird, des Papstes Geistlichkeit mit Füßen zu treten. Das ist von dem ersten Stück geredet, daß der Mann übt gegen GOtt im gleisenden Wesen, um des erdichteten Gottesdienstes willen. Also handelt er erstlich wider den Glauben.

27. Zum andern können es solche Heilige nicht lassen, sie müssen auch wider die Liebe handeln. Denn wie sie GOtt lästern durch falschen Gottesdienst, so beleidigen sie auch den Nächsten. Also hast du zwei Creuz gegen einander: Jakob geht in reinem Glauben, christlicher Liebe, und trägt das Kreuz; Laban geht in seinem gleisenden Wesen, ohne Glauben und Liebe, und will kein Kreuz tragen, also, daß er auch seinen eigenen Kindern abbricht, und das Ihre vorenthält.\footnote{In den alten Ausgaben: „fürhält“.} Damit hat die Schrift eben abgemalt, was sie für Leute sind, die falschen Heiligen. Er hat der Liebe so gar vergessen, daß er nicht allein den Jakob zehnmal um seinen Lohn täuscht, sondern auch den Töchtern nimmt, was sie haben, und läßt sie mit ihrem Schweiß und Arbeit ihm dienen. Das ist es, daß die falschen\footnote{Erlanger: „diese“ statt: die falschen.} Heiligen nichts thun, denn ihren Wanst und Genuß suchen, und drücken die rechten Heiligen unter.

28. Deß lasse uns ein Exempel sehen. Gib mir ein Kloster, ja, einen Mönch, da man christlich und brüderlich gegen einander lebt, oder einer den andern treulich meint. Man sieht vor Augen, wie man bisher dem Teufel so viel gegeben hat; ja, daß wir nur dem tollen Volke konnten viel geben, haben wir alles zu uns gerissen und gescharrt, dem Nächsten nicht mit einem Heller geholfen; darnach so viel Testamente gemacht, wenn wir sterben sollten, nur die Bäuche zu füllen. Also haben die Pfaffen auch gethan, die haben müssen voll sein, und genug haben, daß ihnen je nichts fehlte, armer Leute Schweiß und Blut gefressen, und niemand umsonst gedient.
\quellseite{254}

29. Und sonderlich mercke das, daß Laban nicht allein alles zu sich reißt, den Töchtern nichts gibt, sondern auch sie darüber zwingt zum Dienste, und Mägde aus ihnen macht, als die er um Geld kaufte. Das ist die Tyrannei der Menschengesetze, wie die Juden auch thaten, wenn sie die Heiden unter sich brachten, und drangen sie dahin, daß sie ihre Gesetze mußten halten, so ihnen zuträglich waren. Als das war, davon Christus [Matth. 23, 16.] sagt: „Wehe euch, verblendete Leiter, die ihr suget: Wer da schwört bei dem Tempel, das ist nichts; wer aber schwört bei dem Golde am Tempel, der ist schuldig.“ Item, V. 18.: „Wer da schwört bei dem Altare, das ist nichts; wer aber schwört bei dem Opfer, das droben ist, der ist schuldig.“ Das Opfer, sagen sie, ist ein großer, köstlicher Gottesdienst, viel besser denn der Altar; denn um des Opfers willen ist er gemacht. Also ist auch das Gold viel köstlicher, denn der Tempel, denn es ziert den Tempel. Das thaten sie um ihres Nutzens willen, und brachten die Leute dahin, daß sie dachten: Nun will ich viel lieber Gold und Opfer geben, denn zu dem Altare oder Tempel.

30. Also auch zu unsern Zeiten: Wie haben wir müssen leiden und uns lassen martern, daß man nichts gepredigt hat, denn was ihnen gedient hat, hier und da zu geben. Das haben sie also gefasset und geordnet mit Gesetzen, daß, wer anders thäte, der wäre in GOttes Ungnade. Also sind wir ihre Gefangenen gewesen, und haben müssen thun, wozu sie uns zwangen. Doch hat der Heilige Geist etliche erhalten, die andern haben sich müssen leiden. Deß hat niemand denn der Laban genossen; [er] ist reich worden mit unserm Schaden.

31. Aber als [es] Rahel und Lea gewahr wurden, sprachen sie: Mache dich nur davon, und nimm uns mit, wir sind doch hier nicht mehr denn gefangene Mägde. Solches predigen wir jetzt auch, daß wir der Lehre, damit sie uns gezwungen haben und nur ihren Nutzen geschafft, entlaufen und heraus reißen. Es haben auch die Propheten hin und wieder viel geklagt, wie die falschen Propheten das Volk nur so ausgezogen, und an Leib, Gut und Seele verderbt haben, und sie davon sind fette Wänste worden.

32. Siehe, das hat GOtt in diesem Laban angezeigt, daß es sich reime, wie ich [§ 27] gesagt habe, auf den Glauben, wie beide, Glaube und Unglaube, handeln, sonderlich wie der Unglaube immer will den Wanst nur voll haben, niemand dienen, und nichts leiden; der Glaube aber handelt nach GOttes Worte, dient jedermann durch die Liebe, und geht im Leiden und Kreuz.

33. Weiter. Die zwei Töchter und Mägde müssen wir auch deuten auf das jüdische Volk; also, daß sie, wie wir oben [Cap. 17, § 11. Cap. 25, § 26] gehört haben von Sara und Rebekka, eine Figur, dazu auch Glieder der christlichen Kirche sind. Denn Jakob muß eine Figur tragen des Königes, der ohne Gesetze regiert, darum er auch ein armer König auf Erden in seinem Reich ist, wiewohl er alle Dinge in der Hand hat, führt aber sein Volk in Armuth, Schwachheit, Tod und allen Jammer. Gewalt und Reichthum hindert sein Reich; darum spart er es in jenes Leben, übt und treibt uns hier mit Kummer und Unglück, wie er Jakob zwanzig Jahre dienen läßt, und viel Ungemach leiden.

34. So sind nun die vier Weiber, das Volk, aus den Juden genommen. Denn ich wollte gerne aus den Vieren nur Eine machen; denn Christus soll nur Eine Braut haben, wie wir im Glauben bekennen: Ich glaube eine gemeine christliche Kirche. Und Paullus [Eph. 4, 5. 6.] sagt, daß alles sei Ein Glaube, Eine Taufe, Ein HErr, und Ein GOtt, und der einige Glaube ist die Jungfrauschaft und Reinigkeit, und der Schatz miteinander dieser Braut, hat auch keine andere Speise, Samen, Milch \&c., denn dasjenige einige Wort GOttes, damit sie alle Christen nährt. Darum müssen wir die Viere zusammen in Eines schließen.

35. Zum ersten ist Rahel die Liebste und Schönste, [sie] war ihm verheißen, und [er] hat auch um ihretwillen allein gedient: Lea aber ist nicht die Liebste, doch ist sie die Erste. Das wollen wir so deuten, auf das Beste wir können; weiß nicht, ob wir es so genau treffen; [es] ist genug, wenn wir den Grund halten. Rahel heißt auf deutsch ein Schaf; Lea heißt müde, bekümmert, jämmerlich, die war bleich, mager, dürre und siech von Angesicht, Rahel aber fett und völlig. Das soll sein die Christenheit nach zweierlei Wesen, innerlich und äußerlich. Daher wir genannt werden alte und neue Menschen; doch nicht, daß der Mensch zwei Dinge sei; so viel wir den Glauben haben, sind wir
\quellseite{255}
nen, so viel mir nicht glauben, sind wir alt. Darum kann ich den alten Menschen nicht deuten, als sei er allein Fleisch und Blut, denn der, wie auch der neue Mensch, der ganze Mensch ist. Äußerlich bin ich wohl anders, wenn ich Lea bin, innerlich bin ich anders, wenn ich Rahel bin. Doch ist beides Ein Mensch, denn ich fühle es beides, Glauben und Unglauben, im Grunde des Herzens, und bricht heraus im ganzen Menschen. Darum ist es nur so viel geredet: So viel du glaubst, bist du ein neuer Mensch durch und durch; und wiederum.

36. So sei nun Rahel das innerliche Wesen, durch den Glauben, die hatte die Verheißung GOttes, ist die gelobte Braut, darinne Christus regiert; wie Paulus Eph. 3, 17. sagt, daß er wohne in unsern Herzen durch den Glauben. Diese Braut hat er lieb, und stellt sich freundlich gegen ihr, und dient um ihretwillen, das ist, setzt sein Leben an sie; wie Paulus sagt Eph. 5, 25. 26.: „Ihr Männer, liebet eure Weiber, gleichwie Christus geliebet hat die Gemeine, und sich selbst für sie gegeben, auf daß er sie heiligte“ 2c. Aber die Lea hat er nicht so lieb; das ist, die Christenheit nach dem äußerlichen Wesen, wie man sie ansieht, da ist sie müde, bekümmert und jämmerlich, liegt in Armut und Elend. Da geht die Klage, so sie führt: „Ich bin geheilet“, und muß das Kreuz tragen, und Nichtenbrödel sein. Denn wiewohl er uns durch und durch lieb hat, scheint es doch nicht. Inwendig haben wir wohl Frieden, Lust und Freude, da liegen wir ihm in Armen, und [er] thut sich freundlich zu uns, aber äußerlich stellt er sich, als kenne er uns nicht, läßt uns viel Unglücks und Herzeleid anlegen.

37. Aber nichtdestoweniger muß Lea die ersten Kinder mit ihm haben und Früchte bringen; das ist, wiewohl es scheint, er stelle sich nicht freundlich gegen ihr, doch läßt er sie nicht, sondern tröstet sie; daß wir nicht denken, er habe unser vergessen; gibt, daß man desto mehr Werke und gute Früchte überkommt. Denn das leidende Leben, das da so versucht und durchtrieben wird, bringt viel mehr guter Werke, denn das innerliche, kann viel schaffen, rathen und andern nütze sein; welches ein Unerfahrner nicht thun kann. Darum muß Lea die Allerreichste und Fruchtbarste sein, um ihres Elendes willen; wie der Text spricht: „GOtt sahe an, daß sie gehasset war, und machte sie fruchtbar.“

38. Rahel bringt auch zuweilen Frucht, aber nicht so viel, und wird ihr sauer und schwer, [es] geht langsam zu, [sie] muß zuletzt das Leben darüber lassen. [Es] ist alles darum geredet: weil GOtt sieht, daß wir so blöde sind, daß unfreundliche und mühselige Leben zu führen, tröstet er uns also: Du sollst viel Gutes thun und fruchtbar sein, du mußt Lea sein; darum leide dich, und sei getrost; wenn ich dich immer tröstete und mit dir spielte, so würdest du faul werden und unfruchtbar bleiben. Darum versucht er sie, und läßt es ihr sauer werden, daß sie nicht eitel Lust und Freude an ihm suche. Also läßt er viel Werke thun von der schwachen, leidenden Lea, in äußerlichem Wesen des Kreuzes; von der andern aber, als Rahel, in dem inwendigen Wesen, wenig.\footnote{Erlanger: „von der Schwachen, leidenden Wesen, der Andern wenig“.} Denn das innerliche Wesen muß sich nicht viel herausbrechen, sondern beides gehen, wie es geordnet ist, inwendig im Gewissen gegen GOtt, auswendig vor der Welt, dadurch andere auch herzu kommen, und die Christenheit gemehrt wird. Sie meint wohl, sie wolle die Erste bleiben, und die besten Tage mit ihm haben; es muß ihr aber fehlen.

39. Das geht alles zu durch die Erfahrung, und reimt sich recht mit dem Glauben. Denn wir sind alle also gesinnt, daß wir gerne wollten Frieden und Freude im Gewissen haben, und nach der Rahel greifen; aber GOtt wendet es um, daß wir zum ersten müssen Lea bleiben, wie Jakob zum ersten bei der Lea schläft, so er doch nur um Rahel gedient hatte. [Es] geht aber heimlich zu, daß er nichts darum weiß; denn GOtt braucht auch böse Buben dazu; wie ihm der Laban mit List und Täuscherei die Lea zulegt. Wenn wir das nun hin haben, so gibt er uns Rahel auch; das ist, wenn wir so versucht sind, kommt darnach Trost und Ruhe des Gewissens, welche denn stets und fest bleibt. So müssen wir beide nacheinander haben; sonst, wo Jakob die Rahel allein hätte gekriegt, hätte er Lea nicht angesehen. Also wird aus den zwei Weibern Eine Christenheit, die da recht geht, beide nach innerlichem und äußerlichem Wesen.

40. Was wollen wir denn aus den beiden Mägden machen? Wir haben gehört, wie der Vater jeglicher Tochter eine Magd zugab, welche sie beide ihrem Manne gaben. Das deutet nun
\quellseite{256}
den äußerlichen Menschen, davon Paulus [Röm. 8, 14. Gal. 5, 18.] sagt: Das Fleisch soll unterworfen und gehorsam sein dem Geiste; der Leib muß dienen, und im Zwange gehen. Denn wiewohl das Leben geistlich ist im Glauben, wird es doch nicht ausgerichtet denn durch die Glieder. Die Mägde bringen nun auch Kinder, das ist, ein jeglicher Mensch hat Werke, die beide aus dem Trost, und aus dem Herzeleid kommen (denn Früchte oder Kinder müssen wir immer auf Werke deuten). Lea und ihrer Magd gibt er viel Kinder, das ist, des Leibes Gliedmaßen müssen auch viel Werke thun; Rahel aber und ihre Magd haben ihrer nicht viel, das ist, wenn wir auswendig viel leiden und versucht werden, thun wir mehr Werke, denn wenn wir ruhen, und inwendig im Glauben stehen. Also werden alle vier Weiber Eine Kirche, wie sie gestaltet ist auswendig und inwendig.

41. [Dies] ist nun die Summa: Christus, der König, regiert also, daß er uns läßt betrübt sein, und wiederum tröstet. Davon kommen denn gute Werke, läßt uns viel Unglück und Leid tragen, daß er uns wohl versuche und stark mache. Er ist aber nicht so ein Narr, daß er sich betrügen lasse, und nichts darum wisse, wie Jakob im Texte betrogen wird; sondern es ist nichts mehr, denn daß er sich so stellt, und von uns nicht anders erkannt wird.

42. Das sei die Deutung von Jakob und seinen Weibern. Folgt nun, das beschrieben ist, wie er habe die Schafe zu sich gebraucht durch die Kunst, daß er die Stäbe oder Stecken in die Tränkrinnen legte, daß die Schafe darüber empfingen, und kriegte die bunten oder einfarbigen, wie er sie haben wollte; also, daß Laban die Spätlinge kriegte, Jakob aber die Frühlinge und besten Schafe. Darinne ist bedeutet das Predigtamt, was das ausrichtet in der Welt. Wir haben gehört, wie Laban ein Gleisner ist, der sich auch GOttes Wortes rühmt, gibt vor, als auch wahr ist, daß er GOttes Gebot predigt; aber es fehlt daran, daß er nicht predigt den rechten Glauben und Evangelium, sondern führt die Leute auf die Werke und selbsterwählten Gottesdienst. Also zanken sich nun die zweierlei Prediger mit einander in der Welt von Werken und vom Glauben; da hilft kein Wehren, kein Feuer, keine Gewalt für, man muß es lassen gehen, wie diese und alle anderen Figuren ausweisen. Denn, wenn man es gleich dahin brächte, daß nur einerlei Predigt ginge (wie oft versucht ist), so werden doch mancherlei Secten daraus, wie wir unter dem Papste gesehen haben. Da ist wohl einerlei Predigt gegangen, daß sie uns alle vom Glauben auf Werke geführt haben, doch sind daneben so viel Secten eingerissen, unter welchen jegliche hat wollen die beste sein, und ihr Ding vorgezogen, daß es doch zertrennt ist worden.

43. Darum ist's die beste Weise, man lasse es gehen, und wider einander predigen, falsche und rechte Predigten, wie Paulus [1 Cor. 11, 19.] sagt: „Es müssen Secten oder Spaltungen sein, auf daß bekannt und offenbar werden, die da rechtschaffen sind.“ Also geht es hier auch: Jakob schafft so viel, daß er vor seine Schafe bunte Stäbe legt in die Tränke, davon kriegt er bunte Schafe; aber weil Laban sieht, daß der bunten so viel ist, will er die bunten haben, und wechselt also zehnmal um seinen Lohn. Das ist so viel gesagt: Wenn der Geist predigt, so will das Fleisch immer nachahmen und nachthun, und trifft es doch nicht. Predigt das Evangelium von bunten, so will es das Fleisch auch bunt machen; und wiederum. Also will der Teufel immerdar GOttes Affe sein, und trifft es doch nimmer.

44. Das sehen wir durch und durch in der Welt. Die Apostel haben gepredigt allein den Glauben an Christum, und die Liebe gegen dem Nächsten [zu] üben. Das ist die rechte Predigt, und die Stäbe, davon die Seelen bunt werden, das ist, geziert und geschmückt durch mancherlei Weise und Werke, so in der Christenheit gehen. Denn die christliche Kirche hat durch und durch nur Ein Wort, Glauben, Liebe und Hoffnung, und einerlei Gut; doch bleibt es Buntwerk in äußerlichem Wandel, daß sie mancherlei Amt und Gaben von GOtt haben; einer hat Gnade zu predigen, der andere zu regieren etc. [Röm. 12, 6. f. 1 Cor. 12, 5. f.] Denn wir können nicht alle nach dem äußerlichen Wesen gleich sein, wiewohl wir alle gleich und eins sind im Glauben. Darnach GOtt seine Gaben austheilt, sonst oder so, darnach sind wir bunt, und sprenklicht, und bleiben gleichwohl alle Schafe in der Einigkeit. So heißt nun, Schafe tränken in den Rinnen über den Stäben, solche Sprüche vorlegen, davon man Einen Glauben faßt, aber mancherlei Gaben und Werke.

45. Den will nun der andere Haufe nach, und wenn er sieht, daß so mancherlei Gaben
\quellseite{257}
und Weisen unter den Christen im äußerlichen Wandel sind, will er es auch also machen, richtet mancherlei Stände und Secten an; der wird ein Pfaffe, jener ein Mönch, und so fortan. Aber es geräth nicht also, wird nichts daraus denn Trennung von der Einigkeit des Glaubens; denn ihr Ding ist nicht GOttes Wort und Gabe, sondern Menschentand, aus ihren eigenen Köpfen vorgenommen. Da hebt Jakob nun die bunten Stäbe auf, das ist, wenn man die Leute auf Werke führt, und Secten will machen, so werfen sie die Stäbe weg, und predigen das Widerspiel, daß die Werke ohne Glauben nicht gelten.

46. Also haben uns unsere Heiligen den 45. Psalm, V. 10., ausgelegt: Astitit regina a dextris tuis in vestimento deaurato, circumamicta varietate, die Königin stehet zu deiner Rechten, in köstlichem Golde, bekleidet mit Buntwerk. Das haben sie also gedeutet: Die christliche Kirche hat einen güldenen Rock, das ist, der Glaube, oder Liebe, wie sie wollen; aber der bunte Rock soll sein, daß man in einem Bisthume so viel Messen hält, im andern so viel Lectionen in der Messe liest; diese führen die Chorkappen oben, die andern unten; in dem Stifte oder Orden hält man es sonst, im andern so. Wenn man solche Possen reißt in der Kirche, was sie gelüstet, das soll hier Buntwerk heißen. Das ist eben Labans Buntwerk. Es gilt nicht, daß du es nach deinem Traume anrichtest, sondern es muß von GOtt kommen, wie es Jakob aus dem heiligen Geiste hatte. Das ist nun der rechte güldene Rock, wenn wir alle Einen Glauben haben; daneben der bunte Rock, daß wir in allerlei Glauben mancherlei Gnade haben, daß alle Werke und Gaben in dem Glauben und Liebe gehen. [Es] ist nicht darum zu thun, daß [es] bunte Farbe sei; sondern daß [es] vor GOtt bunt sei, das ist, daß [es] GOtt also wohlgefällt. Das ist das Beste daran.

47. Darum, daß Laban hier also verwechselt, und Jakob zehnmal täuscht, ist, daß es immer die falschen Lehrer wollen meistern und besser machen, bunt oder einfarbig. Da es soll Buntwerk sein, ist es des Teufels Buntwerk, denn sie sollten einen einfältigen Glauben haben. Wenn wir denn solches verwerfen, und vom einfältigen Glauben predigen, wollen sie es auch einfarbig haben, und kein Werk thun. Das gilt abermal nicht. Denn also sagen wir: Der Glaube ist der Art, daß er Buntwerk treibt. So müssen wir immer wechseln, und anders predigen, auf allen Seiten wehren; noch wollen sie nimmer recht hinan, betrügen nur die Welt und sich selbst mit ihren Predigen.

48. Das ist nun, daß Laban behält die Spätlinge, das ist, die letzten Schafe, so im Herbste kommen, Jakob aber die Frühlinge. Die Spätlinge sind die schwächsten Schafe; das sind die Leute, die durch Menschenlehre aufgezogen werden in diesen und jenen Werken, sind arme, elende Seelen. Die aber im Glauben wandeln, sind die rechten, starken Gewissen, da wird Christus gepredigt, daß die Leute zunehmen, und die Christenheit gemehrt wird.

49. Also ist dies alles auf das Predigtamt zu deuten, so, daß die Stäbe nichts anders sind, denn die Sprüche der heiligen Schrift, sonderlich, die da trösten, als des heiligen Evangelii. Denn es ist nicht umsonst dabei gesetzt, daß sie die Schafe ansehen, wenn sie zusammen kommen zu trinken in der Hitze, und darüber empfingen. Denn so haben wir in der Schrift, daß zur Nahrung gehört zweierlei, Essen und Trinken. Das muß beides beieinander sein; wo eines ohne das andere ist, kann man das Leben nicht erhalten. Als, daß wir ein Exempel sehen, wie Maria, der Jungfrau, widerfuhr; die hat gefastet, daß ihr Sohn sollte ein König werden, das war die Speise; aber wo sie eitel geblieben wäre, wäre es ihr schädlich gewesen; darum muß er ihr auch zu trinken geben, welches geschah, als er am Kreuze hing, da gewann sie einen andern Geschmack.

50. Das Trinken ist nichts anders, denn wenn man in großem Herzeleid steckt, und er uns tröstet. Als, wenn er uns läßt sagen, daß wir der Sünde sollen los sein, und der Tod soll uns nicht schaden; das ist das Essen, das schmeckt uns wohl. Aber wir müssen auch zu trinken haben; das geschieht, wenn er herkommt und will mich würgen, da bedarf ich des Trostes, daß ich fest halte am Worte, und mir die Speise nicht mache; das ist der Trank dazu, damit man sich fühlt und erquickt.

51. Siehe, das ist es, das der Text anzeigt, in dem, daß die Schafe in der Hitze mit einander laufen, und darüber trinken. Denn da gibt das Wort erst einen rechten Schmack; wie man sagt: Auf einen guten Bissen gehört ein guter Trunk, da kommt Leib und Seele zusammen
\quellseite{258}
. Darum muß das Kreuz auch da sein, und uns solche Kost schmackhaftig und süß machen. Das heißt auch recht, Christi Fleisch und Blut trinken, wie im Johanne [Cap. 6, 54. 56.] steht. Also verstehe es durch und durch in der Schrift und Figuren, in welcher viel von Essen und Trinken geredet wird. So haben wir die Figur von dem Laban, das ist, von heuchlerischen Heiligen, und den Schafen, und vom ganzen Reiche Christi.

52. Zuletzt ist noch da, wie Laban dem Jakob nacheilt, und ergreift ihn auf dem Berge Gilead, und sucht seine Götzen, die ihm seine Tochter gestohlen hatte; sie aber hatte sich darauf gesetzt, und stellte sich, als wäre sie krank und könnte nicht aufstehen. Das ist, daß die heilige christliche Kirche bekennt, daß sie eine arme Sünderin sei, stiehlt aber nichtsdestoweniger die falsche Lehre, worauf die Heuchler sich verlassen, und versteckt sie also, daß sie das Aergerniß beiseite thue, daß die falschen Lehrer darum kommen. Denn das sieht man: wenn man GOttes Wort predigt, daß die Trügerei untergeht, daß niemand weiß wie; und doch die Christen schwach und krank scheinen, bekennen auch, daß sie arme Sünder sind; aber durch die Schwachheit, die man sieht, und wir fühlen, geht nichtsdestoweniger, weil wir am Worte bleiben hangen, die falsche Lehre mit ihrem Vertrauen unter.

53. Also kann man falscher Lehre nicht baß rathen, denn man verachte sie nur, und halte nichts darauf; ob wir gleich schwach im Glauben sind, dennoch muß jenes untergehen, so man nicht daran hanget. Also ist es gegangen in der ersten Christenheit, zur heiligen Märtyrer Zeit, daß sich die Christen täglich zusehens mehrten, und des Teufels Dienst dahin fiel und zu Boden ging, als [es] denn noch heute auch bei denen geht, die keinen rechten Glauben haben. Denn wie viel sind ihrer, die dem Evangelio nicht glauben, und doch von dem Papstthume nichts halten, als, daß es von aller Welt verachtet wird.

54. Darnach sagt er von dem Berge, darauf sie haben einen Bund gemacht. Derselbige ist hier wohl zu merken, denn die Propheten viel davon geredet haben, sonderlich Jeremias, welcher ihn anziehet im 22. Capitel, V. 6.: Gilead, tu mihi caput Libani. [Es] ist ein kleiner Berg gewesen, als er da geopfert hat, aber hernach ein großes Land da worden, daß alles Gilead geheißen hat von dem Berge, und von der Geschichte, daß die Zwei mit einander eins worden; und gibt ihm den Namen daher, daß er hieß seine Brüder Steine lesen, und einen Haufen machen, der da sollte Zeuge sein des Bundes, und ein Zeichen, daß kein Theil darüber sollte fahren, und ein Schiedsmann, daß keiner dem andern schadete. Denn das Wort Gilead heißt ein Zeugenhaufen, aus zweien Wörtlein gemacht: Gal, das ist, ein Haufe, und Ed, ein Zeuge. Der Zeugenhaufe bedeutet die heilige Schrift.

55. Bisher hat man alle Sachen, die sich entspinnen über dem Glauben zwischen rechter und falscher Lehre, auf ein Concilium geschoben, oder dem Pabste zu Rom, oder den hohen Schulen heimgestellt, die haben sollen Schiedsleute sein. Aber sie sind nicht Gilead, sie haben uns verführt und betrogen, sondern die heilige Schrift, darinne Sprüche über einander liegen mit Haufen, wie dort die Steine; darinne sollen wir das Urtheil holen, wer recht oder unrecht lehrt. Denn wiewohl der Heilige Geist jedermann selbst lehrt im Herzen, daß er weiß, was recht ist, so muß man dennoch die Schrift gebrauchen, damit zu beweisen, daß [es] also sei, wie wir im Herzen glauben. So urtheilt sie denn, und scheidet, was recht und unrecht geglaubt sei. Darum darf man kein Zeugniß weiter suchen, weder von Vätern, noch Concilien, sondern soll allein bei der hellen und klaren Schrift bleiben. Das sind die drei Capitel, wie es Jakob gegangen ist bei seinem Vetter Laban zwanzig Jahre über, mit ihrer Deutung.
\quellseite{259}

\section{Das zweiunddreißigste Capitel}
\begin{schriftzitat}
V. 3—24. Jakob aber schickte Boten vor ihm her zu seinem Bruder Esau, in das Land Seir des Feldes Edom, und befahl ihnen, und sprach: Also saget meinem Herrn Esau: Dein Knecht Jakob läßt dir sagen: Ich bin bis daher außen gewesen, und bin bisher unter den Fremden gewesen, und habe Rinder und Esel, Schafe, Knechte und Mägde, und habe ausgesandt, dir, meinem Herrn, anzusagen, daß ich Gnade vor deinen Augen fände. Die Boten kamen wieder zu Jakob, und sprachen: Wir kamen zu deinem Bruder Esau, und er zeucht dir auch entgegen mit vierhundert Mann. Da fürchtete sich Jakob sehr, und ihm ward bange, und theilte das Volk, das bei ihm war, und die Schafe, und die Rinder, und die Kameele in zwei Heere, und sprach: So Esau kommt auf das eine Heer, und schlägt es, so wird das übrige entrinnen. Weiter sprach Jakob: GOtt meines Vaters Abraham, GOtt meines Vaters Isaak, HErr, der du zu mir gesagt hast: Zeuch wieder in dein Land, und zu deiner Freundschaft, ich will dir wohl thun; ich bin zu geringe aller Barmherzigkeit und aller Treue, die du an deinem Knechte gethan hast. Denn ich hatte nicht mehr, denn diesen Stab, da ich über diesen Jordan ging, und nun bin ich zwei Heere worden. Errette mich von der Hand meines Bruders, von der Hand Esau's, denn ich fürchte mich vor ihm, daß er nicht komme, und schlage mir die Mütter sammt den Kindern. Du hast gesagt: Ich will dir wohl thun, und deinen Samen machen wie den Sand am Meere, den man nicht zählen kann vor der Menge. Und er blieb die Nacht da, und nahm von dem, das er vorhanden hatte, Geschenke seinem Bruder Esau, zweihundert Ziegen, zwanzig Böcke, zweihundert Schafe, zwanzig Widder, und dreißig säugende Kameele mit ihren Füllen, vierzig Kühe und zehn Farren, zwanzig Eselinnen mit zehn Füllen. Und that sie unter die Hand seiner Knechte, je eine Heerde sonderlich, und sprach zu ihnen: Gehet vor mir hin, und lasset Raum zwischen einer Heerde nach der andern. Und gebot dem ersten, und sprach: Wenn dir mein Bruder Esau begegnet, und dich fragt: Wem gehörst du an? und wo willst du hin? und wes ist es, das du vor dir treibst? sollst du sagen: Es gehört deinem Knechte Jakob zu, der sendet Geschenke seinem Herrn Esau, und zeucht hinten hernach. Also gebot er auch dem andern und dem dritten, und allen, die den Heerden nachgingen, und sprach: Wie ich euch gesagt habe, so saget zu Esau, wenn ihr auf ihn stoßet. Und saget ja auch: Siehe, dein Knecht Jakob ist hinter uns. Denn er gedachte: Ich will ihn versöhnen mit dem Geschenke, das vor mir hergeht, darnach will ich ihn sehen; vielleicht wird er mich annehmen. Also ging das Geschenk vor ihm her. Aber er blieb dieselbe Nacht im Lager. Und stand auf in der Nacht, und nahm seine zwei Weiber, und die zwo Mägde, und seine elf Kinder, und zog an den Furt Jabot, nahm sie und führte sie über das Wasser, daß hinüber kam, was er hatte; und blieb diesseits allein.
\end{schriftzitat}

1. Bisher ist Jakob noch ein Pilgrim, wie seine Väter gewesen waren, hat kein eigen Haus noch Wohnung, da er sich setzte, lebt mit allem seinem Gesinde unter dem Himmel, auf GOttes Gnade. Esau aber, sein Bruder, hatte sich längst gesetzt, und ein eigen Land eingenommen, war bereits ein mächtiger Herr mit seinen Kindern worden, daß sich Jakob vor ihm mußte fürchten und demüthigen als ein elender Fremdling, so doch GOttes Verheißung also lautete, daß Jakob sollte Herr sein, und Esau ihm dienen. Dies ist nun auch ein seltsam Capitel, darinnen haben wir abermal ein Wunderwerk GOttes, doch anzusehen, als sei es abereins eitel Kinderwerk, damit Jakob umgeht. Am Ende des vorigen Capitels hat Moses geschrieben, wie er von Laban seinen Weg zog, und unter Wege ihm begegneten die Engel GOttes, und da er sie sahe, gab er der Stätte einen Namen davon zum Gedächtniß, und hieß sie Mahanaim, das ist, Lager oder Heer.

2. Da hatte ihn GOtt also gestärkt und frisch gemacht, daß er sich vor Laban nicht fürchten sollte. War es nicht eine kühne That, und ein fester Glaube, daß er durfte aufbrechen mit Weib, Kind, großem Gesinde und allem Gute in ein fremdes Land, seinem Schwäher zu entrinnen, der auf das höchste wider ihn erbittert war, und durch die Leute zu ziehen, da er unbekannt
\quellseite{260}
war? Ueber das zieht ein Heer der Engel GOttes bei ihm her, nur darum, daß er je desto stärker und fester wäre; also, daß er stand in einem hohen, unüberwindlichen Glauben. Nun aber ist es auf einen Hui umgekehrt, daß er so schwach wird; der sich vor Laban gar nichts fürchtete, ja, gewiß war, daß die Engel GOttes bei ihm standen, der verzweifelte jetzt schier vor seinem Bruder, den er nicht sieht.

3. Das ist, daß ich gesagt habe, wie GOtt so wunderlich in seinen Heiligen regiert, läßt sie so stark werden, daß sie sich weder vor Teufel noch Tode fürchten. Und wiederum, wenn man meint, sie sollen am stärksten sein, macht er sie so schwach, daß sie sich schier vor einem Baumblatte entsetzen; wie Christus [Matth. 14, 29. 30.] mit Petro that. Als er aus dem Schiffe trat, und wollte auf dem Wasser zu ihm gehen, da war er fest und sicher, und fürchtete niemand, aber so bald er einen Wind sahe herkommen, hob er an zu zagen und sinken.

4. Es ist ein köstlich Ding um den Glauben, und große Gewalt. Ja, wenn er in unserer Macht stände, [ihn] zu überkommen, oder, wenn wir ihn gleich haben, auch zu behalten. Es ist ohne Zweifel, der Glaube thut es alles, daß ihm Teufel und Tod weichen muß; aber wo ist die Kunst, daß man dabei bleibe, und ihn behalte? GOtt hat es ihm selbst vorbehalten, den Glauben zu geben und zu nehmen, wann er will, und stark oder schwach zu machen, wie es ihm gefällt. Darum muß es dem Jakob gerade sowohl widerfahren als uns, allein darum, daß man GOtt recht erkennen lerne, was er im Sinne hat. Wir dürfen den Jakob nicht zu hoch heben. Er ist ein großer Heiliger, aber es kann bald geschehen, daß er schwächer im Glauben wird denn ich. Also kann mir Gott auf diese Stunde einen hohen, starken Glauben geben, aber wiederum, ehe man sich umsieht, mich sinken lassen, und den Glauben irgend einem großen Sünder geben.

5. Warum thut er denn solches, daß er seine Heiligen also läßt gehen, und nicht immer in starkem Glauben bleiben? Darum, daß sie nicht wieder stolz werden, oder meinen, sie hätten es von ihnen selbst, und sich selbst zum Gott machen. Darum muß er es so starken und mengen, daß sie wissen, daß er GOtt sei, sich selbst erkennen, und bleiben in der Demuth; die will er haben, nicht allein von uns, sondern von den allerhöchsten Heiligen, auch seiner eigenen Mutter; [sie] müssen sich alle auf das tiefste herunter lassen und sagen: Ich bin nichts, und vermag nichts; also, daß wir Gott alle gleich sind; wiewohl ein Unterschied nach den Gaben ist, welche doch nicht unser, sondern auch alle sein sind. Das ist ein Hauptstück dieses Capitels, wie der Glaube dem heiligen Patriarchen schier gar entfällt, und er ihn doch wieder ergreift, wie wir hören werden.

6. Das Land Seir oder Edom, da Esau gewohnt hat, und nach ihm den Namen gegeben, liegt nicht ferne von der Straße. Denn Edom und Juda stoßen an einander. Weil nun Jakob hart an demselbigen Wege herziehen muß, schickt er Boten vor ihm her, um Geleit zu erkunden, wie sein Bruder gegen ihm gesinnt ist. Denn er fürchtet sich noch vor ihm, weil er ihm die Erstgeburt und den Segen genommen hatte, und hatte deshalb vor zwanzig Jahren von Vater und Mutter müssen ziehen, daß er ihn nicht dermal eins ermürgete. Dazu war Esau nun gewaltig worden, und Herr im Lande, also, daß es sich alles ließ ansehen, als wäre das Wort nichts, das GOtt gesagt hatte: „Der Größere soll dem Kleinern dienen“, und muß vor der Vernunft zum Lügner werden. Dazu muß Jakob Esau einen Herrn heißen, und zu Fuße fallen, und um Gnade bitten, mit Weib und Kind, als einem Landesfürsten; ja, er hält ihn noch dazu wie einen Engel. Wie reimt sich das auf die Verheißung? Oder, ist das gehalten, was zugesagt ist? Es ist alles uns geschrieben. Denn wir sind auch Herren über die Welt, durch Christum, und müssen ihr dennoch zu Fuße fallen. Also, daß es alles widersinnig geht, was GOtt macht. Wir sollen sagen, daß wir aller Dinge Herren sind, und sie lassen regieren und Oberhand haben. Es geht aber nach dem Geiste zu; der Leib aber muß sich lassen herunter werfen, und unterthan sein.

7. Daß aber Jakob im Texte sagt: „Daß ich Gnade vor deinen Augen fände“, ist gar hebräisch geredet; deutsche und lateinische Sprache redet nicht also; denn wir geben die Gnade allein den Fürsten; ist aber so viel gesagt: Laß mich Gnade finden vor dir, als wir sprechen: Sei mir günstig und hold, zürne nicht mit mir, oder sei mein Freund, ich will auch dein Freund sein. Also grüßte [Luc. 1, 28.] der Engel Mariam die Jungfrau: „Sei gegrüßet, du Gnadenreiche“ oder „Holdselige“; denn was heißt Gnade haben anders
\quellseite{261}
, denn holdselig sein? Als wir von einer Magd sagen: Du feine, freundliche Metze. Welches unsere Lumpenprediger deuten, es heiße voller Gnaden, so voll als eine Blase voll Windes; gerade als wäre Gnade ein Ding, das GOtt in das Herz gieße, wie man Wein in die Kandel gießt, daß [es] oben übergeht. Wir können es nicht besser nennen denn holdselig, denn ein jedermann hold und günstig ist; davon unser deutscher Name Huldreich gemacht ist. Auf hebräisch heißt es Johannes und Hanna.

8. Danach folgt, wie Esau seinem Bruder entgegen zeucht mit vierhundert Mann, welches er nicht böser Meinung thut, wie das folgende Capitel anzeigt. Noch entsetzt sich Jakob, meint, es sei eitel Mord und Tod, da eitel Friede ist. Welch ein elender, schwacher Glaube ist da geworden! Esau ist in dem Sinne und Gedanken: Ich bin reicher und mächtiger denn mein Bruder (wie gesagt ist, daß er im Lande regierte, und ein mächtiger Fürst war), und bin auf ihn zornig gewesen nun zwanzig Jahre. Indem bin ich so gewachsen, daß ich ein Herr bin; er ist noch ein Bettler, hat weder Haus noch Hof, zeucht mit seinem Vieh so bloß auf der Straße, daß man es ihm alle Stunden könnte nehmen. Was soll ich mich nun an ihm rächen? Ich will es ihm schenken, und lassen gut sein. So zeucht er ihm entgegen mit seiner Pracht, daß man sehen soll, wie mächtig er ist, daß es gleich als ein Trotz anzusehen ist. Als sollte er sagen: Ja, wie fein hast du mir meine Erstgeburt genommen! Du hast mir den Segen heimlich vom Vater geraubt, so hat mich GOtt dafür gesegnet. Also vergibt er des Zorns, und zeucht dem Bruder entgegen, daß er ihn empfähe. Aber Jakob ist so blöde und verzagt, daß er nicht anders meint, denn er werde seinen Groll auf ihn und sein Gesinde ausschütten.

9. In solcher Angst und Noth weiß er keine Hülfe noch Zuflucht, denn zu seinem lieben GOtt, und schüttet alle sein Herz vor ihm aus. Das ist auch das beste Stück in diesem Capitel. Wenn der Glaube will schwach werden, und anhebt zu zappeln, ist kein anderer Rath noch Zuflucht denn zu dem, der uns läßt zappeln; wie auch St. Petrus [Matth. 14, 30.] auf dem Wasser schrie, als er wollte sinken: „O HErr, hilf mir!“ Das ist aber eine große Kunst, wenn die Gnade gegeben wird in der Noth, daß er kann fliehen alle menschliche Hülfe und Trost, und allein ergreifen und sich halten an dem Wort und Zusagung. Da ist das Schaf dürftig von Hitze, und trinkt auch. Nun haben wir [Cap. 28, 15.] gehört, wie ihm GOtt hatte zugesagt, er wollte bei ihm sein und ihn behüten, wo er hin käme, und wieder heimbringen; jetzt aber läßt er ihn so tief sinken, und in der Angst stecken, daß er keine Hülfe weiß zu suchen denn bei ihm.

10. So kommt er nun, und legt GOtt seine Noth vor, thut nicht mehr, denn daß er ihm seine Verheißung aufrückt, und thut ein recht stark, feurig Gebet. Denn das heißt nicht beten, wenn man in der Kirche steht, plärrt und plappert; sondern Angst lehrt recht beten; wie man spricht: Hunger ist ein guter Koch. Dieselbe treibt ihn, daß er fürchtet, er müsse mit Weib und Kind umkommen; denn an ihm allein war ihm nicht so viel gelegen. Er hatte die Zusagung, die mußte wahr werden. Ja, wenn er gleich erschlagen wäre, sollten auch aus Steinen Jakobs Kinder werden. GOtt läßt ihn aber in der Angst stecken, auf daß er erzeige die Kraft seines Wortes in unserer Schwachheit.

11. Er reißt ihm das natürliche Herz nicht aus, sondern läßt die Furcht in ihm stecken, mehr für die Seinen denn für sich; er hält sich aber an das Wort und bittet. Es sind nicht viel Worte oder viel Stunden, wie die Werkheiligen ihre Gebete nach der Länge messen; aber siehe, wie gewaltig er bittet. Willst du beten, so komme nur nicht mit vollem Bauche, sondern kriege zuvor eine\footnote{„eine“ fehlt in der Wittenberger und in der Jenaer.} Angst und Noth, die dich dringe zu beten, oder laß es anstehen, und lege ihm solche Noth vor, und ergreife ihn, da du ihn halten kannst, nämlich bei seinem Worte, wie Jakob, und sprich: HErr, da ist das Jammer und Unglück, das mich drückt und drängt, daß wäre ich gerne los; so hast du gesagt [Matth. 7, 7. Luc. 11, 9.]: Bittet, so werdet ihr nehmen, das sind deine Worte, darauf komme ich und bitte, so hast du recht gebetet, und bist gewiß erhört, weil es GOtt zugesagt hat, und die Noth zwingt.

12. Es ist nur der Jammer, daß wir unser Herzeleid nicht fühlen, und nicht zu Herzen lassen gehen, sonst hätten wir Ursache und Noth genug, alle Augenblicke zu beten, sonderlich so wir rechte GOttes Kinder wären, und sähen, daß kein heiliger, edler Name, darinne wir getauft sind, und
\quellseite{262}
darnach genannt, damit er uns ziert und segnet, so greulich und schändlich gelästert wird. Wo sind da Christen, die da klagen und schreien? Also auch hätten wir täglich Ursache, und hohe Zeit zu bitten, daß sein Reich zukomme, weil wir sehen, daß [es] so schwach in uns geht, und vom Teufel und Fleisch und Blut gehindert wird; und so fortan in allen andern Stücken. Das will niemand zu Herzen gehen, welches ein Zeichen ist, daß wir nicht in Angst stehen. Darum habe ich mehr gesagt: Wir finden in den lieben Vätern nicht so tolle, närrische Werke, als wir treiben, sondern eitel ernstliche, rechtschaffene Werke, wiewohl sie gering scheinen.

13. Also hat nun dies Gebet müssen erhört werden; denn es alle Eigenschaften eines rechten Gebets an sich hat, daß es im rechten Glauben auf GOttes Güte und Zuversicht geht, und GOtt sein Wort anzeigt und damit dringt, und darnach mit Ernst seine Noth vorlegt. [Es] ist so kräftig gewesen, wenn auch Esau noch sein Todfeind wäre, und mit aller seiner Macht aufgebrochen, ihn mit allem seinem Gesinde zu schlagen, hätte er doch müssen sein bester Freund werden, oder selbst zu Scheitern gehen. Was nun folgt im Texte von dem Geschenke, wie er es geordnet hat, ist nicht viel von zu reden, ohne daß es anzeigt, daß Jakob ein eben\footnote{eben = eigentlich bedeutend.} Gut habe gehabt, daß er fünferlei Geschenke vorher schickt, welches allein jetzt wohl eines armen Mannes Gut wäre. [Wir] wollen das letzte Stück des Capitels handeln.

\begin{schriftzitat}
V. 24—28. Da rang ein Mann mit ihm bis die Morgenröthe anbrach. Und da er sahe, daß er ihn nicht überwmochte, rührte er das Gelenk seiner Hüfte an; und das Gelenk seiner Hüfte ward über dem Ringen mit ihm verrenkt. Und er sprach: Laß mich gehen, denn die Morgenröthe bricht an. Aber er antwortete: Ich lasse dich nicht, du segnest mich denn. Er sprach: Wie heißest du? Er antwortete: Jakob. Er sprach: Du sollst nicht mehr Jakob heißen, sondern Israel,\footnote{Handglosse: Israel. Israel kommt von Sarab, das heißt kämpfen oder überwältigen; daher auch Sar ein Fürst oder Herr, [und Sarah] eine Fürstin oder Frau heißt, und Israel ein Fürst oder Kämpfer GOttes, das ist, der mit GOtt ringt und angewinnet; welches geschieht durch den Glauben, der so fest an GOttes Worte hält, bis er GOttes Zorn überwindet, und GOtt zu eigen erlangt zum gnädigen Vater.} denn du hast mit GOtt und mit Menschen gekämpft, und bist obgelegen.
\end{schriftzitat}

14. Daß das vorige Gebet erhört ist, erweist sich wohl aus dem Ringen und Kampfe, den er hier treibt. Da siehe abermal, wie sich GOtt stellt, wenn er jemand erhört, daß dies Sprüchwort wahr ist: GOtt begegnet uns oft, wer ihn könnte grüßen. Jakob ist blöde und verzagt vor dem Bruder, so kommt GOtt und will ihn stärken, und macht ihn noch stärker. Das wäre mir ein guter Trost, wenn ich im Schlamme arbeite, und er steckte mich vollends gar hinein. GOtt schreckt ihn und macht ihm angst, er aber hebt an zu rufen und beten; so kommt er, und will ihn gar erwürgen.

15. Wie aber der Kampf sei zugegangen, werden wir nicht mit Worten erreichen. Es wird wohl bleiben allein bei denen, die solches erfahren und geschmeckt haben, sonst versteht es niemand. Doch können wir nicht wissen, was es gewesen sei. So lasset uns je das lernen, wie er sich stellt, wenn er erhört, nämlich also, als wollte er uns gar verderben. So that er mit den Kindern Israel auch [2 Mos. 5, 9.]; da er sie aus Egypten führen wollte, ließ er sie durch Pharaonen ärger angreifen und plagen, denn je vormals. Item [2 Mos. 14, 8. 9.], da er sie heraus brachte an das rothe Meer, und sie meinten, sie wären nun hindurch, kam Pharao hinter sie, daß sie nirgendes aus kommen konnten, daß ihnen allen der Tod vor Augen war, und die Egyptier sicher waren, sie hätten sie so gewiß, als einen Dieb an einem Stricke oder Bande.

16. Also auch mit Christo: Als er ihn wollte zu einem Könige machen, hob er es an mit großem Wunder und Kraft; aber da er es am besten sollte sein, läßt er ihn am Kreuze sterben, als einen verzweifelten Bösewicht; doch schafft er es so, daß er eben im selbigen, da alle Vernunft an ihm verzweifelt hatte, hervor kommt und ewig König wird. Wie den Kindern Israel, als sie mitten im Tode standen, ohne alle Hülfe und Rath, reißt er das Meer voneinander, daß sie trocken hindurch gehen, daß, wenn gleich das Meer nicht Raum gegeben hätte, ehe hätten müssen eine Wolke herab sein gelassen, und müssen Platz machen. Also fährt er mit uns auch; wenn wir ihn anrufen, errettet zu werden vom Tode, so führt er uns erst hinein. Solches thut er nun\footnote{Jenaer: nur.}
\quellseite{263}
darum, daß er die Vernunft zu Schanden mache, welche nicht glauben, sondern wissen will, wie, wo und wann; auf daß der Glaube Raum habe, und lasse GOtt machen.

17. Also hier auch hofft nun Jakob, er solle errettet werden von seinem Bruder: so kommt er erst in die Noth, daß dieser mit ihm kämpft, und will ihm Leib und Leben nehmen; er mußte auch am ersten nicht, wer der Kämpfer war. Nun denke du selbst, wenn ein blöder, erschrockener Mensch soll mit einem Geiste streiten, wie ihm zu Muthe ist; wenn er ein wenig einen Teufel hört rauschen, erblaßt und biomet\footnote{biomen = beben.} er, daß ihm die Welt zu enge wird. Wie greulich wäre es nun, wenn einer allein mit dem Teufel also kämpfen sollte, und sich wehren? Da bildet er [sich] solche Gedanken ein, davon kein Mensch sagen kann. Da lehrt die Natur wohl, was für Angst und Noth Todeskampf ist.

18. Also ist es mit ihm auch gestanden. Er hat gesehen, daß sich der an ihn legte, und mit ihm rang, der ihn mit einem Finger könnte würgen, und sein nicht konnte los werden. So hat er auch nicht gewußt, ob es ein guter oder böser Geist wäre, denn sonst wäre es kein Kampf gewesen, so er gewußt hätte, daß es GOtt, oder ein guter Engel wäre, sondern mehr ein Trost, und wäre ihm nicht sauer worden. Er meint aber nicht anders, denn es sei ein Geist, der ihn denke zu erwürgen. Zuvor sorgt er für Weib und Kind, daß sie ihm nicht umkämen; jetzt muß er es alles lassen fahren und vergessen, und denken, wo er bleibe. Es wird ihm nicht allein das Leben haben gegolten, sondern mit einander beide Leib und Seele, daß er mit ihm ringen mußte um zeitliches und ewiges Leben. [Er] ist da gestanden und hat ihm vorgegeben, als wollte GOtt sein nicht, und müßte des Teufels sein. Doch ist es ein richtig Bild gewesen, daß er hat angenommen, und mit ihm gefochten, daß nicht allein der Leib müde und matt, sondern auch der Muth und Seele im Schrecken und Zagen gelegen ist.

19. Was ist nun die Kraft, damit er sich schützt und wehrt, und endlich gewinnt? Es ist je\footnote{„je“ fehlt in der Wittenberger.} eine schlechte, ohnmächtige Kraft unseres Leibes gegen einen Geist; jedoch, wenn der Muth bleibt und nicht gebrochen wird, so bleibt der Leib auch stark. So lange das Herz nicht verzagt, fließt seine Macht und Kraft auch in [den] Leib; aber wenn der Muth dahin ist, so ist es alles aus, so, daß der Leib nicht kann auf den Füßen stehen.

20. Also ist das Herz noch in ihm fest geblieben auf GOttes Wort, wie er sprach: „Du hast gesagt: Ich will dir wohlthun, und [dich] zum großen Samen machen.“ Das Wort ist das Leben, Stärke und Kraft des Mannes, das hat er im Herzen gefaßt und so fest gehalten, daß es mußte wahr bleiben, und [hat] gedacht: Der will dich auf der Stätte erwürgen, geht so mit mir um, als sei er von GOtt geordnet; nun laß ihn machen, GOtt hat gesagt, er wolle mich wieder heim zu Lande bringen; das muß geschehen, sollte gleich Himmel und Erde reißen; es komme Teufel oder Engel, oder er selbst, und schlage mir es anders vor, so glaube ich es nicht.

21. Da hat er die alte Haut fein müssen ausziehen, und sich brechen, hat nichts gefühlt in Leib und Leben, denn daß er sich hat gestönt\footnote{stönen = stützen.} auf die Wahrheit, die nicht trügen konnte. Wenn wir auch in Anfechtung versucht hätten, wie GOttes Wort stärkt und einen Muth macht, so könnten wir es verstehen; weil es aber nicht versucht ist, so ist es kalt, und schmeckt nicht. Aber wenn es kommt, daß man Hände und Füße muß lassen gehen, und kann nur das Wort erhalten, so sieht man, was es für eine Kraft ist, daß kein Teufel so stark ist, daß er es umstoße, ob er es wohl beißt und will’s verschlingen; [es] ist ihm aber ein glühender, feuriger Spieß. Es ist ein geringes Wort, das weder Schein noch Ansehen hat; doch, so es durch den Mund oder Rede in das Herz gefaßt wird und versucht, wird man es gewahr, was es kann.

22. Also ist Jakob gestanden und hat ritterlich gekämpft, hat durch die einige Kraft Leib und Seele errettet und gewonnen, wie der Text sagt, daß der Mann sahe, daß er ihn nicht konnte überwältigen, nämlich, daß er solchen Muth und Sinn hätte, und nicht verzagt war. Also ist er durch und in der Schwachheit zur höchsten Stärke gekommen, als ihn das Zappeln in der letzten Noth trieb, allein am Worte zu halten. Das ist, wie oft gesagt, GOttes Kunst, daß er aus Nichts Etwas macht, aus Sünde Frömmigkeit, aus Tod Leben und Seligkeit. So muß der gute Jakob mit aller seiner Kraft zu
\quellseite{264}
Boden sinken, und neue Kraft und Stärke überkommen, daß, der zuvor nicht überwinden konnte überwinden, jedem Engel abgewinnt.

23. Denn man muß das bleiben lassen, daß es ein rechter Kampf gewesen ist, daß ihn der Engel mit aller Kraft angegriffen hat, nicht mit ihm geschimpft; so hat er auch mit aller Kraft gefochten, und hat überwunden; nicht daß er ihn niedergeschlagen habe, sondern also, daß er so feste ist stehen geblieben, daß ihn der Engel nicht konnte dämpfen. Denn GOttes Wortes Kraft ist über alle Welt, Teufel und Engel.

24. Doch hat der Kampf nicht lange müssen währen; denn der Text sagt, daß er in der Nacht sei aufgestanden, und habe die Weiber und Kinder mit allem Heer übergeführt; [es] mag irgend eine Stunde oder halbe vor Tage gewährt haben. Nun ist die Nacht an ihr selbst schrecklich; so läßt ihn GOtt allein, daß er keinen Trost oder Hülfe kann suchen; [er] muß bloß und nackt stehen und kämpfen. Solche Exempel findet man nicht viel mehr in der Schrift; [es] ist hoch und seltsam, aber darum geschehen und geschrieben, daß man lerne, wie unaussprechliche Kraft das Wort GOttes habe.

25. So spricht nun der Engel, da der Kampf aus war: „Du sollst nicht mehr Jakob heißen, sondern Israel; denn du hast mit GOtt und mit Menschen gekämpft, und bist obgelegen.“ Der Name Israel kommt von dem Worte Sarah und El. Denn die Juden, wenn sie wollen nomina propria machen, setzen sie gemeiniglich den Buchstaben Jod, das ist J, vorne an. Als wenn sie sagen Jsaak, Jismael, Jakob, Josaphat, Jezechias, Jeheziel. Sarah aber heißt auf deutsch kämpfen und überwältigen, und El heißt GOtt, so wird daraus Israel, das ist, ein Gotteskämpfer, qui luctans cum Deo vincit, daß es beides in sich begreife, kämpfen und gewinnen. Wie ist es denn möglich, daß man Gott könne abgewinnen? Ist er nicht allmächtig? Zum ersten zeigt der Text damit an, was es für ein Kampf gewesen sei, daß ihm hat gegolten das höllische Feuer, und er auch die Hölle gefühlt hat, und [es] läßt sich ansehen, als sei der Engel GOtt selbst gewesen. Der Prophet Hoseas hat es [Cap. 12, 4.] auch fast also angegeben, und Christus ist auch enthalten in solchen Figuren vorgebildet, daß er sich hat lassen als einen Engel sehen. Als droben im 18. Capitel, V. 20., von dem Engel, der mit Abraham redete. Item im andern Buche [Cap. 3, 4.] folget, wie der Engel Mosi erschienen in dem Busche, und sagt, wie GOtt mit ihm geredet habe. Es sei ihm aber, wie ihm wolle, so ist es doch an GOttes Statt geschehen.

26. So ist nun mit GOtt ringen nichts anders, denn mit dem zornigen GOtt ringen, der sich wider den Menschen setzt als ein Feind [Luc. 22, 44.]; [er] will nicht allein Richter sein, sondern, daß noch greulicher ist, auch der Stockmeister sein und um das Leben bringen. Wenn er nach dem Leben steht, greift er nicht nach der Haut, sondern hinein, daß das Mark verschmachtet, und die Beine so mürbe werden wie das Fleisch, wie auch Christo widerfuhr im Zelgarten [Matth. 26, 37. 38.]. Darum muß man es groß machen. Wenn er einen Menschen angreift, ist [es] hart und schwer, daß [es] niemand begreift, denn der es versucht; weil er selbst mit dem Menschen fechten will, so ist eitel Unfriede und der Hölle Angst da; dazu sind ihm alle Creaturen der Tod, denn sie halten es alle mit GOtt. Also nimmt er dem Menschen das Herz, daß er nirgend nichts sieht, das auf seiner Seite stünde. Wie soll er denn in solcher Angst gewinnen?

27. Die Sophisten können nichts dazu reden, denn es lautet in ihren Ohren nicht, als wider den Artikel, daß er allmächtig ist; darum sich auch der lateinische Text gescheuet hat zu sagen: Du hast mit GOtt gekämpft, und bist obgelegen. Es geht aber also zu, wie ich sonst gesagt habe: GOtt hat der Welt zweierlei Wort gegeben, das Gesetz, das da zürnt und würgt, und das Evangelium, damit er tröstet und lebendig macht. Wenn nun das Wort fällt, dadurch er uns Gnade zusagt, soll man nimmermehr davon lassen, das Gesetz wende es vor oder hinter; ob auch darnach eitel Ungnade, Hölle und Sünde daher fallen, so laß dir es in keinen Weg nehmen, und sprich nur frei: Nach dem Evangelio kommt kein Zorn. Denn so bald das angegangen ist, hat sich GOtt recht zu erkennen gegeben; denn er ist von Natur eitel Güte, darum mußt du ohne Unterlaß daran halten. Fällt aber ein ander Wort darauf, so denke, es müsse entweder falsch sein, oder GOtt wolle dich versuchen. So thun nun Gott, denn er will die Seinen völlig stark machen: [er] legt sich wider das gütige, tröstliche Wort, und will sehen, wie fest sie daran hangen, stellt sich, als wollte er ihnen nimmer
\quellseite{265}
Gutes thun; so fühlt der Mensch, daß es GOtt thut; das thut auch am allerwenigsten. Wenn man meinte, daß [es] der Teufel oder ein Mensch thäte, so hätte es nicht so große Noth, aber wenn das Gewissen sagt, daß [es] GOtt thut, so ist Angst und Noth da.

28. So heißt nun das „GOtt überwinden“ nicht: seine Gewalt überwinden; sondern dasjenige, das er in unserm Gewissen ist und gefühlt wird, überwinden; wie die Schrift redet, daß sich GOtt verwandelt, wenn wir verwandelt werden. Er ist ohne Wandel an ihm selbst, noch verwandelt er sich uns so wunderlich; das macht, daß sich unser Gewissen verwandelt; wie der 18. Psalm, V. 26. 27., sagt: „Bei den Heiligen bist du heilig, und bei denen ohne Wandel, bist du ohne Wandel, und bei den Verkehrten, bist du verkehrt“; er bleibt immer gütig, noch ist in meinem Gewissen nicht anders, denn daß er zornig ist. Also ist er den Verdammten nichts denn eitel Zorn, straft sie nur mit ihrem eigenen Gewissen.

29. Also wäre es Jakob auch gegangen, wenn er das Wort hätte lassen fahren; weil er aber nicht schließen will nach seinem Fühlen, sondern hält fest, und ließ sich ehe zerreißen, hat er gewonnen. Wenn ich ihn also in mir überwinde, so habe ich GOtt überwunden, dadurch, daß ich das Wort von seiner Güte ergreife und halte, und schlage das hinweg, das ihn zornig will machen. Also überwindet man, nicht seine Majestät, sondern sein Werk, das er an uns thut. Derhalben lässet uns lernen, daß solches uns zu unterweisen geschrieben ist, ob uns auch solches begegnet, daß wir müßten GOtt also zu halten, daß wir auch Israel würden.

\begin{schriftzitat}
V. 29—32. Und Jakob fragte ihn, und sprach: Sage mir an, wie heißest du? Er aber sprach: Warum fragst du, wie ich heiße? Und er segnete ihn daselbst. Und Jakob hieß die Stätte Pniel,\footnote{Randglosse: Pniel oder Pnuel heißt GOttes Angesicht oder Erkenntniß; denn durch den Glauben im Streite des Kreuzes lernt man GOtt recht erkennen und erfahren, so hat es denn keine Noth mehr, so geht die Sonne auf.} denn ich habe GOtt von Angesicht gesehen, und meine Seele ist genesen. Und als er vor Pnuel über kam, ging ihm die Sonne auf, und er hinkte an seiner Hüfte. Daher essen die Kinder von Israel keine hohle Ader auf dem Gelenk der Hüfte, bis auf den heutigen Tag, darum, daß die hohle Ader an dem Gelenk der Hüfte Jakobs gerühret ward.
\end{schriftzitat}


30. Da ist nun der Kampf aus, und als der Engel von ihm kommt, nennt er die Stätte darnach Pniel, das ist, GOttes Angesicht. Was heißt aber GOttes Angesicht? Die Schrift hat dies Wort in großem Brauche hin und wieder. Die Propheten treiben es und preisen es mächtig hoch, und sagen, man müsse GOttes Angesicht in diesem Leben sehen; wie David in 4. Psalm, V. 7., sagt: „Erhebe über uns das Licht deines Angesichtes.“ Item, Psalm 67, 2.: „GOtt sei uns gnädig und segne uns, und lasse uns sein Angesicht leuchten.“ So ist nun GOttes Angesicht sehen, wenn man ihn erkennt ohne alle Larve, gerade, als wer eines Menschen Angesicht sieht, der kennt ihn recht. Denn wenn man eines Rücken sieht, kann man nicht wissen, ob er Feind oder Freund ist. Denn wer mit mir zürnt, wendet das Angesicht von mir; wiederum, wer mir günstig und hold ist, läßt mich ihm unter [die] Augen sehen. Darum heißt GOttes Angesicht nichts anders, denn die Offenbarung und das klare Erkenntniß seines Wesens, wie er gesinnt ist, das ist, seiner Güte. Daraus kommt denn auch Erkenntniß unser selbst, also, daß wir von allem unserm Wesen nichts halten.

31. Darum sehen alle, die da glauben, GOttes Angesicht ohne Unterlaß, das ist, sie erkennen, wie GOtt eitel Güte sei, und sie mit barmherzigen Augen ansieht. Das meint nun Jakob, da GOtt aufhört, und sich nicht mehr zornig stellt, müßte es eitel Gnade sein; [er] sieht und erkennt ihn nun recht wohl, wie es wahr ist, daß er geredet hatte, er sollte sein liebes Kind sein; davon wird er so voll Freude und Lust, so groß vor die Betrübniß war, daß er sagt: Diese Stätte muß GOttes Angesicht heißen, darum, daß ich GOtt allda recht gesehen habe, und meine Seele ist genesen, daß ich nun niemand mehr fürchte; denn hernach fürchtete er sich nicht mehr vor Esau, ob er sich wohl vor ihm demüthigte.

32. Zuletzt sagt der Text, wie der Engel, als er mit ihm kämpfte, das Gelenk seiner Hüfte anrührte, daß es über dem Ringen verrenkt und er an der Hüfte hinkend ward; daher die Kinder Israel keine hohe Ader essen auf dem Gelenk der Hüfte. Was er für eine hohe Ader meine, weiß ich nicht; ich halte, er meine die Ader, die auf dem Gelenk herüber geht. In dem Ringen bringt er das davon, daß er hinten oder lahm in der Hüfte wird, also, daß es dennoch ein leiblich
\quellseite{266}
Ringen gewesen ist. Daher ist die Hüftader geheiliget oder verbannt, daß man keine gegessen hat.

33. In der Figur hat GOtt so lange zuvor angezeigt, daß GOttes Kinder nicht sollen gezählt oder gerechnet werden nach Fleisch und Blut, sondern nach dem Geist und Wiedergeburt. Darum mußte der Erzvater des ganzen GOttes Volks an der Hüfte lahm werden, nachdem er GOtt gesehen hatte, und ein geistlicher Mann worden war; denn dem Theile des Leibes, nämlich den Lenden oder Hüften, pfleget die Schrift aller Menschen Geburt zu geben.

34. Essen aber ist nichts, denn predigen und lehren, daß [es] so viel gesagt sei: Man muß dem jüdischen Volke so predigen und sagen, daß niemand darum GOttes Kind sei, daß er von [den] Patriarchen geboren oder gekommen ist. Der Ruhm ist gar niedergelegt, daß ein jeglicher für sich, ohne allen Ruhm, lauter auf GOttes Gnade bestehe, und dadurch vor GOtt gehe; also, daß die Juden bekennen, wollen sie zu GOtt kommen, daß sie es nicht daher haben, daß sie von Abrahams Samen sind, sondern daher, davon es Abraham hat, das ist, daß sie denselbigen Glauben haben. Dadurch fällt nun nieder alles, was man von unsern Werken und freiem Willen predigt. Solches hat Gott zuvor gebeutet, weil es zukünftig war, daß sich das Volk der Väter hoch annehmen, und auf sie trotzen würde; also, daß die Geburt nicht mehr vor ihm gelten sollte, doch eine Zeitlang stehen, bis der verheißene Christus käme, und darnach aufhören. So haben wir dies edle Capitel, darinne du siehst den wunderbaren Rath, ja Gott braucht mit seinen Heiligen, uns zum Troste und zu einem Exempel, daß wir solches täglich im Gedächtnisse behalten, ob er also auch mit uns spielete, daß wir gerüstet dazu wären.

\section{Das dreiunddreißigste Capitel}
\begin{schriftzitat}
V. 1—20. Jakob hub seine Augen auf, und sahe seinen Bruder Esau kommen mit vierhundert Mann. Und theilete seine Kinder zu Lea, und Rahel, und zu [den] beiden Mägden, und stellte die Mägde mit ihren Kindern vorne an, und Lea mit ihren Kindern hernach, und Rahel mit Joseph zuletzt. Und er ging vor ihnen her, und bückte sich siebenmal auf die Erde, bis er zu seinem Bruder kam. Esau aber lief ihm entgegen und herzte ihn, und fiel ihm um den Hals, und küssete ihn, und sie weineten. Und hub seine Augen auf, und sahe die Weiber mit den Kindern, und sprach: Wer sind diese bei dir? Er antwortete: Es sind Kinder, die GOtt deinem Knechte bescheret hat. Und die Mägde traten herzu mit ihren Kindern, und bückten sich vor ihm. Lea trat auch herzu mit ihren Kindern, und bückten sich vor ihm. Darnach trat Joseph und Rahel herzu, und bückten sich auch vor ihm. Und er sprach: Was willst du mit alle dem Heere, dem ich begegnet bin? Er antwortete: Daß ich Gnade fände vor meinem Herrn. Esau sprach: Ich habe genug, mein Bruder, behalte, was du hast. Jakob antwortete: Ach nicht, habe ich Gnade funden vor dir, so nimm mein Geschenk von meiner Hand. Denn ich habe dein Angesicht, als sähe ich GOttes Angesicht, und laß dir es wohlgefallen von mir. Nimm doch den Segen von mir an, den ich dir zugebracht habe; denn GOtt hat mir es bescheret, und ich habe alles genung. Also nöthigte er ihn, daß er es annahm. Und er sprach: Laß uns fortziehen und reisen, ich will mit dir ziehen. Er aber sprach zu ihm: Mein Herr, du erkennest, daß ich zarte Kinder bei mir habe, dazu Schafe und Rinder, die noch jung sind. Wenn sie einen Tag übertrieben würden, würde mir die ganze Heerde sterben. Mein Herr ziehe vor seinem Knechte hin, ich will mäßig\footnote{Handlosse: Mäßig. Merke, daß Rechtsgläubige und Werkheilige nicht können mit einander wandeln; denn die Gläubigen fahren säuberlich, mit stillem Geiste, aber die Werkheiligen fahren mit Vermessenheit ihrer Werke in Gottes Gesetzen.} hinnach treiben, darnach daß Vieh und die Kinder gehen können, bis daß ich komme zu meinem Herrn in Seir. Esau sprach: So will ich doch bei dir lassen etliche vom Volke,
\quellseite{267}
das mit mir ist. Er antwortete: Was ist es vonnöthen? laß mich nur Gnade vor meinem Herrn finden. Also zog des Tages Esau wiederum seines Weges gen Seir. Und Jakob zog gen Suchoth, und baute ihm ein Haus, und machte seinem Viehe Hütten, daher heißt die Stätte Suchoth. Darnach zog Jakob gegen Salem, zu der Stadt Sichem, die im Lande Canaan liegt, nachdem er aus Mesopotamia kommen war, und machte sein Lager vor der Stadt. Und kaufte ein Stück Ackers von den Kindern Hemor, des Vaters Sichem, um hundert Groschen, daselbst richtete er seine Hütten auf. Und richtete daselbst einen Altar zu, und rief an den Namen des starken GOttes Israel.
\end{schriftzitat}

1. Das ist aber eine schlechte Historie, die da wenig scheint, als die andern gemeiniglich auch; denn welche hoch sind, als die nächste vorgeschriebene, sind kaum zu merken. Nun habe ich genug gesagt, warum Moses so viele unnütze Worte macht, und beschreibt die heiligen Väter mit so geringen Geschichten, so die Welt nicht ansähe. Summa: Es geht dahin, daß wir mit unserer Klugheit zu Narren werden und GOtt lernen erkennen, daß er mit solchem Dinge umgeht, so vor der Welt nichts ist.

2. Also hier auch schreibt er, wie die zwei Brüder auf einander stoßen. Eine schlechte Historie; aber nichtsdestoweniger dabei angezeigt, wie GOtt die Weltleute läßt hoch herfahren, aber die Seinen geringe und ungeachtet. Denn was GOttes ist, muß zum ersten gar nichts sein und keinen Schein haben. Esau führet gewaltig daher mit viel Heeres; Jakob allein mit Weib und Kind, welcher doch sollte der Größte sein, laut der Verheißung; aber GOtt läßt es immer so hingehen, bis zu seiner Zeit. Er hat den Esau hervor gezogen und aufgeworfen, daß er empor schwebte bei zweihundert Jahren; indes mußte Jakob immer schwach und gering bleiben; aber da die Zeit kam, hob er ihn mit seinem Volke weit über Esau, daß sie ihm mußten dienen. Welches ist eine Figur und Vorbild Christi; der lag drei Tage im Grabe, und ließ die Welt und Teufel: Gewonnen! schreien, als wäre es mit ihm aus; aber in demselben brach er hervor, und ward ein HErr der Welt und aller Creaturen. Also auch, seine erste Christenheit, die Märtyrer, ließ er ein Jahr oder zwei oder dreihundert unterliegen und sterben; aber bald hernach setzt sie GOtt über alle Gewalt der Welt, so, daß ihnen auch Könige und Fürsten müssen zu Fuße fallen.

3. Als nun Jakob seinen Bruder Esau sieht herkommen, spricht der Text, theilete er seine Kinder, und machte eine Ordnung. Denn wiewohl er sich nicht mehr so fürchtet vor dem Bruder, doch traut er ihm nicht gar viel, sondern weil er ungewiß war, wie es GOtt mit ihm machen wollte, stellte er sich in GOttes Gewalt, und wagte es darauf mit allem, das er hat; denn er weiß, daß Esau in GOttes Hand steht, und ihm nicht schaden kann, er wolle denn. Dazu ist er sicher, daß er einen gnädigen GOtt hat, der ihn schützen und retten will; aber wie, und [auf] waserlei Weise, weiß er nicht. Doch thut er, was er kann, versucht Gott nicht, und bleibt doch im richtigen Glauben. Denn wenn er verzagt hätte, würde er gelaufen haben, und beiseite ausgegangen sein; [er] steht aber still, und wagt sich zum ersten hinan, ob jener ja durch GOttes Verhängniß schalten wollte, daß er vorne an der Spitze stände, ob er möchte die andern retten; setzt sein Leben frei dahin für sein Weib und Kinder.

4. Ist denn das nicht GOtt versucht? Nein, denn der Glaube ist recht, dadurch er sich GOtt befiehlt, [er] sucht nicht, wie ihm zu helfen sei, denkt also: Ob er uns gleich schlüge, müßte dennoch GOttes Wort wahr bleiben; besorgt gleichwohl die Untreue des Bruders, und braucht der Klugheit und Vernunft, die GOtt gegeben hat. Denn das lehrt die Natur, was man lieb hat, daß man dasselbe zum ersten rette.

5. Also richtet er Rahel und Joseph, das rechte Kind, zu hinterst, darnach Lea, und vornean die Mägde, sich aber schlägt er in die Schanze. Das ist die rechte brüderliche Liebe, damit thut er so viel, als die Vernunft lehrt; nicht als die, so GOtt versuchen, und wollen unter dem Haufen ohne Ordnung daran gehen, und denken: Werden wir geschlagen, so sind wir geschlagen.

6. Das gilt aber nicht; denn Glaube und Versuchung leiden sich nicht mit einander. Die da GOtt versuchen, wenden vor, als wollten sie glauben; aber sie fahren ohne GOttes Wort und Befehl. Als, so einer auf einen hohen Thurm und Spitze wollte treten und herab fliegen. Item, daß ein fauler Esel nicht arbeiten wollte, und warten, daß ihn GOtt vom Himmel ernährte.
\quellseite{268}

7. Darum muß man die beiden also scheiden, daß du den Glauben setzest in die Dinge, so über dich sind; Versuchung aber meidest in denen, die unter dir sind. GOtt hat dir den Leib, fünf Sinne und Vernunft gegeben; item, Weib und Kind; welches alles in deiner Hand ist, deß sollst du gebrauchen, so viel du kannst; kannst du aber nicht, magst du es denn lassen anstehen und warten, was dir GOtt gibt. Darum sind wir geschaffen, daß wir regieren über die Thiere auf Erden, Weib, Kind und Haus, und die Nahrung suchen mit Arbeit aus der Erde und sonst; damit sollen wir umgehen. Was aber über uns ist, sollen wir immer lassen gehen, was da geht; wie Jacob hier thut. Denn, daß er die Ordnung macht, war nicht über ihn,\footnote{Zenger = ihm.} sondern unter ihm, nämlich eitel vernünftig Ding; aber der Glaube geht nur mit den Sachen um, die über die Vernunft sind. Also, daß er nur nicht in seiner Hand, sich und die Seinen zu retten vor Esau, sondern allein GOttes Werk; darum braucht er darinne des Glaubens allein, und läßt GOtt walten. Darum sage ich noch einmal: Versuchung GOttes kann nicht geschehen in den Dingen, so über mir sind und den Glauben angehen; wo ich nicht kann helfen, da muß ich glauben. Wo ich mir aber helfen kann, und habe GOttes Güter vorhanden, da soll ich ihn nicht versuchen; denn da darf ich nicht glauben, weil ich es vor Augen und in der Hand habe.

8. Man liest ein Exempel in vitis patrum von zwei Brüdern, die auf eine Zeit mit einander wanderten über Feld; und wiewohl sie nicht zu essen hatten, beschlossen sie doch, nicht Essen zu nehmen von den heidnischen, bösen Leuten, sondern [zu] warten, bis ihnen GOtt Brod vom Himmel gäbe. Da sie nun lange hungerten, sprach der eine: Es wird uns zu lange, ging hin, und nahm, was man ihm gab, und blieb lebend; der andere aber blieb auf seinem Vorsatze, und starb auch darüber. Nun schließe du: Wo ist dieser hingefahren? GOtt gab ihm, daß er wohl konnte dem Hunger rathen, darum durfte er nicht ein anderes warten, und die Baarschaft\footnote{Baarschaft = das Vorhandene an Geld oder Gütern.} lassen liegen; wenn aber niemand gewesen wäre, der ihm hätte ein Stück Brod gegeben, dann wäre es recht gewesen, [es] von GOtt zu [er]warten.

9. Also fasse den Unterschied zwischen versuchen und glauben. Der Glaube soll nicht eher kommen, es sei denn GOttes Gebot und Wort da, welches dann angeht, wenn man in die Noth kommt, daß man ihm selbst nicht helfen kann, als, im Tode und andern Anfechtungen; sondern wo GOtt solch Gut gäbe, soll man ihm danken und zu sich nehmen. Also hat Jakob allein gebraucht seiner Kunst auf das beste, als er wußte, doch im Glauben, daß er nicht wußte, wie es sollte hinaus gehen, wer lebendig oder todt bleiben würde; befehlt GOtt bloß dahin, daß er schütze und rette. So geht er richtig hindurch, weder zur Linken noch zur Rechten, weder im Unglauben noch Versuchung.

10. Lasse nun dies einfältige und schlechte Historien sein, das keinen Schein hat, daß die Väter in solchen losen Werken einhergehen; aber wenn man das Herz sieht und den Glauben, darnach das Werk gerichtet ist, da wird man gewahr, wie es treffliche Exempel sind. Darum wollen solche Historien nicht nach Werken angesehen sein, wie die Welt thut; sonst ist es kalt und unnützlich; aber nach dem Glauben sind es eitel theure Schätze.

11. Weiter sagt der Text: „Und er ging vor ihnen her, und bückte sich siebenmal auf die Erde, bis er zu seinem Bruder kam.“ Darinne, habe ich gesagt, erzeigt er rechte brüderliche Liebe, daß er sein Leben dahin setzt mit gutem, festem Glauben, für Weib und Kind. Das heißt ja nicht gesucht, was sein ist, noch um seinetwillen leben und sterben. So geht das Werk in starkem Glauben und vollkommener Liebe.

12. Wie aber, daß er sich siebenmal vor seinem Bruder bückt, ist das nicht große Ehre? Wie schlecht die Werke sind, darinne er geht, noch zeigen sie so große Tugend. Denn, ist das nicht eine große Demuth, weil er die Zusagung von GOtt hat, daß sein Same gesegnet sollte sein? darum er sich hätte mögen brüsten und sagen: Bruder, du mußt mir zu Füße fallen. Hat er denn auch geheuchelt oder gelogen? Nein; denn St. Paulus hat es auch geboten Röm. 12, 10., daß einer dem andern mit Ehrerbietung zuvorkomme; wiewohl es Esau nicht wieder thut, wird er doch dadurch bewegt.

13. Darum, wie groß Gut wir haben in Christo, und Herren sind auch des Teufels, ist es dennoch GOttes Wille, daß wir uns demüthigen vor denen, so in der Welt etwas sind,
\quellseite{269}
nämlich die Obrigkeit und Herrschaft, wenn du gleich so heilig wärest als der Jakob.

14. Es ist auch nie kein Papst so heilig gewesen, als der Jakob, noch läßt er ihm nicht die Füße küssen, sondern fällt zu Fuße dem Tyrannen, gibt ihm Ehre, mehr denn ihm gebührt, und thut es von rechtem Grunde und gutem Herzen; denkt also, wie er selbst sagt: „Ich habe dein Angesicht gesehen, als sähe ich GOttes Angesicht.“ Wie lügst du so, Jakob? Sollst du zu dem unheiligen Menschen also sagen? Es ist alles dahin gerichtet, weil Esau ein großer Herr war, weit über Jakob, darum sieht er ihn an, als seinen Herrn, und lügt nicht, sondern [es] ist sein Ernst; und zum Exempel. Ist jemand im Regimente, den ist man schuldig zu ehren, nicht um seinetwillen, sondern darum, daß [es] GOttes Ordnung ist [Röm. 13, 1. 7.]. Was fragt Jakob darnach, ob Esau seiner Obrigkeit mißbraucht? [Er] läßt ihm gleichwohl die Ehre, die die Obrigkeit hat, und spricht dir heraus, er habe sein Angesicht gesehen, als GOttes Angesicht.

15. Denn die Obrigkeit sei, wie sie wolle, [so] ist sie nicht von Menschen, sonst wäre sie nicht eine Stunde sicher; wo es GOtt nicht mit Gewalt erhielte, würde sie Herr Dimeß alle todtschlagen. Darum, weil es GOttes Gewalt und Ordnung ist, muß man’s ansehen, als [ob] man GOtt sähe. Wo er sie hinwirft, da soll seine Ehre folgen. Er geht mit seinen Gütern um, als der da Geld in die Napusie wirft: wohin es fällt, da ist GOtt, man brauche sie wohl oder übel, er wird es wohl finden. Fahre du nur fort, demüthige dich, und gib ihr die Ehre.

16. Mehr sind in dem Capitel edle Früchte des Glaubens. Also sagt Salomo in den Sprüchen [Cap. 16, 7.]: „Wenn jemandes Wege dem HErrn wohlgefallen, so macht er auch seine Feinde mit ihm zufrieden.“ Das ist GOttes Wesen und Werk, daß er die höchsten Feinde kann zu Freunden machen, als hier Esau eigentlich Jakobs Todfeind war, hatte auch im Sinne, ihm Schaden zu thun. Da siehst du, wie GOtt dieses Menschen Herz und Gedanken in seiner Hand hat, daß sie gehen müssen, wie er will, und seinen Sinn brechen, daß er selbst nicht weiß, wie er dahin läuft, und fällt ihm um den Hals. Welch ein Wechsel und Wandel ist es, daß er im Tode angesehen mit so viel Mann, und da er hinkommt, hebt er an zu weinen; ist es nicht ein wunderlicher Handel, daß ein solcher Todfeind so bald soll umgekehrt werden?

17. Darum lehrt die Schrift, daß man einen guten Wandel führe, und [wir] uns demüthigen,\footnote{Erlanger: demüthige.} nicht um unserwillen, sondern um der Feinde willen. Denn, wo solche Demuth ist, gibt GOtt Gnade, daß auch die Feinde gebrochen und umgekehrt werden. Daher das Sprüchwort kommt: Ein gut Wort findet eine gute Statt; wenn aber Stolz, und hart gegen hart ist, da wird nichts Gutes aus. Wiewohl auch viel harte Köpfe sind, die weder mit Worten noch Werken zu erreichen sind. Also will nun [G]Ott, daß die Christen sich demüthigen und ehrerbietig seien, auch gegen ihren Feinden; so gibt er den Segen mit zu, daß solche Werke kräftig werden, und bekehren die Herzen.

18. Also gehen beide, Weib und Kind, ihm entgegen, und bücken sich alle vor ihm. Das läßt er geschehen, wiewohl er nun Freund ist; [sie] gewinnen ihm aber damit gar das Herz ab, daß er zufährt und spricht: „Behalte, was du hast, mein Bruder, ich habe genug“; erbeut sich dazu, er wolle ihn geleiten, und als sich Jakob entschuldigt, will er ja ihm seines Volkes einen Theil mitgeben. Das ist uns geschrieben, wenn wir es auch einmal thäten; aber der Teufel ist im Hochmuthe, daß keiner dem andern weichen will. Darum lasset uns solche Exempel behalten für Früchte des Glaubens, wie der heilige Vater glaubt und versucht wird, und sich demüthigt gegen der Obrigkeit.

19. Darnach schreibt Moses, wie Jakob gekommen ist gen Suchoth, und bauete da ein Haus, und seinem Viehe Hütten. Die Stätte ist darum beschrieben, auf daß die Juden, wenn sie hernach das Land einnahmen, wüßten, wo Jakob gewohnt hätte. Darum muß dieser Ort auch den Namen behalten, ist auch hernach eine Stadt daraus geworden, daß der heilige Jakob da gewohnt und gelegen hatte.

20. Darnach ist er gezogen gen Salem, zu der Stadt Sichem. Salem muß nicht Jerusalem sein, wie etliche meinen; denn Sichem ist noch eine gute Meile davon gelegen, da hat er auch eine Weile gewohnt als ein Pilgrim und Fremdling, kauft aber ein Stück Ackers, und richtet seine Hütten an. Warum er von Suchoth gezogen ist, zeigt der Text nicht an, [es] ist aber
\quellseite{270}
aus dem Vorigen wohl zu merken. Denn wir haben oben gehört von Abraham und Isaak, wie sie immerdar haben müssen Fremdlinge sein, und keine gewisse Statt haben, und von einem Orte zum andern weichen, welches ihnen fast schwer worden ist, stetig mit Weib, Kind und allem Gute zwischen Himmel und Erde zu schweben, und nirgend sicher [zu] sein, daß sie zuweilen auch ungeduldig sind worden, immer Gewalt und Unrecht unter Fremden zu leiden. Welches uns immerdar vorgehalten ist, daß man sehe, was das rechte Leben ist, ob es wohl schlecht natürlich Ding scheint.

21. Item, daß er einen Altar hier aufrichtet und ruft GOttes Namen an, haben wir auch gehört, wie die lieben Väter einen Ort gelegt haben, da man zusammen käme, lehrte und predigte, wie man GOtt trauen und anrufen sollte auf die Verheißung von dem Samen, dadurch die Welt gesegnet sollte werden. Denn solch Predigen und Vermahnen ist wohl noth gewesen wider Schwachheit des Glaubens und Ungeduld. Und eben setzt der Text dazu: „Er rief an den Namen des starken GOttes Israel“, anzuzeigen, was, und wie er gepredigt hat, nämlich von dem GOtt, der da stark und mächtig ist, daß er sie schützen und ihnen helfen könnte im Leiden und Kreuz, ob es gleich nicht schiene. So siehst du das einfältige Leben der Väter, und doch so voll Exempel des Glaubens, Liebe und Geduld in den geringsten Werken.

\section{Das vierunddreißigste Capitel}
\begin{schriftzitat}
V. 1—31. Dina aber, Lea’s Tochter, die sie Jakob geboren hatte, ging heraus, die Töchter des Landes\footnote{Landglosse: Töchter des Landes. Was man außer GOttes Wort bei der Vernunft und menschlicher Weisheit sucht, das verderbt gewißlich den Geist und Glauben; darum soll kein Zusatz menschlicher Lehre und Werke zu GOttes Wort gethan werden.} zu besehen. Da die sahe Sichem, des Hemor Sohn, des Heviters, der des Landes Herr war, nahm er sie und beschlief sie, und schwächte sie, und sein Herz hing an ihr, und hatte die Dirne lieb, und redete freundlich mit ihr, und sprach zu seinem Vater Hemor: Nimm mir das Mägdlein zum Weibe. Und Jakob erfuhr, daß seine Tochter Dina geschändet war, und seine Söhne waren mit dem Viehe auf dem Felde, und Jakob schwieg, bis daß sie kamen. Da ging Hemor, Sichems Vater, heraus zu Jakob, mit ihm zu reden, indes kamen die Söhne Jakobs vom Felde, und da sie es hörten, verdroß [es] die Männer, und wurden sehr zornig, daß er eine Thorheit in Israel begangen, und Jakobs Tochter beschlafen hatte; denn so pflegte man nicht zu thun. Da redete Hemor mit ihnen, und sprach: Meines Sohnes Sichems Herz hat Lust zu eurer Tochter; Lieber, gebet sie ihm zum Weibe. Befreundet euch mit uns, gebet uns eure Töchter, und nehmet ihr unsere Tochter, und wohnet bei uns, das Land soll euch offen sein; wohnet und werbet, und arbeitet darinnen. Und Sichem sprach zu ihrem Vater und Brüdern: Lasset mich Gnade bei euch finden; was ihr mir sagt, das will ich geben, fordert nur getrost von mir Morgengabe und Geschenk, ich will es geben, wie ihr es heißt, gebet mir nur die Dirne zum Weibe. Da antworteten Jakobs Söhne dem Sichem und seinem Vater Hemor, und redeten trüglich, darum, daß ihre Schwester Dina geschändet war, und sprachen zu ihnen: Wir können das nicht thun, daß wir unsere Schwester einem unbeschnittenen Manne geben, denn das wäre uns eine Schande. Doch, dann wollen wir euch zu Willen sein, so ihr uns gleich werdet, und alles, was männlich unter euch ist, beschnitten werde; dann wollen wir unsere Töchter euch geben, und eure Töchter uns nehmen, und bei euch wohnen, und Ein Volk sein. Wo ihr aber uns nicht gehorchet, euch zu beschneiden, so wollen wir unsere Tochter nehmen und davon ziehen. Die Rede gefiel Hemor und seinem Sohne wohl. Und der Jüngling verzog nicht, solches zu thun; denn er hatte Lust zu der Tochter Jakobs. Und er war herrlich gehalten über allen in seines Vaters Hause. Da kamen sie nun, Hemor und sein Sohn Sichem, unter der Stadt Thor, und redeten mit den Bürgern der Stadt, und sprachen: Diese Leute sind friedsam bei uns, und wollen im Lande wohnen
\quellseite{271}
und wehren; so ist nun das Land weit umfangen, wir wollen uns ihre Töchter zu Weibern nehmen, und ihnen unsere Töchter geben. Aber dann wollen sie uns zu Willen sein, daß sie bei uns wohnen, und Ein Volk mit uns werden, wo wir alles, was männlich unter uns ist, beschneiden, gleichwie sie beschnitten sind. Ihr Vieh und Güter, und alles, was sie haben, wird unser sein, so wir nur ihnen zu Willen werden, daß sie bei uns wohnen. Und sie gehorchten dem Hemor und Sichem, seinem Sohne, alle, die zu seiner Stadt Thore aus- und eingingen, und beschnitten alles, was männlich war, das zu seiner Stadt aus- und einging. Und am dritten Tage, da sie es schmerzte, nahmen die zwei Söhne Jakobs, Simeon und Levi, der Dina Brüder, ein jeglicher sein Schwert, und gingen in die Stadt tüchtig, und erwürgeten alles, was männlich war, und erwürgten auch Hemor und seinen Sohn Sichem mit der Schärfe des Schwertes. Und nahmen ihre Schwester Dina aus dem Hause Sichem, und gingen davon. Da kamen die Söhne Jakobs über die Erschlagenen, und plünderten die Stadt darum, daß sie hatten ihre Schwester geschändet. Und nahmen ihre Schafe, Kinder, Esel, und was in der Stadt und auf dem Felde war. Und alle ihre Habe, alle Kinder und Weiber nahmen sie gefangen, und plünderten alles, was in den Häusern war. Und Jakob sprach zu Simeon und Levi: Ihr habt es zugerichtet, daß ich stinke vor den Einwohnern dieses Landes, den Cananitern und Pheresitern, und ich bin ein geringer Haufe. Wenn sie sich nun versammeln über mich, so werden sie mich schlagen; also werde ich vertilgt sammt meinem Hause. Sie antworteten aber: Sollten sie dann mit unserer Schwester als mit einer Hure handeln?
\end{schriftzitat}

1. Da siehst du abermal, daß die Schrift nichts Gutes von dem guten und frommen Jakob schreibt, das ihm begegnet sei, und seine Legende uns immer vorhält, wie er aus einem Unglück und Elend in das andere geworfen ist; als wir gesehen haben, wie sauer es ihm geworden war mit den Weibern und Schafen, als er zwanzig Jahre diente. Item, wie er sich vor Esau fürchtete. Jetzt, als er nun meint, es soll des Unglücks ein Ende sein, widerfährt ihm das Herzeleid, daß seine Tochter geschändet wird, und seine Söhne auch ein Unglück anrichten, daß er es mit niemand mehr gut hat im Lande. Also geht seine Legende hinaus bis auf das Letzte, daß ihm GOtt immer ein Kreuz über das andere auf den Hals legt; wie wir weiter hören werden.

2. Also ist dies nicht der geringsten Leiden eines, daß er die Schmach und Schande leiden muß, die ihm mit solchem Muthwillen widerfährt, als müsse er es gut lassen sein, und könne sich nicht rächen, weil es der Junker im Lande thut, und denkt, es liege nicht viel daran, ob er ihm seine Tochter nehme und zu Schanden mache. Aber das ist noch das Größte, daß er mit seinen Kindern die Schande muß tragen, daß das ganze Land von ihm sagt, und jedermann seines Unfalls lacht, und muß den Spott zum Schaden haben.

3. Hier denke du, wie sein Glaube versucht ist, daß er hätte mögen sagen: Wo ist nun die große Verheißung, daß mir GOtt wolle wohlthun und mich schützen? Heißt das wohlgethan, daß er mich läßt eine solche Schande übergehen? Der Text redet nicht davon, wie ihm zu Muthe gewesen sei, sondern gibt es einem jeglichen selbst zu bedenken, wie wehe es ihm im Herzen gethan habe, und doch hat müssen stillschweigend leiden. Das ist nun abermal uns zum Trost und Exempel der Geduld geschrieben, wenn wir sein Wort und Verheißung haben; [er] läßt es uns aber anders gehen, daß wir nicht in Zweifel fallen, sondern wissen, daß uns GOtt nicht verlasse; wie er am Worte gehalten hat, daß es endlich ist erfüllt worden.

4. Siehe nun die Geschichte an, woher es kommt: „Sie ging heraus (spricht der Text), die Töchter oder Mägde des Landes zu besuchen.“ Das steht einer Jungfrau übel an, daß sie hin und wieder laufe; [es] ist nicht ein gut Zeichen, wenn sie nicht können daheim bleiben. Jakob hätte es ihr auch nicht zugegeben, so herum zu laufen. Aber es geht also, wenn GOtt nicht hütet, so hütet niemand, wie der 127. Psalm, V. 1., sagt: „Wo der HErr nicht die Stadt bewahret, so wachen die Wächter umsonst.“ Er hat die Tochter ohne Zweifel auf das allerbeste erzogen, aber es hilft nicht. Also widerfährt noch vielen Eltern, wenn sie die Kinder wohl aufgezogen haben, daß sie dennoch übel gerathen. Er will nicht, daß man sie lasse umherlaufen, und ihren Willen gestatte, sondern mit allem Fleiße ziehe und lehre; geräth es, so dancke man
\quellseite{272}
Gott, geräth es nicht, so haben wir das Unsere gethan. Denn, daß Kinder wohlgerathen, ist nicht in unserer, sondern GOttes Gewalt und Macht; wo er nicht mit im Schiffe ist, da fährt man nimmer wohl.

5. Weiter siehe, was das Uebel für Unglück mit sich bringt, daß Jakobs Söhne, Simeon und Levi, über die Stadt Sichem fallen, und alles, was männlich ist, ermorden, wider ihres Vaters Willen, also, daß sie sich allzumal versündigen, Simeon und Levi, Sichem, Hemor und Dina, und alle Brüder; [es] sind eitel Bubenstücke, das sie treiben; dazu die zwei Söhne auch meinleidig werden, daß sie nicht Frieden halten, wiewohl sie es gelobt hatten, wenn sich die Bürger beschnitten; also, daß die ganze Stadt des einigen Buben muß entgelten.

6. Wie aber, daß Jakob hernach am Ende des 48. Capitels, V. 22., sich rühmt, er habe das Land und Gut mit Recht eingenommen, so es doch hier viel anders sieht? Denn seine Söhne fallen je mit Gewalt und Unrecht über die Stadt, schlagen und nehmen es alles ein. Aber es geht also zu, wie GOtt pflegt aus Unrecht Recht zu machen. Die Kinder thun Unrecht wider den Vater und ihre Zuziehung, doch kriegt es Jakob mit Recht zu sich; also, daß Sichem und Hemor billig und recht gestraft werden, wiewohl jene Unrecht thun. Wie GOtt gemeiniglich Buben durch Buben strafen läßt, zuweilen auch fromme Leute durch Buben, und wiederum; als der heilige Hiob durch böse Buben aus Arabien gestraft ward [Hiob 1, 15.], dagegen der fromme Abraham die vier Könige schlug, wie droben [1 Mos. 14, 15.] gehört ist. Summa: Es ist alles in GOttes Hand, es komme her, wie und wo es wolle, wenn er es nun gibt, der hat es. Daß nun die Sicheniter geschlagen werden, und um das Ihre kommen, ist auch recht, wiewohl Jakobs Söhne mit Unrecht hinein fallen; doch nimmt es GOtt, weil sie geschlagen sind, und gibt es dem Jakob, daß er es mit gutem Rechte behält.

7. Das sind seine Werke, so er einem ein Land gibt, oder nimmt, und entwendet, wohin er will; wie er die Hebräer den Egyptern ließ entwenden ihre güldenen und silbernen Geräthe, wie im andern Buche Mosis [Cap. 12, 36.] steht. Das war vor den Leuten unrecht gehandelt, vor GOtt aber nicht, denn er wollte jene also strafen, wie sie verdient hatten. Also geht es jetzt in der Welt zu, daß ein Fürst den andern, ein Land das andere beraubt oder schlägt, wie ein Bube den andern.

8. Doch müssen wir es dafür halten, daß die zwölf Söhne des Patriarchen Jakob auch fromm und heilig sind gewesen, weil sie die Schrift oft hoch hebt und preist. Doch haben sie oft gar weidlich gestrauchelt, und manche grobe Stücke begangen, wie wir hernach mehr sehen werden; das ja viel ist von den großen Patriarchen. Jetzt wäre es ein Bubenstück, so falsch und mit Untreue zu handeln, wie sie hier thun, und mit Hinterlist über die wehrlosen Sicheniter fallen, und alles schlagen und wegnehmen, so sie doch zu den Bürgern keine Schuld hatten. Dazu sich es nicht gebührt, zu empören wider Hemor und Sichem, ob sie schon mit Unrecht und Gewalt führen, weil sie Herren des Landes waren, Jakob aber mit den Seinen als Fremdlinge darinne wohnten. Also läßt GOtt die Seinen vorher zu großen Sündern werden, ehe er sie fromm macht. So mußte St. Paulus auf das erste die Christen auf das ärgste verfolgen, ehe er der vornehmste Apostel ward [Apost. 8, 3. 9, 1. 2. 1 Tim. 1, 13. 14.].

9. Das ist alles uns zum Exempel geschrieben, daß, ob ein Mensch gleich strauchelt mit groben Werken, [er] nicht darum zu verdammen sei. Sollten die Werke vor GOtt fromm machen, so müßte man diese Patriarchen gar verwerfen; denn man findet in der Legende schier nichts von ihnen, denn solche groben Stücke. Darum muß vor GOtt nichts gelten, denn lauter Gnade und Schonen. Derhalben viel besser ist, daß man also von den Heiligen predige, denn wie sie gefastet, diese oder jene Werke gethan haben. Denn das sind die heimlichen und wunderbaren Gottes Werke, daß er will Sünder zu Heiligen haben, auf daß alle unser Ruhm von Frömmigkeit und guten Werken niedergeschlagen werde. Denn das ist die Summa davon: So lange er uns hält, so lange sind wir fromm; wenn er aber die Hand abzeucht, und läßt uns selbst hingehen, so sind wir verzweifelte Buben. Doch soll niemand verzweifeln, ob wir fallen, allein, daß wir das Wort nicht fahren lassen; denn sein Wort und Gnade ist größer und mehr denn aller Menschen Sünde. Das sei kürzlich dies Capitel nach der Historie; wollen auch ein wenig von der heimlichen Deutung sehen.
\quellseite{273}

10. Daß Dina, Jakobs Tochter, herausgehet unter die fremden Weiber, bedeutet die leichtfertigen Seelen, die nicht daheim in GOttes Worte bleiben, sondern sich herausgeben und mengen unter Menschenlehre, und fallen auf das, so der Vernunft gefällt. Wenn Dina daheim bei dem Vater geblieben wäre, so wäre sie auch unverrückt geblieben; also bleibt das Gewissen auch rein, weil es allein am Worte Gottes hangt; sobald es aber etwas mehr sucht, und geht nach Menschenlehren, wird es gewißlich verderbt. Denn GOtt will, daß wir allein seine Braut seien, und an ihm hangen; so ist es nicht möglich, wo Menschengesetze in das Gewissen gerathen, daß es unverrückt bleibe. Ich sage aber nur von den Menschenlehren und Gesetzen, die man auf das Gewissen schlägt und treibt, nicht von äußerlichen Satzungen, die allein auf die Hand gelegt werden.

11. Als sie nun herausgegangen ist, wird ihrer der Landesherr Sichem gewahr, und buhlt mit ihr; das ist, wie Habakuk [Cap. 1, 16.] sagt: Cibus ejus electus, der Teufel, der Weltfürst, sucht nur niedliche Bißlein; das bedeutet, daß er durch Menschenlehrer nicht seinen Haufen, sondern nur die, die Gottes Wort gehört haben, und GOttes Volk sind, zu sich bringe und betrüge. Also beschläft er sie nun, und schmächt sie; denn der Teufel sucht nicht die Seelen zu verderben, die er [zu]vor unter ihm hat, sondern nur die noch rein sind, wie er durch das Papstthum gethan hat.

12. Daß sich aber Sichem und Hemor mit Jakob und seinen Söhnen verträgt, und ein Verbündniß macht, auf daß er sie behalte und wieder zu Ehren mache, bedeutet die Tücke und Heuchelei der Menschenlehrer, so sie brauchen, die Leute zu betrügen, und zu sich zu reißen. Als die Papisten thun und vorgeben, sie wollen es mit dem Evangelio halten, und alles predigen, was darinne steht; wie die Sichemiter bewilligen, sich zu beschneiden, und sprechen: Wir wollen nichts lehren, das Schaden bringt, sondern was zur Seligkeit dient. Also haben sie eben die Schrift und das Evangelium, so wir haben, aber sie meinen es falsch, streichen ihm nur so eine Farbe an, daß sie uns zu sich bringen und vom rechten christlichen Glauben führen.

13. Solches lassen nun ihre Brüder zum ersten geschehen, aber bald fallen sie herein mit Gewalt, nehmen die Schwester wieder, schlagen und verstören alles; das ist, die rechten Prediger sind da, und gewinnen die Seelen und Gewissen, bringen sie GOtt wiederum heim. Da sie aber die Bürger der Stadt täuschen, und stellen sich, als wollten sie bei ihnen wohnen und mit ihnen halten, bedeutet, daß wir auch in den Menschenlehren gesteckt sind; aber darnach sind wir zugefahren, haben das rechte Schwert ergriffen, das ist, GOttes Wort, damit wir niederlegen und zerstören alles, was Menschentand ist, damit sie die Gewissen verderben, hören nicht auf, bis wir es gar gedämpft und alle gemacht haben. Und daß die zwei Brüder allein hingehen und alles ermürgen, ist, daß derer wenig sind, die GOttes Wort predigen, aber der andern falschen Heiligen sind der größte Haufe; doch GOttes Wort so mächtig ist, daß es ihnen den Schaden thut, und ihr Ding zunichte macht.

\section{Das fünfunddreißigste Capitel}
\begin{schriftzitat}
V. 1—8. Und GOtt sprach zu Jakob: Mache dich auf, und zeuch gen Bethel, und wohne daselbst, und mache daselbst einen Altar dem GOtt, der dir erschien, da du flohest vor deinem Bruder Esau. Da sprach Jakob zu seinem Hause und zu allen, die mit ihm waren: Thut von euch die fremden Götter, so unter euch sind, und reiniget euch, und ändert eure Kleider, und lasset uns auf sein, und gen Bethel ziehen, daß ich daselbst einen Altar mache dem GOtt, der mich erhöret hat, zur Zeit meines Trübsals, und ist mit mir gewesen auf dem Wege, den ich gezogen bin. Da gaben sie ihm alle fremden Götter, die unter ihren Händen waren, und ihre Ohrenringe; und er vergrub sie unter eine Eiche, die neben Sichem stund, und sie zogen aus. Und es kam die Furcht GOttes über die Städte, die um sie her lagen, daß sie den Söhnen Jakobs nicht nachjagten. Also kam Jakob
\quellseite{274}
gen Lus, im Lande Canaan, die da Bethel heißt, sammt alle dem Volke, das mit ihm war. Und baute daselbst einen Altar, und hieß die Stätte El-Bethel, darum, daß ihm daselbst GOtt offenbaret war, da er floh vor seiner Bruder. Da starb Debora, der Rebekka Amme, und ward begraben unter Bethel, unter der Eiche, und ward genannt die Klageiche.
\end{schriftzitat}

1. In dieser Gegend hören wir nichts denn Jammer und Herzeleid, ohne daß zuweilen gesagt wird, wie er einen Altar bauet und betet. Aber, wie ich mehr gesagt habe, das Hauptstück darinne ist, daß er GOttes Wort oft gehört hat. Wo das ist, da gelten die Werke auch, wie geringe sie sind. So erzählt er hier, wie GOtt ihn heißt einen Altar bauen, um der Ursache willen, daß ihn GOtt erlöst, daß seine Kinder nicht erschlagen waren. Denn sie hatten es also ausgerichtet, daß das ganze Land ihnen feind ward, wie er auch klagt im vorigen Capitel, V. 30.: „Ihr habt es zugerichtet, daß ich stinke vor den Einwohnern dieses Landes; wenn sie sich nun versammeln über mich, so werden sie mich schlagen, daß ich sammt meinem Hause vertilgt werde.“ Darum ist er nicht mehr im Lande sicher, daß ihn GOtt muß heißen hinwegziehen.

2. Also wirft GOtt die Seinen immerdar in Noth und allerlei Fahr, läßt sie aber nicht darinne verderben. Er hatte die Verheißung, daß sein Same sollte groß werden, und das Land überkommen; und hier fehlt es nicht weit, daß sie alle erschlagen wären; und [es] wäre auch geschehen, wo nicht GOtt eine Furcht hätte lassen kommen, wie der Text sagt, über die Städte, die umher lagen, daß sie den Söhnen Jakobs nicht nachjagten; also, daß er sie gleich als in einem Augenblicke errettet und [ihnen] davon hilft. Das sind Jakobs Wehre, damit er sich schützt, und die eiserne Mauer, die GOtt um ihn her bauet. Denn er hat gar bald geholfen; es kostet nur ein Winken mit dem Auge, so hat er die ganze Welt erschlagen. Wenn er schlagen will, so hebt er inwendig an, und erschlägt von ersten das Herz; wenn das genommen ist, so ist seine Kraft mehr da.

3. Nun, auf dies Wort, daß ihn GOtt heißt ziehen und einen Altar bauen, machte sich Jakob auf in fröhlichem Glauben, und fürchtete sich nicht, mitten durch die Feinde zu ziehen. Denn da ist angezeigt, daß Jakob abermal in großen Furchten gestanden ist, daß er nicht wußte, wo er hinaus sollte, weil er mitten unter den Feinden war, und doch an dem Orte zu bleiben nicht bleiben durfte, daß GOtt selbst muß kommen, und ihn durch sein Wort stärken und trösten, daß er sich nicht mehr fürchtet.

4. Da thut er nun eine Predigt, und sagt, was ihn Gott befohlen hat, welches das rechte Amt eines Priesters ist. Und hier siehst du, daß er nicht allein seinem Gesinde gepredigt hat, sondern auch andern, so ihm anhingen um des Wortes willen; wie hernach immerdar etliche Heiden sich zu den Juden gehalten haben, und GOttes Wort von ihnen gelernt. Denn das ist es, das der Text sagt: „Da sprach Jakob zu seinem Hause, und zu allen, die mit ihm waren“; also, daß Jakob zur selben Zeit der oberste Rabbi oder Bischof ist gewesen, und fast allein das Wort GOttes erhalten, und das geistliche Regiment geführt hat.

5. Was hat er aber gepredigt? Das, wie im Texte steht: „Thut von euch die fremden Götter“ etc., wir wollen einen Altar bauen, GOtt loben und danken, der uns geholfen und erlöst hat. GOtt einen Altar bauen, ist nicht, GOtt mit unsern Werken finden oder ergreifen, sondern [es] ist allein ein Raum gewesen, da man zusammen käme zu predigen und GOtt [zu] loben; wie auch GOttes Haus nicht heißt, daß wir ihm ein Haus bauen, sondern darum, wie oben [Cap. 28, § 12 f.] gesagt, daß er zuvor da wohnt, und wir da zu ihm kommen. Wie wohnt er denn da? Also, daß er da redet. Denn wo sein Wort gepredigt wird, da will er sich gewißlich finden lassen; wo aber dasselbige nicht ist, da ist sein Haus nicht, wenn man eitel Kirchen auf einander bauete.

6. Also bauete Jakob nicht aus seiner Andacht einen Altar, sondern aus GOttes Befehl. Also thue du auch. Denn es ist nicht in unserer Gewalt, seinen Namen zu ehren, oder ihm einen Dienst aufzurichten, sondern er muß es vorhin sagen und heißen, hier oder da soll man predigen, und mir dienen; wie er 2 Mos. 20, 24. sagt: „An welchem Orte ich meines Namens Gedächtniß mache, da will ich zu dir kommen, und dich segnen“, das ist, wo ich mein Evangelium von meiner Gnade und Güte lasse predigen, meinen Namen zu preisen, nicht wo du Kirchen bauest oder weihest. Denn wer bist
\quellseite{275}
du, daß du dich unterstehest, GOtt an einen Ort zu binden? Er will zuvor kommen, aus lauter Gnade und Güte, und heißen aufrichten. Darum ist der bloße Altar viel köstlicher, denn wenn wir von eitel Golde Kirchen baueten, und alle Heiligen dahin brächten; denn hier ist GOttes Wort, welches mehr gilt, denn alle Welt.

7. Was ist es aber, daß Jakob heißt die fremden Götter von sich thun, und sich reinigen, und andere Kleider anziehen? Was die fremden Götter sind, habe ich droben [Cap. 31, § 23 ff.] genug gesagt. Jakob aber, wiewohl er heilig war, und voll Geistes, doch hat er es nicht können dahin bringen, daß keine Götzen unter seinem Gesinde wären, wie wenig ihrer waren. Das war nun nichts anders, denn daß etliche mit falschem Gottesdienste GOtt haben wollen dienen, wie wir durch unsere selbsterdachten Werke. Denn er will keinen Gottesdienst haben denn den, so er selbst befohlen und geboten hat. So hat er sie nun geheißen, alle solche falsche Weise, GOtt zu dienen, fahren [zu] lassen und ab[zu]legen; denn das wäre nichts gewesen, wenn sie gleich die Götzen von sich gethan hätten, und das Herz gleichwohl abgöttisch geblieben wäre.

8. Die andern zwei Stücke gehören auf äußerliche Weise und Gebärden. Möchtest du aber sprechen: Siehst du nicht, daß hier geboten wird, äußerlich reinigen und waschen, und andere Kleider anziehen? darum muß man je die Ceremonien und unsern Gottesdienst auch gelten lassen. Es gilt nicht, daß man sagt, man müsse es geistlich verstehen, und auf die Deutung sehen, weil es keine Noth zwingt, sondern muß es lassen äußerlich und leiblich bleiben. Antwort: Wir wollen es auch lassen stehen, sagen aber also: Wiewohl die Seligkeit und der Glaube oder christlich Wesen nicht steht in solchen äußerlichen Dingen, so kann es doch ohne solches nicht geschehen, gleichwie Paulus 1 Cor. 8, 8. sagt: „Die Speise fördert uns nicht vor GOtt; essen wir, so sind wir darum nicht besser; essen wir nicht, so sind wir darum nicht weniger.“ Doch müssen wir essen und trinken, das Leben zu erhalten, auf daß wir können predigen und dem Nächsten dienen.

9. Also mögen wir auch sagen: Kein äußerlich Ding fördert oder hilft uns vor GOtt; dennoch müssen wir äußerlich Ding und Geberde halten, so dazu dienen, daß man die Leute zum Worte GOttes halte; als, daß man einen Ort wählt, da man GOttes Wort predigt und die Sacramente handelt etc. Darum soll man solches alles lassen frei bleiben. Also ist diese Kirche nichts besser denn ein ander Haus; doch schickt es sich nicht, daß man hier wollte fressen, saufen, oder tanzen, und dergleichen; denn es muß je eine Ordnung und Unterschied sein unter den Leuten, wo man dies oder jenes thue und treibe, welches vor GOtt weder fördert noch hindert.

10. Darum ist es eben so viel, daß Jakob sagt: „Reiniget euch, und ändert eure Kleider“, als wenn ich spräche: Wenn ihr wollt Predigt hören, so lasset eure Leichtfertigkeit, Fressen und Saufen daheim, und stellet euch tapfer. Also will er auch sagen: Das Hauptstück, das wir thun, ist, daß wir hinziehen, bauen einen Altar, predigen und GOtt anrufen; dazu ist nun gut, daß ihr die Götzen hinwegwerfet, und feine, ehrliche Geberden führet.

11. Darnach sagt er nun, wie sie die Götzen von sich gegeben haben, und gedenkt der andern zwei Stücke nicht mehr. [Es] ist auch zu wissen, daß die Kleider, so die Schrift heißt ablegen, sind die gewesen, damit sie sich geschmückt haben den Götzen zu Ehren, als, die Ehrenringe, die sie hier hinweggegeben, und Jakob mit den Götzen vergräbt unter einen Baum. Denn es ist in den Landen eine gemeine Weise und Schmuck gewesen, Ehrenringe und Armspangen zu tragen, und solchen Schmuck ihren Götzen zu Ehren anzulegen. Was es bedeutete, werden wir hernach hören. Nach der Historie ist es darum zu thun gewesen, daß sie kein Zeichen an sich hätten, damit sie anzeigten, daß sie irgend an einem Götzen hingen, als ieztund die Pfaffen, Mönche und Nonnen an ihren Kleidern, Kappen und Platten hangen.

12. Also kommt er mit allem, was er bei sich hat, in das Land Canaan, gen Bethel, nicht weit von Jerusalem, da ihm zuvor auch GOtt erschienen war, und bauet allda einen Altar, daher die Stätte hernach, wie droben [Cap. 28, § 17] gesagt, eine Ursache ist worden der größten Abgötterei unter den Juden. Denn es hat der falsche Gottesdienst nie aufgehört, so lange bis Israel hinweg geführt ward. Das machte allein, daß sie auf dies Exempel fielen, und sagten: Gott wohnte sonderlich an dem Orte, weil er da zweimal dem Jakob erschienen war, und hatte ihm geheißen, einen Altar bauen. Wenn der Papst jetzt so starke Exempel und Schrift
\quellseite{276}
für sich hätte, wer wollte vor ihm bleiben? Als er nun dahin kommt, starb der Rebekka Amme; da erhebt sich das Trauern und Klagen. Er ist nun aus der Feinde Händen erlöst, aber nun folgt sein eigen Unglück, wie wir sehen werden.

\begin{schriftzitat}
V. 9—29. Und GOtt erschien Jakob abermal, nachdem er aus Mesopotamia gekommen war, und segnetet ihn, und sprach zu ihm: Du heißest Jakob, aber du sollst nicht mehr Jakob heißen, sondern Israel sollst du heißen; und also heißt man ihn Israel. Und GOtt sprach zu ihm: Ich bin der allmächtige GOtt, sei fruchtbar und mehre dich, Völker und Völker-Haufen sollen von dir kommen, und Könige sollen aus deinen Lenden kommen. Und das Land, das ich Abraham und Isaak gegeben habe, will ich dir geben, und will es deinem Samen nach dir geben. Also fuhr GOtt auf von ihm, von dem Orte, da er mit ihm geredet hatte. Jakob aber richtete ein steinern Mal auf an dem Orte, da er mit ihm geredet hatte, und goß Trankopfer\footnote{Randglosse: Trankopfer. Das war Wein, wie das in den folgenden Büchern genugsam gesehen wird.} darauf, und bezog ihn mit Oele. Und Jakob hieß den Ort, da GOtt mit ihm geredet hatte, Bethel. Und sie zogen von Bethel, und da noch ein Feldweges war von Ephrath, da gebar Rahel, und es kam sie hart an über der Geburt. Da es ihr aber so sauer ward in der Geburt, sprach die Wehmutter zu ihr: Fürchte dich nicht, denn diesen Sohn wirst du auch haben. Da ihr aber die Seele ausging, daß sie sterben mußte, hieß sie ihn Benoni\footnote{Randglosse: Ben Oni heißt meines Schmerzens Sohn.}. Aber sein Vater hieß ihn Benjamin\footnote{Randglosse: Ben Jamin heißt der Rechten Sohn.}. Also starb Rahel, und ward begraben an dem Wege gen Ephrath, die nun heißt Bethlehem. Und Jakob richtete ein Mal auf über ihrem Grabe, dasselbe ist das Grabmal Rahel, bis auf diesen Tag. Und Israel zog aus, und richtete eine Hütte auf jenseit dem Thurm Eder. Und es begab sich, daß Israel im Lande wohnte, ging Ruben hin, und schlief bei Bilha, seines Vaters Kebsweibe. Und das kam vor Israel. Es hatte aber Jakob zwölf Söhne. Die Söhne Lea waren diese: Ruben, der erstgeborne Sohn Jakobs, Simeon, Levi, Juda, Jisachar und Sebulon. Die Söhne Rahel waren: Joseph und Benjamin. Die Söhne Bilha, Rahels Magd: Dan und Naphtali. Die Söhne Silpa, Lea Magd: Gad und Asser. Das sind die Söhne Jakobs, die ihm geboren sind in Mesopotamia. Und Jakob kam zu seinem Vater Isaak gen Mamre in die Hauptstadt, die da heißt Hebron, da Abraham und Isaak Fremdlinge innen gewesen sind. Und Isaak ward hundert und achtzig Jahre alt, und ward krank und starb, und ward versammelt zu seinem Volke, alt und des Lebens satt. Und seine Söhne, Esau und Jakob, begruben ihn.
\end{schriftzitat}




13. Da ist aber mal eine andere Predigt, die GOtt mit ihm hält, wie er zuvor auch mit ihm geredet hatte an demselben Orte, als er auszog. Warum sagt er aber also: „Du sollst nicht mehr Jakob heißen, sondern Israel sollst du heißen“? Hatte er ihm doch den Namen nun eben lange zuvor gegeben, als er mit dem Engel kämpfte. Es war aber der Name noch nicht ausgekommen, hier aber wird er ausgeschrien und lautbar, daß er mit alle seinem Geschlechte Israel ward\footnote{Walch und die Erlanger: wird.} geheißen. Doch hat man ihn immerdar nicht allein Israel, sondern auch Jakob geheißen; wie man in den Propheten sieht. Nun thut er ihm eine neue reiche Verheißung, daß von dem Jakob, der so gering und elend war vor der Welt, noch sollten mächtige Könige kommen. Es ward aber langsam hernach erfüllt, denn sein Same lange Zeit mußte die geringsten Knechte sein im Lande Egypten.

14. Als er nun fortzeucht, widerfährt ihm aber ein Herzeleid, daß Rahel über der Geburt stirbt, seine liebste Frau, um welcher willen er so viele Jahre gedient hatte, daß er sie nicht kann heimbringen, sondern stirbt ihm auf dem Wege. Das ist ihm auch ein hartes Kreuz gewesen. [Er] hat großen Jammer an ihr müssen sehen, wie der Text anzeigt, daß es ihr sehr sauer ist worden über dem Kinde, daß es gar genau lebendig ist von ihr gekommen; darum gibt sie dem Kinde noch einen Namen von dem Schmerzen, ehe sie stirbt, und heißt ihn Ben Oni, ein Sohn meines Schmerzens. Als sollte sie sagen: Der soll das Zeichen tragen, daß ich ihn mit Wehe und Schmerzen geboren habe, daß es mich mein Leben kostet. Der Vater gibt ihm auch einen Namen, und heißt ihn Benjamin, einen Sohn der rechten Hand, das ist, seinen allerliebsten Sohn, weil er nun in seinem Alter geboren wird, und an der Mutter Statt tritt.
\quellseite{277}

15. Ehe er nun dieses Leides vergißt, kommt ihm aber ein ander Unglück, daß sein erster Sohn Ruben seine Stiefmutter beschläft. Ich meine ja, das sei ein Stück von des heiligen Mannes Söhne; wie könnte er größere Sünde und Schande anrichten, und dem Vater mehr zu Leide thun? St. Paulus hat das Laster groß angezogen [1 Cor. 5, 1.]: „Eine solche Hurerei, da auch die Heiden nicht wissen von zu sagen, daß einer seines Vaters Weib habe“,\footnote{Erlanger: hat.} hat sich aber gemäßigt, daß er nicht durfte sagen von den Juden. Es ist auch diesem Ruben hernach nicht wohl gelungen, daß er am Ende ein Hartes vom Vater hören mußte, und beide, des Priesterthums und Regimentes, so dem ersten Sohne gebührte, beraubt sein.

16. Da läßt uns GOtt sehen, was wir für fromme Kinder sind, wenn er uns nicht ohne Unterlaß in seiner Hand hält, daß uns keine Büberei noch Untugend zu viel ist. Aber der fromme Jakob muß es alles in sich fressen und leiden, daß ihm seine eigenen Kinder mehr Leides thun, denn die Fremden und seine Feinde. Denke du selbst, wie wehe es ihm gethan habe, daß ihm weder Tochter, noch sein eigen Weib unerrückt bleibt. Daraus sollen wir sehen, was ein recht christlich Wesen ist, daß immer GOtt ein Kreuz über das andere den Seinen zuschickt, auf daß sie geschickt und bereit seien zu leiden, daß der alte Adam wohl gedämpft, und des Lebens müde werde.

17. Aufs letzte werden die zwölf Söhne Jakobs noch einmal erzählt, um Benjamin willen, welcher zuletzt geboren war; darnach schreibt er, wie Jakob wieder zum Vater Isaak kam, daß er ihn noch sieht vor seinem Tode, und begräbt ihn selbst mit seinem Bruder. Aber von Rebekka schreibt er nichts mehr; [sie] wird vielleicht vorhin gestorben sein, denn er war nun drei und zwanzig Jahre außen gewesen, und fast hundert Jahre alt. So haben wir den Isaak dahin; ist noch ein wenig übrig von Jakobs Legende.

18. Nun lasset uns auch die heimliche Deutung dieses Capitels durchlaufen. Zum ersten zeucht Jakob gen BethEl, das heißt GOttes Haus, oder da GOtt wohnt; ehe sie dahin kommen, predigt er ihnen vor, wie sie sollen die Götzen von sich thun, die nimmt er sammt den Ohrringen und Spangen, und vergräbt sie unter eine Eiche. In dieser Geschichte ist lieblich bedeutet die Predigt des Evangelii, welche nirgend geschieht, denn da GOtt wohnt. Dasselbige Evangelium verwirft untereinander die silbernen Götzen, das ist, die Menschenlehre, oder solche Gottesdienste, so Gott nicht befohlen hat, welche nichts denn eine lautere Abgötterei sind, wie wir gehört haben. Denn er will keinen andern Dienst leiden, der vor ihm gelten soll, denn welchen er uns selbst lehrt und heißt, auf daß unser Herz gewiß sei, daß es ihm gefalle. Solchen Gottesdienst können die nicht thun, die sein Wort nicht haben, sondern was sie thun, das thun sie aus ihrem Dünkel und Traume. Das sind eigentlich Bilder mit Händen gemacht, das ist, eigene Werke und erdichtete Menschenlehre, welche alle das Evangelium verdammt und nicht leiden will.

19. Welche nun das Evangelium hören und annehmen, die werfen solche Götzen von sich, dazu auch ihre Ohrringe, das ist, wenn sie hören, daß ihre Werke und Gottesdienste nichts gelten, so helfen sie selbst die Abgötterei verdammen. Die Ohrringe, welche man in den Ländern gemeiniglich pflegte zu tragen, bedeuten nichts denn die Predigten, die sie zuvor in Ohren getragen haben, und sich damit gerüstet, die reißen sie nun hinweg, und begraben sie unter die Erde, daß sie nichts mehr schaffen noch schaden können.

20. Als nun Jakob gen BethEl kommt, richtet er ein steinern Mal auf, da GOtt mit ihm geredet hatte, und gießt Trankopfer und Oel darauf. Ein köstlich Gebäu ist das gewesen, da weder Dach noch Gewölbe oder Haus war, denn ein schlechter Stein, da Jakob gelegen hatte; noch hat es GOtt besser gefallen, denn unsere allerherrlichsten Tempel und Kirchen, darum, daß er es Jakob befohlen hatte. Der Stein nun, wie auch droben gesagt ist am 28. Capitel [§ 38], bedeutet Christum, darauf wir uns legen, und sicher auf ihm ruhen und schlafen. Der wird nun aufgerichtet zu BethEl, das ist, wo man GOttes Wort predigt, da richtet man den HErrn Christum auf zum Grundsteine. Darum ist das Aufrichten nichts anders, denn das Evangelium predigen, welches allein auf den Christum weist; derhalben auch nichts Anderes soll gepredigt werden. Daraus schließe nun, wo Christus recht gepredigt wird, daß da gewiß ein Gotteshaus sei.
\quellseite{278}

21. Was bedeutet aber das Trankopfer? In folgenden Büchern Mosis, nämlich im dritten, wird oft von Trankopfern oder Gießopfern gesagt, da die Priester Semmelmehl oder Kuchen nahmen, und gossen Wein darauf, thaten gemeiniglich auch Oel dazu [2 Mos. 29, 40. 3 Mos. 23, 13. 18. 37.]. Wie geht nun das im Evangelio zu? Was in der Schrift von Opfern steht, können wir nicht das deuten, denn auf das Predigtamt in der Christenheit, und nehmen den Grund aus St. Paulo, Röm. 15, 15. f.: „Ich habe euch (spricht er) ein wenig dürftiglich geschrieben, euch zu erinnern, um der Gnade willen, die mir von GOtt gegeben ist, daß ich soll sein ein Diener Christi unter die Heiden, zu opfern das Evangelium GOttes, auf daß die Heiden ein Opfer werden, GOtt angenehm, geheiliget durch den Heiligen Geist.“ Wenn ich nun so lehre, daß der alte Adam in Christo soll getödtet werden, und untergehen mit allem, was von ihm kommt, so gieße ich Wein auf den Stein, das ist, ich bekenne, daß unser Ding nichts gilt, sondern müssen alle nach dem Fleische getödtet werden; daß, wie der Wein die Leute trunken macht, also macht das Evangelium, wenn man es recht versteht und faßt, daß man nicht achtet Gut, Ehre und Leben, läßt ihm alles nehmen und sein Ding verdammen; die hält denn die Welt für Narren.

22. Das ist nun eine Predigt des Evangelii, daß alles, was wir sind, verdammlich ist, wie Christus in Johanne [Cap. 16, 8.] sagt: „Der Heilige Geist wird die Welt strafen um die Sünde, Gerechtigkeit und Gericht“; die das nun glauben, die werden begossen mit dem scharfen Weine. Darnach folgt das andere, daß man auch Oel darauf gießt. Das ist die andere Predigt von der Barmherzigkeit und Gnade GOttes, damit wir getröstet werden. Wo das nun geschieht, da ist ein recht Bethel, da GOtt wahrhaftig wohnt. Das ist die Summa des ganzen Predigtamtes, das die zwei Stücke führt. Denn es kann kein Mensch ein gutes Gewissen überkommen, daß GOtt in seinem Herzen wohne, es lasse denn zuvor den alten Adam tödten, daß er nicht fragt nach dem, das Fleisch und Blut groß achtet. Solches Gießen, das ist, diese Predigt soll nun immerdar währen, so lange wir leben und die Welt steht. Siehe, das ist uns hier in den Vätern mit lieblichen Figuren vorgespielt, und hernach durch die Propheten, allermeist aber im neuen Testamente durch die Apostel getrieben.

23. Weiter: Wenn sich das Evangelium durch Christum und die Apostel anhebt, so legt sich Rahel und stirbt, das ist, die Synagoga, oder das Judenthum. Das Kind kommt hervor, aber die Mutter muß darüber bleiben. Aus dem Sohne wird nun ein neu Volk, dadurch, daß das alte Volk ist untergegangen. Das zeigt auch an der Name, so sie ihm gibt, Ben-Oni, Schmerzenssohn, das ist, daß wir müssen Kinder des Kreuzes sein, und nur davon predigen. Der Vater aber gibt ihm einen andern Namen, und heißt ihn Benjamin, das ist, der Rechten Sohn. Denn nach dem Adam sind wir Schmerzenskinder, arme, elende Leute, aber nach dem Geiste sind wir Benjamin, das ist, Kinder des ewigen Lebens, und aller Güter GOttes. Also hat der Geist gespielt mit den Vätern, anzuzeigen, wie es zur Zeit des Evangelii gehen sollte.

24. Weil nun das Evangelium im Schwange geht, und das Judenthum umkommt, da fährt Ruben zu, und begeht einen Ehebruch mit seiner Stiefmutter. Er war der erste Sohn, der da sollte der andern aller Herr sein, wie die Freiheit\footnote{Freiheit = Vorrecht, Privilegium.} der ersten Geburt war; der verderbt es und schläft bei der Mutter; das ist, die Besten im jüdischen Volke, die Weisesten und Heiligsten haben die Synagoga verderbt und zu Schanden gemacht mit Menschenlehren, und des Vaters Bett befleckt, welches ist das Gewissen, darinne Christus ruht und bei seiner Braut liegt, und nicht leiden will, daß jemand anders darinne regiere. Wenn man nun predigt, daß unser Glaube allein an Christo hangen soll, so ist das Gewissen rein und unbesudelt; wer aber durch Werke die Gewissen will führen, und auf Menschenlehre bauen, der legt sich auf des Vaters Bette, und verderbt die Seelen. Wie die Heuchler und Schriftgelehrten zu Christi Zeiten das Volk auf ihre Werke führten, und hernach das Papstthum die Welt voll Menschenstand genommen hat, und nach ihnen aber andere aufkommen, die uns von Christo und reinem Glauben führen. Also geht diese Deutung eigentlich auf falsche Lehre, welche allein den Schaden thut in dem Gewissen.

25. Derhalben ich oft vermahnt habe, daß man weit solle von einander scheiden Leben und
\quellseite{279}
Lehre. Die Lehre ist, daß ich an Christum glaube, mein Werth, Leiden und Tod für nichts achte, und den Nächsten diene, und darüber lasse fahren, was ich soll. Das Leben aber ist, ob ich so oder so wandle und darnach thue. So liegt nun weit nicht so viel am Leben, als an der Lehre, daß, wenn schon das Leben nicht so rein ist, kann dennoch die Lehre wohl rein bleiben, und mit dem Leben Geduld getragen werden. Solches lehren unsere Kinder nicht, sondern welche die Besten unter ihnen sind, und aufs höchste kommen, die predigen nur von strengem Leben, und legen uns große Exempel der Heiligen vor, die große, wunderliche Werke gethan haben, mit Lachen in [den] Tod gegangen sind, und führen die Leute dahin, daß sie der Lehre nicht wahrnehmen. Denn ein Mensch ist nicht besser zu verführen, denn durch solch scheinend Leben; wo nicht wackere Prediger sind, ist es eine hohe Gnade, wer das Leben von der Lehre scheiden kann.

26. Wahr ist es, daß wir also sollen leben; aber ich lebe, wie ich lebe, so wird darum die Lehre nicht falsch. Darum müssen wir nicht nach dem Leben, sondern nach der Lehre sehen und richten. Aber der Haufe hält allezeit mehr vom Leben, denn von der Lehre; [sie] sagen alle also: Ei, was ist es, wenn man gleich lange lehrt vom Glauben, es muß ja höher kommen. Höher kann ich nicht predigen, denn daß man den alten Adam tödte, und ein neuer Mensch werde. Sagen sie: Ja, thust du es doch nicht? Antwort: Ich sollte es wohl thun, ja, wenn mir es GOtt auch gibt; aber so hoch wird es niemand bringen, es wird noch wohl viel daran fehlen.

27. Derhalben laß das Leben herunter bleiben auf Erden, die Lehre hebe hinauf in Himmel. Die Lehre bleibt allezeit ihr selbst gleich, daß sie will ganz rein und vollkommen sein; das Leben aber kann wohl höher kommen. Ich wollte auch, daß mein Leben höher wäre, aber es wird nichts daraus; das Leben wird die Lehre nicht erreichen, so lange wir hier leben. Also sehen wir auch hier in allen Exempeln, wie weit es den heiligen Leuten noch am Leben gefehlt hat; als, wie sich Jakob gefürchtet hat; item, wie schwach seine Weiber gewesen sind und zuweilen auch gestrauchelt haben, die doch der Patriarchen Mütter sind gewesen. Siehe im Evangelio, wie manchmal die Apostel straucheln und marren, also, daß man gar nicht viel Exempel in der Schrift liest, da nicht viel Schwachheit und Fehler am Leben mit untergelaufen sind.

\section{Das sechsunddreißigste Capitel}
36, 1—43. Das ist das Geschlecht Esau, der da heißt Edom. Esau nahm Weiber von den Töchtern Canaan, Ada, die Tochter Elon, des Hethiters; und Aholibama, die Tochter des Ana, die Nefte Zibeons, des Heviters; und Basemath, Ismaels Tochter, Nebajoths Schwester. Und Ada gebar dem Esau Eliphas. Aber Basemath gebar Reguel. Aholibama gebar Jeus, Jaelam und Korah. Das sind Esau’s Kinder, die ihm geboren sind im Lande Canaan. Und Esau nahm seine Weiber, Söhne und Töchter, und alle Seelen seines Hauses, seine Habe, und alles Vieh mit allen Gütern, so er im Lande Canaan erworben hatte, und zog in ein Land von seinem Bruder Jakob. Denn ihre Habe war zu\footnote{Jenaer: so.} groß, daß sie nicht konnten bei einander wohnen, und das Land, darinnen sie Fremdlinge waren, mochte sie nicht ertragen vor ihren Gütern. Also wohnte Esau auf dem Gebirge Seir. Und (Esau ist der Edom. Dies ist das Geschlecht Esau, von dem die Edomiter herkommen auf dem Gebirge Seir, und so heißen die Kinder Esau: Eliphas, der Sohn Ada, Esau’s Weib. Reguel, der Sohn Basemath, Esau’s Weib. Eliphas Söhne aber waren diese: Theman, Omar, Zepho, Gatham und Kenas. Und Thimna war ein Kebsweib Eliphas, Esau’s Sohn, die gebar ihm Amalek. Das sind die Kinder von Ada, Esau’s Weibe. Die Kinder aber Reguel sind diese: Nahath, Serah, Samma, Missa. Das sind die Kinder von Basemath, Esau’s Weibe. Die Kinder aber von Aholibama, Esau’s Weibe, der Tochter des Ana, der Nefte Zibeons, sind diese,
\quellseite{280}
die sie dem Esau gebar: Jeus, Jaelam und Korah. Das sind die Fürsten unter den Kindern Esau: Die Kinder Eliphas, des ersten Sohnes Esau, waren diese: Der Fürst Theman, der Fürst Omar, der Fürst Zepho, der Fürst Kenaz, der Fürst Korah, der Fürst Gaetham, der Fürst Amalek. Das sind die Fürsten von Eliphas im Lande Edom, und sind Kinder von der Ada. Und das sind die Kinder Reguel, Esau’s Sohn: Der Fürst Nahath, der Fürst Serah, der Fürst Samma, der Fürst Misa. Das sind die Fürsten von Reguel im Lande der Edomiter und sind Kinder von der Basemath, Esau’s Weibe. Das sind die Kinder Ahalibama, Esau’s Weib: Der Fürst Jeus, der Fürst Jaelam, der Fürst Korah. Das sind die Fürsten von Ahalibama, der Tochter des Ana, Esau’s Weib. Das sind Esau’s Kinder und ihre Fürsten. Er ist der Edom. Die Kinder aber von Seir, des Horiten, der im Lande wohnete, sind diese: Lothan, Sobal, Zibeon, Ana, Dison, Ezer und Disan. Das sind die Fürsten der Horiten, Kinder des Seir im Lande Edom. Aber des Lothans Kinder waren diese: Hori und Heman. Und Lothans Schwester hieß Thimna. Die Kinder von Sobal waren diese: Alwan, Manahath, Ebal, Sepho und Onam. Die Kinder von Zibeon waren: Aja und Ana. Das ist der Ana, der in der Wüste Maulthiere erfand, da er seines Vaters Zibeon Esel hütete. Die Kinder aber Ana waren: Dison und Ahalibama; das ist die Tochter Ana. Die Kinder Dison waren: Hemdan, Esban, Jethran und Charan. Die Kinder Ezer waren: Bilhan, Saawan und Akan. Die Kinder Disan waren: Uz und Aran. Dies sind die Fürsten der Horiten: Der Fürst Lothan, der Fürst Sobal, der Fürst Zibeon, der Fürst Ana, der Fürst Dison, der Fürst Ezer, der Fürst Disan. Das sind die Fürsten der Horiten, die regieret haben im Lande Seir. Die Könige aber, die im Lande Edom regieret haben, ehe denn die Kinder Israel Könige hatten, sind diese: Bela war König in Edom, ein Sohn Beor, und seine Stadt hieß Dinhaba. Und da Bela starb, ward König an seine Statt Jobab, ein Sohn Serah von Bazra. Da Jobab starb, ward an seine Statt König Husam, aus der Themaniter Lande. Da Husam starb, ward König an seine Statt Hadad, ein Sohn Bedad, der die Midianiter schlug auf der Moabiter Felde, und seine Stadt hieß Avith. Da Hadad starb, regierete Samla von Masrek. Da Samla starb, ward Saul König von Rehoboth am Wasser. Da Saul starb, ward an seine Statt König Baal Hanan, der Sohn Achbor. Da Baal Hanan, Achbors Sohn, starb, ward an seine Statt König Hadar, und seine Stadt hieß Pagu, und sein Weib hieß Mehetabeel, eine Tochter Matreds, und die Tochter Mesahab. Also heißen die Fürsten von Esau in ihren Geschlechtern, Oertern und Namen: Der Fürst Thimna, der Fürst Alwa, der Fürst Jetheth, der Fürst Ahalibama, der Fürst Ela, der Fürst Pinon, der Fürst Kenaz, der Fürst Theman, der Fürst Mibzar, der Fürst Magdiel, der Fürst Iram. Das sind die Fürsten in Edom, wie sie gewohnet haben in ihrem Erblande. Und Esau ist der Vater der Edomiter.

Dies Capitel mag ein jeglicher bei sich selbst ansehen und studiren, denn es sind eitel Namen darinnen; darum wollen wir weiter fahren.

\section{Das siebenunddreißigste Capitel}
\begin{schriftzitat}
V. 1—36. Jakob aber wohnete im Lande, da sein Vater ein Fremdling innen war, nämlich im Lande Canaan. Und das sind die Geschlechter Jacobs: Joseph war siebenzehn Jahr alt, da er ein Hirte des Viehes ward mit seinen Brüdern. Und der Knabe war bei den Kindern Bilha und Silpa, seines Vaters Weibern, und brachte vor ihren Vater, wo ein böses Geschrei wider sie war. Israel aber hatte Joseph lieber, denn alle seine Kinder, darum, daß er ihn im Alter gezeugt hatte. Und machte ihm einen bunten Rock.\footnote{Randglosse: Der bunte Rock Josephs war von mancherlei Farben Faden gewebt, und bedeutet die mancherlei Gnade und Gaben des einigen Geistes in Christo und seinen Christen.} Da nun seine Brüder sahen, daß ihn ihr Vater lieber hatte, denn alle seine Brüder, waren sie ihm feind, und konnten ihm kein freundlich Wort zusprechen.
\quellseite{281}
Dazu hatte Joseph einmal einen Traum, und sagte seinen Brüdern davon. Da wurden sie ihm noch feinder. Denn er sprach zu ihnen: Höret, Lieber, was mir doch geträumet hat. Mich däuchte, wir banden Garben auf dem Felde, und meine Garbe richtete sich auf und stand, und eure Garben umher bücketen sich gegen meiner Garbe. Da sprachen seine Brüder zu ihm: Solltest du unser König werden, und über uns herrschen? Und wurden ihm noch feinder, um seines Traumes und seiner Rede willen. Und er hatte noch einen andern Traum, den erzählte er seinen Brüdern, und sprach: Sehet, ich habe noch einen Traum gehabt: Mich däuchte, die Sonne und der Mond, und elf Sterne bücketen sich vor mir. Und da das seinem Vater und seinen Brüdern gesagt ward, strafte ihn sein Vater, und sprach zu ihm: Was ist das für ein Traum, der dir geträumet hat? Soll ich, und deine Mutter, und deine Brüder kommen, und vor dich auf die Erde fallen? Und seine Brüder neideten ihn, aber sein Vater behielt diese Worte. Da nun seine Brüder hingingen, zu weiden das Vieh ihres Vaters in Sichem, sprach Israel zu Joseph: Hüten nicht deine Brüder des Viehes in Sichem? Komm, ich will dich zu ihnen senden. Er aber sprach: Hier bin ich. Und er sprach: Gehe hin, und siehe, ob es wohl stehe um deine Brüder und um das Vieh, und sage mir es wieder, wie sich’s hält. Und er sandte ihn aus dem Thale Hebron, daß er gen Sichem ginge. Da fand ihn ein Mann, daß er irre ging auf dem Felde, der fragte ihn, und sprach: Wen suchest du? Er antwortete: Ich suche meine Brüder, Lieber, sage mir an, wo sie hüten. Der Mann sprach: Sie sind von dannen gezogen; denn ich hörte, daß sie sagten: Lasset uns gen Dothan gehen. Da folgete Joseph seinen Brüdern nach, und fand sie zu Dothan. Als sie ihn nun sahen von ferne, ehe denn er nahe bei sie kam, schlugen sie an, daß sie ihn tödteten, und sprachen untereinander: Sehet, der Träumer kommt daher, so kommet nun, und lasset uns ihn erwürgen, und in eine Grube werfen, und sagen, ein böses Thier habe ihn gefressen; so wird man sehen, was seine Träume sind. Da das Ruben hörte, wollte er ihn aus ihren Händen erretten, und sprach: Lasset uns nicht eine Seele schlagen. Und weiter sprach Ruben zu ihnen: Lasset uns nicht Blut vergießen, sondern lasset uns ihn in die Grube werfen, die in der Wüste ist, und die Hand nicht an ihn legen. Er wollte ihn aber aus ihrer Hand erretten, daß er ihn seinem Vater wiederbrächte. Als nun Joseph zu seinen Brüdern kam, zogen sie ihm seinen Rock mit dem bunten Rock aus, den er hatte, und nahmen ihn, und warfen ihn in eine Grube; aber dieselbige Grube war leer, und kein Wasser darinnen; und setzten sich nieder zu essen. Indes huben sie ihre Augen auf, und sahen einen Haufen Ismaeliter kommen von Gilead mit ihren Kamelen, die trugen Würze, Balsam und Myrrhen, und zogen hinab in Egypten. Da sprach Juda zu seinen Brüdern: Was hilft es uns, daß wir unsern Bruder erwürgen, und sein Blut verbergen? Kommt, lasset uns ihn den Ismaeliten verkaufen, daß sich unsere Hände nicht an ihm vergreifen, denn er ist unser Bruder, unser Fleisch und Blut. Und sie gehorchten ihm. Und da die Midianiter, die Kaufleute, vorüber reiseten, zogen sie ihn heraus aus der Grube, und verkauften ihn den Ismaeliten um zwanzig Silberlinge, die brachten ihn in Egypten. Als nun Ruben wieder zur Grube kam, und fand Joseph nicht darinnen, zerriß er sein Kleid, und kam wieder zu seinen Brüdern, und sprach: Der Knabe ist nicht da, wo soll ich hin? Da nahmen sie Josephs Rock, und schlachteten einen Ziegenbock, und tunkten den Rock in das Blut, und schickten den bunten Rock hin, und ließen ihn ihrem Vater bringen und sagen: Diesen haben wir gefunden, siehe, ob es deines Sohnes Rock sei, oder nicht? Er kannte ihn aber, und sprach: Es ist meines Sohnes Rock, ein böses Thier hat ihn gefressen, ein reißend Thier hat Joseph zerrissen. Und Jakob zerriß seine Kleider, und legte einen Sack um seine Lenden, und trug Leid um seinen Sohn lange Zeit. Und alle seine Söhne und Töchter traten auf, daß sie ihn trösteten, aber er wollte sich nicht trösten lassen, und sprach: Ich werde mit Leid hinunter fahren in die Grube zu meinem Sohne. Und sein Vater\footnote{Randglosse: Vater, das war Israel.} beweinete ihn. Aber die Midianiter verkauften ihn in Egypten dem Potiphar, des Pharao Hofmeister.
\end{schriftzitat}

1. Da folget nun eine schöne, lustige Historie von dem vierten Patriarchen Joseph. Hier siehe\footnote{Genauer! „sihestu“.} zum ersten, wie fein der Heilige Geist die Untugend der heiligen Patriarchen daher schreibt,
\quellseite{282}
daß er vor der Vernunft billig sollte schweigen, so er nichts Besseres müßte von ihnen zu schreiben. Zuvor haben sie eine ganze Stadt ermordet, darnach beschläft Ruben seines Vaters Weib; hier thun sie zuhauf, und berauben den Vater seines liebsten Sohnes, ihres eigenen Bruders, und das allein darum, daß sie sehen, daß ihn der Vater lieber hatte denn sie alle. Denn Benjamin war noch ein Kind, so war Joseph von der rechten Frau, Rahel, und sollte der rechte Erbe werden; dazu war es ein fromm Kind, wie der Text sagt, daß er nicht konnte leiden, wenn die Brüder unrecht handelten, sondern brachte es vor den Vater; so können sie nicht leiden, daß ihn der Vater lieb hat, und macht ihm einen bunten Rock, sind ihm so feind, daß sie ihm kein freundlich Wort können zusprechen.

2. Das ist je zu viel von des Patriarchen Kindern, daß da die christliche und brüderliche Liebe so gar verlicht. Aber GOtt schenkt ihnen die Sünde auch, allen, so da glauben, zum Troste, daß wir sehen, wie er nicht nach Werken richtet, daß keine Sünde verdammen kann, wenn man nur den Glauben nicht verliert. Darum sind es eitel feine, liebliche und tröstliche Historien. Das andere ist eitel unfreundlich, feindselig Ding, daß unsere Lehrer gepredigt haben von Heiligen, als hätten sie nie Fleisch und Blut gehabt.

3. Das ist nun der letzte und größte Puß, den der Vater in seinem Alter erleiden mußte; denn das hernach folgt, von der Theurung, ist nicht so groß gewesen. Also beschreibt der Heilige Geist dieses heiligen Vaters Elende von Jugend auf. So bald als er vom Vater gezogen ist, hat er immer ein Unglück über das andere müssen haben. Esau, sein Bruder, war der Jüngste und das liebe Kind, er mußte Aschenbrödel sein, darnach zwanzig Jahre lang dienen, und täglich viel leiden, darnach mit dem Engel kämpfen, und zuletzt, als er meint, nun mit Ruhe zu sitzen, hebt sich erst das größte Herzeleid, das er sieht an seinen Kindern. Noch war er der einige Mann in der ganzen Welt, auf den GOtt sonderlich ein Auge hatte, und ihm die reichsten und größten Verheißungen gethan hatte.

4. Aber da war nichts zu sehen, denn das Widerspiel, als hätte GOtt sein gar vergessen, und sollte alles zurück gehen. Denn er hatte ihm nun sein Weib, die Rahel, welche er am liebsten hatte, genommen; jetzt nimmt er ihm auch das liebste Kind, den Joseph, welcher der erste war von der rechten Frau. (Er) hatte nun keinen Trost mehr, denn auf den Sohn, daß er sollte der sein, darauf die Verheißung lautete, davon Christus kommen sollte; der Trost muß nun auch hinweg, daß er keine Hoffnung mehr weiß. Da hat sein Herz müssen zappeln und denken: Wie geht das zu, daß mich GOtt so verläßt, und nimmt mir allen Trost, den ich habe? Vielleicht habe ich ihn erzürnt, daß er mir die Verheißung wieder nehmen will.

5. Siehe, das sind die rechten güldenen Legenden, darinne uns GOtt lehrt, wie er seine Heiligen kocht und brät, und so mit ihnen spielt, als sei es alles erlogen, was er ihnen verheißt. Er hatte ihm gelobt einen großen Samen zu geben, und ihn [zu] segnen; und nimmt ihm alles, was er zum liebsten hat, reißt ihm dazu den Trost und Hoffnung aus dem Herzen, daß er nicht weiß, ob er mit ihm zürnt oder Freund ist: noch muß er an dem Worte halten, und GOtt bloß heimgeben, wie er seine Zusagung erfüllen will. Solches sollten wir auch lernen, ob uns GOtt ließe nehmen, was uns lieb ist, daß wir könnten hoffen, und uns trösten, daß er uns dennoch nicht lassen wolle, sondern viel mehr und Besseres geben. Wie wehe ihm nun das gethan habe, lasse ich einen jeglichen selbst denken, sonderlich dabei, daß die Natur solcher heiligen Leute rechtschaffen und nicht so verderbt ist. Darum hat es das Herz wohl gefühlt, und nicht so in Wind geschlagen, wie etliche lose Weiber thun; darum sagt der Text: „Und alle seine Söhne und Töchter traten auf, daß sie ihn trösteten; aber er wollte sich nicht trösten lassen.“

6. Damit zeigt GOtt an, daß er die Natur nicht gar verwirft, sondern in den Heiligen bleiben läßt, also, daß er sie bessert und rechtfertig macht. Zuweilen läßt er wohl einen so trunken werden in Heiligen Geiste, daß er keines Dinges auf Erden achtet; aber sonst gemeinsiglich läßt er die Natur fühlen und bewegt werden, wie er sie geschaffen hat. Denn sofern nicht Sünde mit regiert, ist keine natürliche Bewegung böse; wie wir in Christo sehen, der allerlei gefühlt und empfunden hat nach der Natur, wie ein anderer Mensch. Also ist eine natürliche Neigung, daß die Eltern die Kinder lieb haben, und sich betrüben, wenn es ihnen übel geht; wie man
\quellseite{283}
in unvernünftigen und wilden Thieren steht, daß sie auch ihr Leben bei ihnen zusetzen. Wiederum ist es auch natürlich, daß die Kinder die Eltern lieb haben, und sich alles Gutes zu ihnen versehen. So war es ein natürlich Vaterherz, daß Jakob das Kind am liebsten hatte, weil er von der rechten Frau war, und in seinem Alter gezeugt.

7. Solche Art sind für sich selbst nicht böse, wenn man’s nicht übel braucht; und zwar, wenn uns GOtt nicht solche Natur eingepflanzet hätte, so ließe man die Kinder in ihrem Stande verderben; wiederum, wo die Natur nicht hielte, die Eltern zu lieben, so würde kein Kind den Eltern gehorchen.

8. Das rede ich darum, daß man nicht meine, die Heiligen seien Holz oder Steine gewesen, wie man uns gepredigt hat. Maria, die Mutter Christi, ist auch stark gewesen; aber als sie ihn sahe am Kreuz hangen, ging ihr ein schneidendes Schwert durch das Herz [Luc. 2, 35. Joh. 19, 25. f.]. Also ist es nicht wohl möglich, daß sich ein Märtyrer nicht fürchte vor dem Tode, denn auch Christus selbst sich davor entsetzt hat, und doch nicht wider GOtt gethan. GOtt will die Seinen also üben mit Wohl- und Wehethun, beide, wenn er gibt und nimmt, daß doch das Herz gegen ihm immerdar recht stehe.

9. Aber es sind viel närrischer Heiligen gewesen, sonderlich unter den Altvätern in der Wüste, die da alle natürliche Neigung wollten brechen und dämpfen. Als man liest von zwei Brüdern, die nicht wollten ihre Mutter\footnote{Wald und die Erlanger: Mütter.} zulassen, sie zu sehen, und sprachen: Sie sollte warten, bis sie einander im ewigen Leben sähen. Solches lesen wir nicht in dieser rechten Heiligen Legenden, welchen jene nicht werth sind, die Schuhe abzuziehen, und vielleicht viele, so man für Heilige hält, zum Teufel gefahren sind. GOtt will nicht, daß wir die Creaturen verachten und wegwerfen, sondern brauchen sollen, was er uns gibt, doch alles in seinen Willen setzen. Also, daß ein Christ nach dem Geiste und Gewissen nicht achtet zeitlich Ding, aber nach dem Leibe nimmt er sich allerlei an, daß in äußerlichem Wesen kein Unterschied ist unter Christen und andern Leuten. Das sei kurz geredet von der Historie. Nun müssen wir auch die heimliche Deutung handeln.


10. Diese Historie ist reich und lieblich mit Figuren und Deutung, bis an das Ende dieses Buches. Denn in Josephs Person hat GOtt auf das allerfeinste Christum und sein ganzes Reich geistlich abgemalt. So ist nun die Summa von dieser Figur: Wie es Joseph geht mit seinen Brüdern, also geht es Christo mit seinen Brüdern, das ist, mit den Juden. Zum ersten sagt der Text, wie Jakob den Joseph lieber hatte, denn seine Brüder, darum machte er ihm einen bunten Rock, das war ein langer Rock, von mancherlei Faden gewirkt, und jeglicher von sonderlicher Farbe; und sagt Ursache, warum er ihn lieb hatte, nämlich, daß er ihn in seinem Alter gezeugt hatte. Denn wir haben gehört, daß er über siebenzig Jahre war, ehe er ein Weib nahm; so ist Joseph geboren ohngefähr im neunzigsten Jahre seines Alters. Das ist nun Jesus Christus, das rechte liebe Kind GOttes [Eph. 1, 6.] und erstgeborne Sohn [Col. 1, 15.], durch welchen alle anderen zur Gnade und Huld des Vaters kommen, also, daß er niemand will ansehen, denn durch diesen Sohn.

11. Der bunte Rock aber ist die christliche Kirche, welche ist der Schmuck und das Kleid Christi. Die heißt darum ein bunter Rock, daß sie begabt ist mit mancherlei Gaben und Tugenden des Heiligen Geistes. Es ist ein Kleid, und doch nicht einerlei Fadens oder Farbe; fein unterschieden, und doch zusammen gewirkt. Das ist es, das Paulus auslegt zum Corinthern [1. Ep. 12, 4. ff.] und Ephesern [Cap. 4, 4. ff.], daß ein Christ solche Gnade hat, der andere eine andere; etliche sind Propheten, etliche Lehrer, etliche können Wunder thun. Mancherlei sind die Gaben, Aemter und Kräfte, sagt St. Paulus; aber es ist Ein Geist, Ein Glaube, und einerlei Sinn. Denn was Einer von Christo glaubt und gewiß ist, das glauben sie alle. Dieser Rock reicht ihm bis an die Fersen, das ist, die Christenheit mit ihren Gaben währet durch und durch, bis an das Ende der Welt.

12. Daß er aber den Sohn in seinem Alter gezeugt, bedeutet, daß Christus gesandt und öffentlich gepredigt ist am Ende der Welt. Denn das Evangelium ist das Abendmahl und die letzte Predigt, die der Welt verkündigt wird, nach welcher man keine mehr erwarten darf.

13. Daß ihm aber seine Brüder feind waren, sonderlich um der Träume willen, die da vorher
\quellseite{284}
gehen, und von Zukünftigen sagen, ist, daß Christus in dem Volke gewesen ist, ehe er geboren ward; denn die Propheten lange zuvor von ihm predigten, daß er kommen und König werden würde, und Herr über alle Brüder [Matth. 22, 43. 44.]. Darum hob sich der Hader, daß die Propheten immer verfolgt wurden, bis auf Christum; aber als Christus selbst kam, waren sie ihm erst todfeind, und trachteten ihm nach, ihn zu tödten.

14. Die zwei Träume aber, von den Garben und von der Sonne und Monden, deuten nichts anders, denn die Schrift und Historien des Alten Testamentes, welche allenthalben auf Christum weisen, daß er sollte der sein, vor dem sie sich sollten bücken, und [die] Kniee beugen alle Patriarchen und Propheten.

15. Item, daß Joseph die Missethat seiner Brüder vor den Vater bringt, ist auch eine Ursache, darum sie ihm feind waren. Da zeigt auch die Schrift, daß die Patriarchen dennoch viel Büberei haben getrieben, daß sie nicht viel gutes Geschrei hatten. Hier ist aber ein Unterschied zwischen den Afterredern und Frommen. Denn der fromme Joseph, hat sie nicht ausgetragen, noch ein böse Geschrei von ihnen gemacht, sondern, wenn er von andern ein böse Geschrei hörte, sagt er es dem Vater heimlich, und bringt es nicht weiter aus. Also sollen wir auch thun: wenn du etwas Böses weißest von deinem Nächsten, so schweige, und mache es nicht rüchtig, sondern lasse es zwischen dir und ihm bleiben; ist es aber schon rüchtig worden, kann man nicht anders rathen, denn daß man es der Obrigkeit ansage, so die Macht hat zu strafen. Also hat nun auch Christus gethan durch die Propheten und seine Predigt. Wenn das Volk Abgötterei trieb und den falschen Propheten anhing, ließ er sie strafen und vermahnen, daß sie davon ließen, und auf Christum warteten.

16. Daß ihn nun der Vater ausschickt zu den Brüdern gen Sichem, und er sie nicht findet; denn sie waren nicht da geblieben, sondern fortgezogen gen Dothan; ist, daß die Juden dazu geordnet gewesen sind, daß sie sollten GOttes Satzung und Gebote tragen; aber sie sind davon gewichen, und haben ihnen eigene Aufsätze gemacht.

17. Als sie nun Joseph von ferne sehen, sagen sie unter einander: „Siehe, da kommt der Träumer her“; das ist, da sie die Schrift hörten, daß er vorhanden wäre, konnten sie es nicht leiden, dachten ihn bald unterzudrücken, und umzubringen; darnach griffen sie ihn, und zogen ihm den Rock aus, machten, daß alle Jünger und das ganze Volk von Christo fiel, daß er allein und bloß mußte stehen, und warfen ihn in eine Grube, das ist, sie brachten ihn zum Tode.

18. Aber als die Ismaeliter kamen, zogen sie ihn heraus, und verkauften ihn, das ist, da die Juden sahen, daß die Heiden zu Christo kamen, ließen sie ihn hinfahren, und verloren ihn mit allen Gütern, um ihres Genießes und Bauches willen. Also kam er in ein fernes, fremdes Land, das ist, unter die Heiden, bis an der Welt Ende, die zuvor nie GOttes Volk gewesen waren.

19. Daß aber Ruben wiederkommt und sucht ihn, findet ihn aber nicht, bedeutet, daß die Juden, als sie ihn verworfen hatten, froh gewesen wären, daß sie ihn hätten mögen wieder haben; aber es war zu lange geharrt.

20. Als er nun hinweg ist, nehmen sie seinen Rock, und tunken ihn in Bocksblut, und schicken ihn dem Vater heim, und er nicht anders meint, denn ein wildes Thier habe ihn zerrissen, bedeutet nichts anders denn die Schmach und Leid, so Christo widerfährt, welche wir noch heute müssen leiden. Denn der bunte Rock muß ihm noch immer ausgezogen, und in Bocksblut gefüllt, dazu dem Vater von der Welt vorgebildet werden, als habe ihn ein wildes Thier zerrissen. Denn was wir von Christo predigen, sammt seinem Geiste, Gnade und allen Gütern, damit der Vater ihn und uns geziert hat, das ziehen sie uns aus, unsere Brüder, die uns die Nächsten und die Besten in der Christenheit wollen sein, und sagen, es sei nichts, und beschudeln es im Blute, das ist, schänden’s und lästern’s auf das ärgste; wie er von den Juden geschändet und geschmäht ward, als wäre sein Ding des Teufels. Deß müssen wir uns erwägen. Denn die Schrift hat es beides zuvor verkündigt, und mit Figuren abgemalt. Es ist ein mercklicher Punkt, daß ihm seine eigenen Brüder also mitfahren. Das thut auch am meisten wehe, daß, die das Evangelium sollen regieren, handhaben und ehren, müssen es also beschudeln, schänden und ausrichten\footnote{ausrichten = üble Nachrede machen.} daß es nichts werth ist.

21. Dazu sind sie noch so vermessen, und
